% L’impérialisme en Afrique — Histoire 1re Tech — Cours signé OnBuch+
% Programme officiel MINESEC. Compiler avec : tectonic N02-imperialisme-afrique.tex
\newcommand{\DOCMATIERE}{Histoire}
\newcommand{\DOCNIVEAU}{1re Tech}
\newcommand{\DOCMODULE}{Histoire – classe de Première technique}
\newcommand{\DOCLECON}{N02}
\newcommand{\DOCTITRE}{L’impérialisme en Afrique}
% OnBuch+ — préambule « cours signés » ENRICHI (compilé avec tectonic / XeLaTeX)
% Système visuel « Pop » d'OnBuch+ : fond crème (jamais blanc), bordures encre
% épaisses, ombres « dures » décalées, titres Archivo Black en capitales.
% Version MATHÉMATIQUES : graphes (pgfplots/tikz), boîtes pédagogiques
% (définition, propriété, méthode, attention, le savais-tu, exercice résolu,
% exercice, TP, bilan), page de garde et signature OnBuch+ ancrée en bas.
%
% Macros attendues dans main.tex AVANT \input{preamble} :
%   \DOCMATIERE  \DOCNIVEAU  \DOCMODULE  \DOCLECON (numéro)  \DOCTITRE
\documentclass[11pt]{article}

\usepackage{fontspec}
\usepackage[french]{babel}
\usepackage[a4paper,margin=2cm,top=3.4cm,bottom=2.8cm,headheight=1.4cm,footskip=1.3cm]{geometry}
\usepackage{amsmath,amssymb,amsfonts}
\usepackage{mathtools} % \xmapsto, \coloneqq... souvent utilisés par les modèles
\usepackage{cancel} % simplifications barrées (bilans de forces)
% \abs{z} (module, valeur absolue) et \norm{u} (norme), utilisés par les modèles
\providecommand{\abs}{}\let\abs\relax\DeclarePairedDelimiter{\abs}{\lvert}{\rvert}
\providecommand{\norm}{}\let\norm\relax\DeclarePairedDelimiter{\norm}{\lVert}{\rVert}
\usepackage[table]{xcolor} % [table] : fournit aussi \rowcolors (alternance de lignes)
\usepackage{pagecolor}
\usepackage{tikz}
\usetikzlibrary{arrows.meta,positioning,calc,shapes.geometric,shapes.misc,shapes.arrows,decorations.pathmorphing,decorations.pathreplacing,decorations.markings,patterns,shadows,mindmap,trees,fit,backgrounds,matrix,intersections,angles,quotes}
\usepackage{pgfplots}
\usepackage[version=4]{mhchem} % équations nucléaires \ce{^{235}_{92}U}
\usepackage[european,siunitx]{circuitikz} % schémas électriques
\pgfplotsset{compat=1.18}
\usepgfplotslibrary{fillbetween,groupplots} % aires entre courbes (calcul intégral)
\usepackage{enumitem}
\usepackage{fancyhdr}
% [most] charge toutes les bibliothèques tcolorbox (listings, minted, raster,
% documentation...) et gonfle inutilement la mémoire TeX sur un document long
% (~300 boîtes) : on ne charge que ce qui est réellement utilisé.
\usepackage[skins,breakable]{tcolorbox}
\usepackage{graphicx}
\usepackage{colortbl}
\usepackage{booktabs}
\usepackage{tabularx}
\usepackage{multirow}
\usepackage{diagbox} % case d'en-tête diagonale des tableaux à double entrée
% Colonnes extensibles centrée / gauche / droite (C, L, R), souvent utilisées par les modèles
\newcolumntype{C}{>{\centering\arraybackslash}X}
\newcolumntype{L}{>{\raggedright\arraybackslash}X}
\newcolumntype{R}{>{\raggedleft\arraybackslash}X}
\usepackage{array}
\usepackage{multicol}
\usepackage{siunitx}
% parse-numbers=false : accepte aussi \SI{2,5 \times 10^{-3}}{...}
\sisetup{locale=FR,output-decimal-marker={,},group-minimum-digits=4,parse-numbers=false}
\DeclareSIUnit\ohm{\ensuremath{\symup{\Omega}}} % U+2126 absent de la police
\DeclareSIUnit\division{div} % réglages d'oscilloscope (ms/div, V/div)
\DeclareSIUnit\tr{tr} % tours (vitesse de rotation, ex. \SI{1500}{\tr\per\minute})
\DeclareSIUnit\molar{M} % concentration molaire (ex. \SI{3}{\molar}, \SI{1}{\micro\molar})
\DeclareSIUnit\an{an} % année (le modèle confond parfois avec \year, un compteur TeX) — ex. \SI{5}{\kilo\meter\per\an}
\DeclareSIUnit\FCFA{FCFA} % franc CFA (ex. \SI{500}{\FCFA\per\kilo\gram})
\DeclareSIUnit\degreeBrix{\textdegree Brix} % teneur en sucre d'un sirop/jus (réfractométrie)
\DeclareSIUnit\month{mois} % le modèle utilise parfois \month, un compteur TeX, comme unité de durée
\DeclareSIUnit\microliter{\micro\liter} % le modèle écrit parfois \microliter en un seul token
\DeclareSIUnit\million{millions} % dénombrement (ex. \SI{15}{\million\per\mL} spermatozoïdes)
\usepackage{qrcode}
% \degree : commande fréquemment utilisée par les modèles pour un angle ou une
% température en dehors de \SI{}{} (non fournie par les paquets chargés).
\providecommand{\degree}{^{\circ}}
% \qed (fin de démonstration), utilisé par les modèles sans amsthm
\providecommand{\qed}{\hfill$\square$}
% \barre : double barre des valeurs interdites dans un tableau de variations (tkz-tab)
\providecommand{\barre}{\ensuremath{\|}}
% environnement proof (amsthm non chargé) utilisé par les modèles
\ifdefined\proof\else\newenvironment{proof}[1][Démonstration]{\par\noindent\textit{#1.}\ }{\qed\par}\fi
\usepackage{lastpage}
\usepackage{hyperref}
\hypersetup{hidelinks}

% ============================================================
% Palette « Pop » OnBuch+ — mêmes hex que la classe P (plus_ui.dart)
% ============================================================
\definecolor{poporange}{HTML}{F59321}
\definecolor{popdark}{HTML}{E07A0C}
\definecolor{popcream}{HTML}{FFF6E8}
\definecolor{popink}{HTML}{1C1714}
\definecolor{poppurple}{HTML}{7A5AE0}
\definecolor{popgreen}{HTML}{2E9E6A}
\definecolor{poppink}{HTML}{E0457B}
\definecolor{popblue}{HTML}{2F7DE1}
\definecolor{popgold}{HTML}{C98A12}
\definecolor{popmuted}{HTML}{8A7F76}
\definecolor{popline}{HTML}{EADFD2}
\definecolor{popgray}{HTML}{8A7F76}
\colorlet{popgrey}{popgray}
% Alias tolérants (le modèle nomme parfois les couleurs différemment)
\colorlet{popwhite}{white}
\colorlet{popblack}{popink}
\colorlet{popred}{poppink}
\colorlet{popyellow}{popgold}
% Teintes claires (fonds de tableaux/encadrés) — le modèle nomme parfois
% les couleurs avec un suffixe L (ex. popblueL) sans les avoir définies.
\colorlet{poporangeL}{poporange!20}
\colorlet{popdarkL}{popdark!20}
\colorlet{poppurpleL}{poppurple!20}
\colorlet{popgreenL}{popgreen!20}
\colorlet{poppinkL}{poppink!20}
\colorlet{popblueL}{popblue!20}
\colorlet{popgoldL}{popgold!20}
\colorlet{popredL}{popred!20}
\colorlet{popgrayL}{popgray!20}
% Teintes claires pour fonds de tableaux / zones
\colorlet{popblueL}{popblue!10!white}
\colorlet{poporangeL}{poporange!14!white}
\colorlet{popgreenL}{popgreen!12!white}
\colorlet{poppurpleL}{poppurple!12!white}
\colorlet{poppinkL}{poppink!12!white}
\colorlet{popgoldL}{popgold!14!white}

\pagecolor{popcream}

% ============================================================
% Typographie
% ============================================================
\defaultfontfeatures{Ligatures=TeX}
\setmainfont{PlusJakartaSans-Medium.ttf}[Path=../../../fonts/,BoldFont=PlusJakartaSans-Bold.ttf]
% Maths en sans-serif (Fira Math) avec chiffres et lettres droites de Plus
% Jakarta, pour que les formules se fondent dans le texte.
\usepackage[math-style=ISO,bold-style=ISO]{unicode-math}
\setmathfont{FiraMath-Regular.otf}[Path=../../../fonts/]
\setmathfont{PlusJakartaSans-Medium.ttf}[Path=../../../fonts/,range={up/num,up/{Latin,latin},\mathup}]
\usepackage{icomma} % virgule décimale sans espace en mode math
% Caractères mathématiques Unicode absents des polices de texte (Plus Jakarta,
% Archivo Black) : rendus par leur équivalent mathématique, en texte comme en
% formule — sinon ils s'affichent en carré vide (ex. « Arithmétique dans ℤ »).
\usepackage{newunicodechar}
\newunicodechar{ℝ}{\ensuremath{\mathbb{R}}}
\newunicodechar{ℤ}{\ensuremath{\mathbb{Z}}}
\newunicodechar{ℂ}{\ensuremath{\mathbb{C}}}
\newunicodechar{ℕ}{\ensuremath{\mathbb{N}}}
\newunicodechar{ℚ}{\ensuremath{\mathbb{Q}}}
\newunicodechar{𝔻}{\ensuremath{\mathbb{D}}}
\newunicodechar{α}{\ensuremath{\alpha}}
\newunicodechar{β}{\ensuremath{\beta}}
\newunicodechar{γ}{\ensuremath{\gamma}}
\newunicodechar{δ}{\ensuremath{\delta}}
\newunicodechar{ε}{\ensuremath{\varepsilon}}
\newunicodechar{θ}{\ensuremath{\theta}}
\newunicodechar{λ}{\ensuremath{\lambda}}
\newunicodechar{μ}{\ensuremath{\mu}}
\newunicodechar{σ}{\ensuremath{\sigma}}
\newunicodechar{τ}{\ensuremath{\tau}}
\newunicodechar{φ}{\ensuremath{\varphi}}
\newunicodechar{ψ}{\ensuremath{\psi}}
\newunicodechar{ω}{\ensuremath{\omega}}
\newunicodechar{Σ}{\ensuremath{\Sigma}}
% majuscules produites par \MakeUppercase dans les titres (α-aminés -> Α)
\newunicodechar{Α}{\ensuremath{\alpha}}
\newunicodechar{Β}{\ensuremath{\beta}}
\newunicodechar{Γ}{\ensuremath{\Gamma}}
\newunicodechar{Δ}{\ensuremath{\Delta}}
\newunicodechar{Ε}{\ensuremath{\varepsilon}}
\newunicodechar{Θ}{\ensuremath{\Theta}}
\newunicodechar{Λ}{\ensuremath{\Lambda}}
\newunicodechar{Μ}{\ensuremath{\mu}}
\newunicodechar{Τ}{\ensuremath{\tau}}
\newunicodechar{Φ}{\ensuremath{\Phi}}
\newunicodechar{Ψ}{\ensuremath{\Psi}}
\newunicodechar{Ω}{\ensuremath{\Omega}}
\newunicodechar{↦}{\ensuremath{\mapsto}}
\newunicodechar{∈}{\ensuremath{\in}}
\newunicodechar{∉}{\ensuremath{\notin}}
\newunicodechar{∧}{\ensuremath{\wedge}}
\newunicodechar{⇔}{\ensuremath{\Leftrightarrow}}
\newunicodechar{⟺}{\ensuremath{\Longleftrightarrow}}
\newunicodechar{⇒}{\ensuremath{\Rightarrow}}
\newunicodechar{⟹}{\ensuremath{\Longrightarrow}}
\newunicodechar{∘}{\ensuremath{\circ}}
\newunicodechar{∪}{\ensuremath{\cup}}
\newunicodechar{∩}{\ensuremath{\cap}}
\newunicodechar{≡}{\ensuremath{\equiv}}
\newunicodechar{⊂}{\ensuremath{\subset}}
\newunicodechar{⊥}{\ensuremath{\perp}}
\newunicodechar{∃}{\ensuremath{\exists}}
\newunicodechar{∀}{\ensuremath{\forall}}
\newunicodechar{∛}{\ensuremath{\sqrt[3]{\,}}}
\newunicodechar{∜}{\ensuremath{\sqrt[4]{\,}}}
\newunicodechar{⃗}{\ensuremath{{}^{\to}}}
\newunicodechar{ⁿ}{\ifmmode^{n}\else\textsuperscript{\ensuremath{n}}\fi}
\newunicodechar{ˣ}{\ifmmode^{x}\else\textsuperscript{\ensuremath{x}}\fi}
\newunicodechar{ᵇ}{\ifmmode^{b}\else\textsuperscript{\ensuremath{b}}\fi}
\newunicodechar{ʳ}{\ifmmode^{r}\else\textsuperscript{\ensuremath{r}}\fi}
\newunicodechar{ᵘ}{\ifmmode^{u}\else\textsuperscript{\ensuremath{u}}\fi}
\newunicodechar{ʸ}{\ifmmode^{y}\else\textsuperscript{\ensuremath{y}}\fi}
\newunicodechar{ᵃ}{\ifmmode^{a}\else\textsuperscript{\ensuremath{a}}\fi}
\newunicodechar{ᵉ}{\ifmmode^{e}\else\textsuperscript{\ensuremath{e}}\fi}
\newunicodechar{ᵏ}{\ifmmode^{k}\else\textsuperscript{\ensuremath{k}}\fi}
\newunicodechar{ᵖ}{\ifmmode^{p}\else\textsuperscript{\ensuremath{p}}\fi}
\newunicodechar{ᴺ}{\ifmmode^{N}\else\textsuperscript{\ensuremath{N}}\fi}
\newunicodechar{ᶜ}{\ifmmode^{c}\else\textsuperscript{\ensuremath{c}}\fi}
\newunicodechar{ʰ}{\ifmmode^{h}\else\textsuperscript{\ensuremath{h}}\fi}
\newunicodechar{ᵠ}{\ifmmode^{q}\else\textsuperscript{\ensuremath{q}}\fi}
\newunicodechar{ₙ}{\ifmmode_{n}\else\textsubscript{\ensuremath{n}}\fi}
\newunicodechar{ₐ}{\ifmmode_{a}\else\textsubscript{\ensuremath{a}}\fi}
\newunicodechar{ᵢ}{\ifmmode_{i}\else\textsubscript{\ensuremath{i}}\fi}
\newunicodechar{ᵣ}{\ifmmode_{r}\else\textsubscript{\ensuremath{r}}\fi}
\newunicodechar{ₖ}{\ifmmode_{k}\else\textsubscript{\ensuremath{k}}\fi}
\newunicodechar{ᵦ}{\ifmmode_{\beta}\else\textsubscript{\ensuremath{\beta}}\fi}
\newunicodechar{ₓ}{\ifmmode_{x}\else\textsubscript{\ensuremath{x}}\fi}
\newfontfamily\popdisplay{ArchivoBlack-Regular.ttf}[Path=../../../fonts/]
\newfontfamily\poplogo{SpaceGrotesk-Bold.ttf}[Path=../../../fonts/]
\newfontfamily\popsemi{PlusJakartaSans-SemiBold.ttf}[Path=../../../fonts/]
\newfontfamily\popxbold{PlusJakartaSans-ExtraBold.ttf}[Path=../../../fonts/]
\newfontfamily\popxboldls{PlusJakartaSans-ExtraBold.ttf}[Path=../../../fonts/,LetterSpace=3.0]

\setlist[enumerate]{leftmargin=1.7em,itemsep=5pt,topsep=4pt}
\setlist[itemize]{leftmargin=1.7em,itemsep=4pt,topsep=4pt,label={\color{poporange}\textbullet}}
\setlength{\parindent}{0pt}
\setlength{\parskip}{5pt}
\renewcommand{\arraystretch}{1.25}

% pgfplots : style Pop par défaut
\pgfplotsset{
  every axis/.append style={axis line style={popink,line width=1pt},tick style={popink},
    grid style={popline,line width=0.5pt},label style={font=\small},tick label style={font=\footnotesize}},
  popcurve/.style={poporange,line width=1.8pt,no markers,smooth},
}
% Même style disponible dans un \draw TikZ simple (hors pgfplots)
\tikzset{popcurve/.style={poporange,line width=1.8pt}}
\tikzset{popgrid/.style={popline,very thin}}
\colorlet{popcurve}{poporange} % le nom du style est parfois utilisé comme couleur

% ============================================================
% Pill / squiggle
% ============================================================
\newcommand{\poppill}[3]{%
  \tikz[baseline=-0.55ex]\node[draw=popink,line width=1.2pt,rounded corners=2.6pt,%
    fill=#2,inner xsep=6.5pt,inner ysep=3pt]{%
    \popxboldls\fontsize{7}{7}\selectfont\color{#3}\MakeUppercase{#1}};%
}
\newcommand{\popsquiggle}[1]{%
  \tikz[baseline=-0.4ex]{
    \draw[#1,line width=2.6pt,line cap=round]
      plot [smooth] coordinates {(0,0.09) (0.34,0.01) (0.68,0.21) (1.02,0.01) (1.36,0.10)};
  }%
}
% Mot-clé surligné (marqueur orange sous le texte)
\newcommand{\cle}[1]{{\popxbold\color{popdark}#1}}

% ============================================================
% En-tête / pied
% ============================================================
\pagestyle{fancy}
\fancyhf{}
\renewcommand{\headrulewidth}{0pt}
\renewcommand{\footrulewidth}{0pt}
\lhead{\raisebox{-0.15ex}{\poplogo\fontsize{13}{13}\selectfont\color{popink}OnBuch\color{poporange}+}}
\rhead{\poppill{\DOCMATIERE\ \textbullet\ \DOCNIVEAU\ \textbullet\ Leçon \DOCLECON}{white}{popink}}
\lfoot{\popxbold\fontsize{7.6}{7.6}\selectfont\color{popmuted}\MakeUppercase{Cours signé OnBuch+}\ \textbullet\ {\popsemi\DOCTITRE}}
\rfoot{\tikz[baseline=-0.6ex]\node[rounded corners=8.5pt,fill=popink,minimum height=17pt,minimum width=17pt,inner xsep=5pt,inner sep=0]{\popxbold\fontsize{8}{8}\selectfont\color{popcream}\thepage\,/\,\pageref*{LastPage}};}

% ============================================================
% Titres
% ============================================================
\newcommand{\chaptitre}[2]{%
  \begin{tcolorbox}[enhanced,colback=poporange,colframe=popink,boxrule=2.6pt,arc=11pt,
      drop shadow=popink,left=15pt,right=15pt,top=12pt,bottom=11pt,before skip=3pt,after skip=18pt]
    \poppill{#2}{popink}{popcream}\par\vspace{9pt}
    {\raggedright\hyphenpenalty=10000\exhyphenpenalty=10000\popdisplay\fontsize{22}{24}\selectfont\color{popcream}\MakeUppercase{#1}\par}
  \end{tcolorbox}
}
% Section numérotée : pastille orange avec le numéro + titre Archivo
\newcounter{popsec}
\newcommand{\coursec}[1]{%
  \stepcounter{popsec}\setcounter{popsub}{0}%
  \par\vspace{16pt}\noindent
  \begin{minipage}{\linewidth}
  \tikz[baseline=-0.7ex]\node[draw=popink,line width=1.6pt,fill=poporange,rounded corners=4pt,minimum size=22pt,inner sep=0]{\popdisplay\fontsize{12}{12}\selectfont\color{popcream}\thepopsec};%
  \hspace{8pt}{\popdisplay\fontsize{15}{17}\selectfont\color{popink}\MakeUppercase{#1}}\par
  \vspace{4pt}\noindent{\color{poporange}\rule{36pt}{3.6pt}}
  \end{minipage}\par\nopagebreak\vspace{8pt}
}
\newcounter{popsub}
\newcommand{\courssub}[1]{%
  \stepcounter{popsub}%
  \par\vspace{9pt}\noindent{\popxbold\fontsize{11.5}{13}\selectfont\color{popink}{\color{poporange}\thepopsec.\thepopsub}\hspace{6pt}#1}\par\nopagebreak\vspace{4pt}
}

% ============================================================
% Boîtes pédagogiques — même recette : fond blanc, bordure épaisse colorée,
% ombre dure encre, badge plein. Titre optionnel [#1].
% ============================================================
\tcbset{popbase/.style={breakable,enhanced,colback=white,boxrule=2.3pt,arc=9pt,
  shadow={3pt}{-3pt}{0pt}{popink},left=14pt,right=14pt,top=11pt,bottom=11pt,before skip=11pt,after skip=15pt,
  fontupper=\popsemi\fontsize{10.6}{14.5}\selectfont\color{popink},
  fontlower=\popsemi\fontsize{10.6}{14.5}\selectfont\color{popink}}}
% Coche dessinée (le glyphe ✓ n'existe pas dans les polices du document)
\newcommand{\popcheck}[1]{\tikz[baseline=-0.5ex]\draw[#1,line width=1.6pt,line cap=round,line join=round](-0.13,0.0)--(-0.04,-0.09)--(0.13,0.1);}
\providecommand{\coche}{\popcheck{popgreen}}
\providecommand{\resultat}[1]{\par\smallskip\noindent\fcolorbox{popgreen}{popgreenL}{\parbox{\dimexpr\linewidth-2\fboxsep-2\fboxrule\relax}{\textbf{Résultat.} #1}}\par\smallskip}
\newcommand{\popboxhead}[3]{\poppill{#1}{#2}{white}\ifx\relax#3\relax\else\hspace{6pt}{\popxbold\fontsize{10.6}{12}\selectfont\color{popink}#3}\fi\par\vspace{7pt}}

\newtcolorbox{aretenir}[1][]{popbase,colframe=popgreen,before upper={\popboxhead{À retenir}{popgreen}{#1}}}
\newtcolorbox{exemplebox}[1][]{popbase,colframe=poporange,before upper={\popboxhead{Exemple}{poporange}{#1}}}
\newtcolorbox{definition}[1][]{popbase,colframe=popblue,before upper={\popboxhead{Définition}{popblue}{#1}}}
\newtcolorbox{propriete}[1][]{popbase,colframe=poppurple,before upper={\popboxhead{Propriété}{poppurple}{#1}}}
\newtcolorbox{methode}[1][]{popbase,colframe=popgold,before upper={\popboxhead{Méthode}{popgold}{#1}}}
\newtcolorbox{attention}[1][]{popbase,colframe=poppink,before upper={\popboxhead{Attention}{poppink}{#1}}}
\newtcolorbox{experience}[1][]{popbase,colframe=popblue,colback=popblueL,before upper={\popboxhead{Expérience}{popblue}{#1}}}
% « Le savais-tu ? » : carte violette pleine (contexte, culture, Cameroun)
\newtcolorbox{savaistu}[1][]{popbase,colback=poppurple,colframe=popink,
  fontupper=\popsemi\fontsize{10.4}{14.2}\selectfont\color{white},
  before upper={\poppill{Le savais-tu\,?}{white}{poppurple}\ifx\relax#1\relax\else\hspace{6pt}{\popxbold\color{white}#1}\fi\par\vspace{7pt}}}
% Exercice résolu : énoncé en haut, \tcblower puis la solution sur fond crème
\newcounter{popexo}\newcounter{popexr}
\newtcolorbox{exoresolu}[1][]{popbase,colframe=poporange,colbacklower=popcream,
  segmentation style={popink,line width=1.2pt,dashed},
  before upper={\stepcounter{popexr}\popboxhead{Exercice résolu \thepopexr}{poporange}{#1}},
  before lower={\poppill{Solution}{popink}{popcream}\par\vspace{6pt}}}
% Exercice (non résolu en place) — la correction est dans la partie Corrigés
\newtcolorbox{exercice}[1][]{popbase,colframe=popink,
  before upper={\stepcounter{popexo}\popboxhead{Exercice \thepopexo}{popink}{#1}}}
% Corrigé
\newtcolorbox{corrige}[1][]{popbase,colframe=popgreen,colback=popgreenL,
  before upper={\popboxhead{Corrigé}{popgreen}{#1}}}
% Objectifs / prérequis / bilan
\newtcolorbox{objectifs}{popbase,colframe=popink,colback=poporangeL,
  before upper={\popboxhead{Objectifs de la leçon}{popink}{}}}
\newtcolorbox{prerequis}{popbase,colframe=popink,colback=popblueL,
  before upper={\popboxhead{Prérequis}{popblue}{}}}
\newtcolorbox{bilan}{popbase,colframe=popink,colback=white,boxrule=2.8pt,
  before upper={\popboxhead{Fiche bilan}{poporange}{L'essentiel à connaître pour l'examen}}}
% Figure encadrée avec légende
\newtcolorbox{popfigure}[1][]{enhanced,breakable=false,colback=white,colframe=popline,boxrule=1.4pt,arc=8pt,
  left=10pt,right=10pt,top=10pt,bottom=8pt,before skip=10pt,after skip=12pt,halign=center,
  lower separated=false,halign lower=center,
  fontlower=\popsemi\fontsize{9}{11.5}\selectfont\color{popmuted},
  #1}
\newcommand{\legende}[1]{\par\vspace{6pt}{\popsemi\fontsize{9}{11.5}\selectfont\color{popmuted}#1}}

% ============================================================
% Signature OnBuch+
% ============================================================
\newcommand{\poplogoinline}[1]{{\poplogo\fontsize{#1}{#1}\selectfont\color{popink}OnBuch\color{poporange}+}}
\newcommand{\signatureonbuch}{%
  \begin{tcolorbox}[enhanced,colback=popink,colframe=popink,boxrule=2.2pt,arc=10pt,
      drop shadow=poporange,left=14pt,right=14pt,top=10pt,bottom=10pt,before skip=0pt,after skip=0pt]
    \tikz[baseline=-0.7ex]\node[circle,fill=popgreen,draw=popcream,line width=1.3pt,minimum size=22pt,inner sep=0]{\popcheck{white}};%
    \hspace{9pt}\begin{minipage}[c]{0.62\linewidth}
      {\popxbold\fontsize{12}{13}\selectfont\color{popcream}Signé\ \textbullet\ {\poplogo OnBuch\color{poporange}+}}\par\vspace{2pt}
      {\raggedright\popsemi\fontsize{8}{10.5}\selectfont\color{popline}\MakeUppercase{Équipe pédagogique OnBuch+}\\\MakeUppercase{Conforme au programme officiel MINESEC}\par}
    \end{minipage}\hfill
    {\poplogo\fontsize{22}{22}\selectfont\color{popcream}OnBuch\color{poporange}+}
  \end{tcolorbox}%
}

% ============================================================
% Page de garde
% #1 titre · #2 matière · #3 niveau · #4 module · #5 n° leçon · #6 durée ·
% #7 description / objectifs (texte ou itemize)
% ============================================================
\newcommand{\pagegarde}[7]{%
  \thispagestyle{empty}
  \newgeometry{a4paper,margin=1.7cm,top=1.6cm,bottom=1.4cm}%
  \begin{tikzpicture}[remember picture,overlay]
    \fill[poporange!18] ([xshift=-3.2cm,yshift=-3.6cm]current page.north east) circle (4.6cm);
    \fill[poppurple!14] ([xshift=2.4cm,yshift=7.6cm]current page.south west) circle (3.8cm);
  \end{tikzpicture}%
  {\poplogo\fontsize{30}{30}\selectfont\color{popink}OnBuch\color{poporange}+}\hfill
  \raisebox{0.6ex}{\poppill{Cours signé \textbullet\ Premium}{popink}{popcream}}\par
  \vspace{1pt}\popsquiggle{poppurple}\par
  \vspace{8pt}
  \begin{tcolorbox}[enhanced,colback=poporange,colframe=popink,boxrule=2.8pt,arc=13pt,
      drop shadow=popink,left=18pt,right=18pt,top=12pt,bottom=13pt,after skip=10pt]
    \poppill{#2 \textbullet\ #3}{popink}{popcream}\hspace{5pt}\poppill{#4}{white}{popink}\par\vspace{8pt}
    {\popdisplay\fontsize{14}{14}\selectfont\color{popink}LEÇON #5}\par\vspace{4pt}
    {\raggedright\hyphenpenalty=10000\exhyphenpenalty=10000\popdisplay\fontsize{25}{27}\selectfont\color{popcream}\MakeUppercase{#1}\par}
    \vspace{8pt}
    {\popxbold\fontsize{9.5}{9.5}\selectfont\color{popink}\MakeUppercase{Durée conseillée : #6}}
  \end{tcolorbox}
  \begin{tcolorbox}[enhanced,colback=white,colframe=popink,boxrule=2.2pt,arc=10pt,
      drop shadow=popink,left=16pt,right=16pt,top=10pt,bottom=10pt,after skip=8pt]
    \poppill{Dans cette leçon}{poppurple}{white}\par\vspace{5pt}
    \popsemi\fontsize{9.3}{12.6}\selectfont\color{popink}#7
  \end{tcolorbox}
  \noindent\begin{minipage}[t]{0.487\linewidth}
    \begin{tcolorbox}[enhanced,colback=popgreen,colframe=popink,boxrule=2.2pt,arc=10pt,
        drop shadow=popink,left=11pt,right=11pt,top=10pt,bottom=10pt,height=3.1cm,valign=center]
      \tcbox[colback=white,colframe=popink,boxrule=1.3pt,arc=5pt,boxsep=0pt,nobeforeafter,
        left=3pt,right=3pt,top=3pt,bottom=3pt,valign=center]{\qrcode[height=1.9cm]{https://play.google.com/store/apps/details?id=com.luvvix.onbuch&hl=fr}}%
      \hspace{8pt}\begin{minipage}[c]{0.5\linewidth}
        {\popdisplay\fontsize{11}{12.5}\selectfont\color{white}\MakeUppercase{Télécharge OnBuch}}\par\vspace{4pt}
        {\raggedright\popsemi\fontsize{8.4}{11}\selectfont\color{white}Gratuit \textbullet\ Google Play\\Tuteur IA, annales, concours, résultats\par}
      \end{minipage}
    \end{tcolorbox}
  \end{minipage}\hfill
  \begin{minipage}[t]{0.487\linewidth}
    \begin{tcolorbox}[enhanced,colback=poppurple,colframe=popink,boxrule=2.2pt,arc=10pt,
        drop shadow=popink,left=11pt,right=11pt,top=10pt,bottom=10pt,height=3.1cm,valign=center]
      \tikz[baseline=-0.7ex]\node[circle,fill=white,draw=popink,line width=1.3pt,minimum size=1.6cm,inner sep=0]{\popdisplay\fontsize{24}{24}\selectfont\color{poppurple}+};%
      \hspace{8pt}\begin{minipage}[c]{0.6\linewidth}
        {\popdisplay\fontsize{11}{12.5}\selectfont\color{white}\MakeUppercase{Passe à OnBuch+}}\par\vspace{4pt}
        {\raggedright\popsemi\fontsize{8.4}{11}\selectfont\color{white}Tous les cours signés, Tuteur IA illimité, fiches PDF — onglet OnBuch+ de l'app\par}
      \end{minipage}
    \end{tcolorbox}
  \end{minipage}
  \vspace{8pt}
  \popcontenu
  \vfill
  \signatureonbuch
  \restoregeometry
  \newpage
}

% Rangée « Ce cours contient » (page de garde)
\newcommand{\poptile}[3]{% #1 couleur #2 gros texte #3 légende
  \begin{tcolorbox}[enhanced,colback=#1,colframe=popink,boxrule=2pt,arc=9pt,shadow={2.5pt}{-2.5pt}{0pt}{popink},
      left=6pt,right=6pt,top=9pt,bottom=9pt,halign=center,valign=center,height=1.75cm,width=\linewidth,nobeforeafter]
    {\popdisplay\fontsize{12.5}{13}\selectfont\color{white}\MakeUppercase{#2}}\par\vspace{3pt}
    {\centering\popsemi\fontsize{7.8}{9.5}\selectfont\color{white}#3\par}
  \end{tcolorbox}}
\newcommand{\popcontenu}{%
  \par\noindent{\popxboldls\fontsize{7.5}{7.5}\selectfont\color{popmuted}CE COURS SIGNÉ CONTIENT}\par\vspace{7pt}
  \noindent\begin{minipage}[t]{0.235\linewidth}\poptile{popblue}{Cours}{complet, illustré, pas à pas}\end{minipage}\hfill
  \begin{minipage}[t]{0.235\linewidth}\poptile{poporange}{Méthodes}{et exercices résolus}\end{minipage}\hfill
  \begin{minipage}[t]{0.235\linewidth}\poptile{poppink}{Exercices}{type examen + corrigés}\end{minipage}\hfill
  \begin{minipage}[t]{0.235\linewidth}\poptile{popgreen}{Fiche bilan}{l'essentiel à retenir}\end{minipage}\par}

% Sommaire
\newtcolorbox{sommaire}{enhanced,colback=white,colframe=popink,boxrule=2.2pt,arc=10pt,
  drop shadow=popink,left=18pt,right=18pt,top=13pt,bottom=13pt,before skip=6pt,after skip=20pt,
  before upper={\poppill{Sommaire}{poppurple}{white}\par\vspace{9pt}\popsemi\fontsize{10.6}{14.5}\selectfont\color{popink}}}

% Signature de fin de cours — poussée tout en bas de la dernière page
\newcommand{\findecours}{%
  \par\vfill\nopagebreak
  \begin{center}\popsquiggle{poporange}\end{center}\vspace{4pt}
  \signatureonbuch
}
\AtBeginDocument{\def\llbracket{\mathopen{[\mkern-3.5mu[}}\def\rrbracket{\mathclose{]\mkern-3.5mu]}}} % intervalles d'entiers (glyphe ⟦ absent de la police)
% Glyphes absents de Fira Math : redéfinis à partir de glyphes disponibles
\AtBeginDocument{%
  \def\setminus{\mathbin{\backslash}}\let\smallsetminus\setminus
  \def\vdots{\mathord{\vbox{\baselineskip3.2pt\lineskiplimit0pt\kern4pt\hbox{.}\hbox{.}\hbox{.}}}}%
  \def\ddots{\mathinner{\mkern1mu\raise6.5pt\hbox{.}\mkern2mu\raise3.6pt\hbox{.}\mkern2mu\raise0.7pt\hbox{.}\mkern1mu}}%
  \def\triangle{\mathord{\scalebox{0.9}{$\symup{\Delta}$}}}%
  \def\nabla{\mathord{\raisebox{\depth}{\scalebox{1}[-1]{$\symup{\Delta}$}}}}%
  \def\checkmark{\popcheck{popdark}}%
  \def\star{\mathbin{\ast}}\def\bigstar{\mathord{\ast}}%
  \def\diamond{\mathbin{\scalebox{0.6}{$\Diamond$}}}%
  \def\bigcup{\mathop{\mathchoice{\raisebox{-0.3em}{\scalebox{1.7}{$\cup$}}}{\scalebox{1.25}{$\cup$}}{\cup}{\cup}}\displaylimits}%
  \def\bigcap{\mathop{\mathchoice{\raisebox{-0.3em}{\scalebox{1.7}{$\cap$}}}{\scalebox{1.25}{$\cap$}}{\cap}{\cap}}\displaylimits}%
  \def\lesseqqgtr{\mathrel{\lessgtr}}\def\bigoplus{\mathop{\oplus}\displaylimits}%
}
\providecommand{\ssi}{si et seulement si} % abréviation fréquente
\AtBeginDocument{\providecommand{\wideparen}{\overparen}} % arc de cercle

%% FIGFIX-BEGIN v5 — correctifs globaux des figures (docs/onbuch-generation/tools/figfix)
\usepackage{adjustbox}\usepackage{etoolbox}\tcbuselibrary{raster}
% toute figure TikZ/pgfplots plus large que la ligne est réduite à la largeur disponible
\BeforeBeginEnvironment{tikzpicture}{\begin{adjustbox}{max width=\linewidth}}
\AfterEndEnvironment{tikzpicture}{\end{adjustbox}}
% légendes pgfplots lisibles par-dessus les courbes
\pgfplotsset{every axis/.append style={legend style={fill=white,fill opacity=0.92,text opacity=1,draw=black!15,inner sep=3pt}}}
% étiquettes d'arêtes (syntaxe edge["..."]) avec fond blanc
\tikzset{every edge quotes/.append style={fill=white,inner sep=1.2pt}}
%% FIGFIX-END
\begin{document}

\pagegarde{L’impérialisme en Afrique}{Histoire}{1re Tech}{Histoire – classe de Première technique}{N02}{3 séances (fiche de progression harmonisée)}{Cette leçon étudie la conquête de l'Afrique par les puissances européennes (France, Grande-Bretagne et autres) à la fin du XIXe siècle, ainsi que les réactions africaines, entre résistances armées et collaborations. À partir de cartes, de photographies et de données chiffrées, l'élève apprend à analyser des documents historiques et à construire des croquis pour rendre compte de ce basculement majeur de l'histoire du continent.
\vspace{4pt}
\begin{multicols}{2}\small
\begin{itemize}
\item À la fin de cette leçon, je sais définir l'impérialisme et situer chronologiquement la conquête de l'Afrique (1870-1914)
\item À la fin de cette leçon, je sais décrire les étapes et les modalités des conquêtes françaises en Afrique
\item À la fin de cette leçon, je sais comparer les conquêtes anglaises et celles des autres puissances (Allemagne, Belgique, Portugal, Italie)
\item À la fin de cette leçon, je sais localiser sur une carte les principales zones de conquête et les empires coloniaux en 1914
\item À la fin de cette leçon, je sais distinguer et illustrer les différentes formes de résistances africaines (armées, diplomatiques, religieuses)
\item À la fin de cette leçon, je sais expliquer les raisons et les formes des collaborations africaines avec les Européens
\end{itemize}\end{multicols}}

\begin{sommaire}
\begin{multicols}{2}
\begin{enumerate}
\item L'impérialisme : définition, causes et contexte
\item Les conquêtes françaises
\item Les conquêtes anglaises et des autres puissances
\item Les résistances africaines
\item Les collaborations africaines
\item Bilan et synthèse : l'Afrique colonisée en 1914
\item Activité d'intégration
\item Méthodes et exercices résolus
\item Exercices
\item Corrigés des exercices
\item Fiche bilan
\end{enumerate}
\end{multicols}
\end{sommaire}

\begin{objectifs}
\begin{itemize}
\item À la fin de cette leçon, je sais définir l'impérialisme et situer chronologiquement la conquête de l'Afrique (1870-1914)
\item À la fin de cette leçon, je sais décrire les étapes et les modalités des conquêtes françaises en Afrique
\item À la fin de cette leçon, je sais comparer les conquêtes anglaises et celles des autres puissances (Allemagne, Belgique, Portugal, Italie)
\item À la fin de cette leçon, je sais localiser sur une carte les principales zones de conquête et les empires coloniaux en 1914
\item À la fin de cette leçon, je sais distinguer et illustrer les différentes formes de résistances africaines (armées, diplomatiques, religieuses)
\item À la fin de cette leçon, je sais expliquer les raisons et les formes des collaborations africaines avec les Européens
\item À la fin de cette leçon, je sais analyser et croiser des documents historiques (cartes, photographies, textes, données chiffrées) pour construire un raisonnement
\item À la fin de cette leçon, je sais réaliser un croquis ou une carte légendée sur le partage de l'Afrique
\item À la fin de cette leçon, je sais rédiger et justifier une démarche, une réponse ou une conclusion en mobilisant les notions de l'unité
\end{itemize}
\end{objectifs}

\begin{prerequis}
\begin{itemize}
\item La traite négrière et ses conséquences sur l'Afrique (classe de 4e)
\item Les grandes explorations européennes de l'Afrique au XIXe siècle (Livingstone, Stanley, Brazza)
\item La révolution industrielle en Europe et la recherche de matières premières
\item La notion de colonie et de protectorat
\item La lecture et la légende d'une carte historique
\end{itemize}
\end{prerequis}

\newpage

\coursec{L'impérialisme : définition, causes et contexte}

\noindent En 1884, alors que les chefs Duala signent le traité Germano-Douala sur les berges du Wouri, ils croient conclure un accord commercial ; quelques années plus tard, leurs terres sont administrées par une puissance étrangère. Comment l'Afrique entière, du Sénégal au Congo, est-elle passée en moins de trente ans sous domination européenne ? Pourquoi certains peuples, comme les Bamiléké ou les guerriers de Samory Touré, ont-ils résisté les armes à la main, tandis que d'autres choisissaient la collaboration ? Comprendre l'impérialisme, c'est comprendre les origines des frontières, des conflits et des inégalités qui marquent encore le Cameroun et l'Afrique d'aujourd'hui.

\courssub{Définition de l'impérialisme}

\begin{definition}[Impérialisme]
L'\cle{impérialisme} est la politique par laquelle un État étend sa domination politique, économique et militaire sur des territoires extérieurs, en y imposant son autorité, ses lois et ses intérêts, souvent au détriment des populations locales et de leur souveraineté.
\end{definition}

\noindent À la fin du \textsc{xix}\textsuperscript{e} siècle, l'impérialisme prend une forme nouvelle : le \cle{partage de l'Afrique}. Contrairement aux comptoirs côtiers de l'époque de la traite (XVe–XIXe siècle), les puissances européennes (France, Royaume-Uni, Allemagne, Belgique, Portugal, Italie, Espagne) pénètrent l'intérieur du continent pour y établir des \cle{colonies} (territoires sous administration directe) ou des \cle{protectorats} (autorité locale maintenue sous contrôle européen). Ce basculement s'explique par la convergence de causes économiques, politiques et idéologiques.

\courssub{Les causes de l'impérialisme}

\textbf{1. Les causes économiques : la révolution industrielle comme moteur.}
L'industrialisation (charbon, acier, chimie, électricité) transforme l'Europe en usine du monde. Elle crée trois besoins vitaux que l'Afrique peut combler :
\begin{itemize}
    \item \textbf{Matières premières :} caoutchouc (pneus, isolation), coton (textile), huile de palme (savon, graissage), minerais (cuivre, étain, or, diamant), bois précieux.
    \item \textbf{Débouchés commerciaux :} les marchés européens se saturant, l'Afrique devient un marché captif pour les produits manufacturés (tissus, armes, alcool, quincaillerie).
    \item \textbf{Lieux d'investissement :} les capitaux excédentaires cherchent des rendements plus élevés (chemins de fer, mines, plantations) que dans les économies nationales matures.
\end{itemize}
\noindent L'économiste \textbf{Hobson} (1902) lie l'impérialisme à la sous-consommation et à la recherche de placements pour les capitaux excédentaires ; \textbf{Lénine} (1917), dans \emph{L'Impérialisme, stade suprême du capitalisme}, en fait le stade monopoliste où l'exportation de \emph{capitaux} prime sur l'exportation de marchandises.

\textbf{2. Les causes politiques et nationalistes : la rivalité des puissances.}
L'unification de l'Allemagne (1871) et de l'Italie (1870) bouleverse l'équilibre européen. La possession de colonies devient un indicateur de \cle{puissance} et de \cle{prestige national}. « Un grand peuple doit avoir un grand empire » (Jules Ferry). Les rivalités franco-allemandes (revanche de 1870), anglo-françaises (Fachoda 1898) et russo-britanniques (Grand Jeu en Asie) se déportent en Afrique. La colonie sert aussi de base stratégique (charbonnage, port militaire) pour la marine de guerre.

\textbf{3. Les justifications idéologiques : la « mission civilisatrice » et le racisme scientifique.}
Pour légitimer la conquête auprès de l'opinion publique métropolitaine, les États produisent un discours idéologique :
\begin{itemize}
    \item \textbf{Mission civilisatrice :} apporter la « civilisation » (écriture, droit, médecine, christianisme) à des peuples jugés « arriérés » ou « enfants ». Jules Ferry : « Les races supérieures ont un devoir vis-à-vis des races inférieures. »
    \item \textbf{Racisme scientifique :} théories pseudo-scientifiques (gobinisme, darwinisme social) classant les « races » humaines en hiérarchie, justifiant la domination par la « supériorité biologique » du Blanc.
    \item \textbf{Œuvre humanitaire et religieuse :} lutte contre l'esclavage interne, traite arabe, sacrifices humains ; missions catholiques (Pères Blancs, Spiritains) et protestantes (Mission de Bâle, Société des Missions de Paris) qui précèdent souvent le drapeau.
\end{itemize}

\begin{attention}[Erreur fréquente : traite négrière vs colonisation]
Ne confondez pas la \cle{traite négrière} (XVe–XIXe siècle : commerce d'êtres humains, razzias, comptoirs côtiers, sans occupation territoriale profonde) et la \cle{colonisation} (fin XIXe–XXe siècle : occupation militaire, administration, exploitation économique systématique, frontières fixes). Ce sont deux phénomènes \emph{distincts et successifs} : la traite a vidé l'Afrique de ses forces vives ; la colonisation a structuré le territoire actuel.
\end{attention}

\courssub{Les acteurs précurseurs : explorateurs, missionnaires et traités}

Avant l'armée, ce sont les \cle{explorateurs} (Livingstone, Stanley, Brazza, Nachtigal, Barth) et les \cle{missionnaires} qui quadrillent le continent. Ils cartographient les fleuves (Congo, Niger, Nil), repèrent les ressources et signent des \cle{traité d'amitié et de commerce} ou de \cle{protectorat} avec les chefs africains. Ces traités, souvent rédigés en français, anglais ou allemand, sont mal compris par les signataires : cession de souveraineté vs simple autorisation de commerce.

\begin{exemplebox}[Le traité Germano-Douala (juillet 1884)]
Le 12 juillet 1884, le Dr \textbf{Gustav Nachtigal}, commissaire du Reich, signe avec les rois \textbf{King Akwa}, \textbf{King Bell} et \textbf{King Deïdo} (chefs Duala du Cameroun) un traité plaçant leurs territoires sous la « protection » de l'Allemagne. Les chefs pensent garantir leurs droits commerciaux et leur autonomie ; l'Allemagne y voit la base juridique de sa colonie du \emph{Kamerun}. Le texte stipule : « Les rois cèdent leurs droits souverains... l'Empire allemand étend sa protection sur le pays. » L'asymétrie linguistique et juridique est totale.
\end{exemplebox}

\courssub{La Conférence de Berlin (novembre 1884 – février 1885) : le partage légalisé}

\begin{definition}[Conférence de Berlin (1884-1885)]
Réunion des puissances européennes (14 États dont Allemagne, France, Royaume-Uni, Portugal, Belgique, États-Unis) convoquée par Bismarck pour fixer les \cle{règles du partage de l'Afrique} sans aucune présence africaine. Elle pose le cadre juridique de la colonisation.
\end{definition}

\noindent Les décisions majeures (Acte général de Berlin, 26 février 1885) :
\begin{enumerate}
    \item \textbf{Principe de l'occupation effective :} un traité signé sur la côte ne suffit plus ; il faut une administration, des troupes, des postes intérieurs pour revendiquer l'arrière-pays (\emph{hinterland}).
    \item \textbf{Liberté de navigation :} sur le \textbf{Congo} et le \textbf{Niger} (et leurs affluents) pour tous les pavillons.
    \item \textbf{Libre-échange :} dans le bassin conventionnel du Congo (vaste zone d'Afrique centrale).
    \item \textbf{Interdiction de la traite des Noirs :} engagement moral (peu appliqué sur le terrain).
    \item \textbf{Notification :} toute nouvelle occupation côtière doit être notifiée aux autres signataires.
\end{enumerate}

\begin{savaistu}[Le grand absent de Berlin]
\emph{Aucun Africain n'a été invité à la conférence de Berlin}, pourtant le sort de tout un continent s'y est joué. Les frontières tracées à la règle (lignes droites, méridiens, parallèles) coupent des peuples, regroupent des ennemis, ignorent les bassins versants et les aires culturelles. Le Cameroun actuel, le Nigeria, la Centrafrique, le Tchad portent encore ces cicatrices géométriques.
\end{savaistu}

\courssub{Chronologie repère : le partage de l'Afrique (1870–1914)}

\begin{popfigure}
\begin{tikzpicture}[
    timeline/.style={very thick, popblue},
    event/.style={circle, fill=popblue, inner sep=2pt},
    date/.style={font=\small, text=popdark, anchor=north},
    desc/.style={font=\small, text=popmuted, anchor=north, align=center, text width=2.5cm},
    scale=0.9
]
% Ligne principale
\draw[timeline] (0,0) -- (15,0);
% Graduations et dates
\foreach \x/\year in {0/1870, 2.14/1875, 4.29/1880, 6.43/1885, 8.57/1890, 10.71/1895, 12.86/1900, 15/1914} {
    \draw (\x,0.1) -- (\x,-0.1);
    \node[date] at (\x,-0.3) {\year};
}
% Événements clés
\node[event] at (0,0) {};
\node[desc] at (0,0.8) {1870 : 10\% colonisé};
\node[event] at (3.2,0) {};
\node[desc, align=center] at (3.2,0.8){Explorations\\ (Livingstone,\\ Stanley, Brazza)};
\node[event] at (6.43,0) {};
\node[desc, text width=3cm] at (6.43,0.8) {Nov 1884 -- Fév 1885 :\\ Conférence de Berlin};
\node[event, fill=poporange] at (5.8,0) {};
\node[desc, text=poporange, align=center] at (5.8,-0.8){Juillet 1884 :\\ Traité Germano-Douala};
\node[event] at (9,0) {};
\node[desc, align=center] at (9,0.8){Conquêtes militaires\\ (France, UK, Allemagne)};
\node[event, fill=poppurple] at (10.2,0) {};
\node[desc, text=poppurple, align=center] at (10.2,-0.8){1898 :\\ Crise de Fachoda};
\node[event] at (15,0) {};
\node[desc, align=center] at (15,0.8){1914 : >90\% colonisé\\ (Éthiopie, Libéria libres)};
\end{tikzpicture}
\legende{Frise chronologique : le partage de l'Afrique de 1870 à 1914. En orange : événement camerounais ; en violet : crise diplomatique majeure.}
\end{popfigure}

\noindent En 1870, les Européens ne contrôlent que les côtes (Algérie, Cap, Angola, Mozambique, Sénégal, Golfe de Guinée) : environ \textbf{10 \%} du continent. La conférence de Berlin (1884-1885) déclenche la \cle{course au partage} (\emph{Scramble for Africa}). En 1914, à la veille de la Grande Guerre, la quasi-totalité du continent est colonisée, à l'exception de l'\textbf{Éthiopie} (victoire d'Adoua, 1896) et du \textbf{Libéria} (créé par d'anciens esclaves américains).

\begin{exemplebox}[Évolution du taux de colonisation]
\begin{itemize}
    \item \textbf{1870 :} $\approx$ 10 \% du territoire africain sous domination européenne.
    \item \textbf{1885 :} $\approx$ 25 \% (début de l'occupation effective post-Berlin).
    \item \textbf{1914 :} $\approx$ 90 \% (seuls l'Éthiopie et le Libéria échappent au partage).
\end{itemize}
Cette bascule fulgurante en 30 ans illustre la violence et la rapidité de l'impérialisme.
\end{exemplebox}

\begin{popfigure}
\begin{tikzpicture}[scale=0.85]
% Styles
\tikzset{
    pie slice/.style={thick, draw=white, line width=1pt},
    label/.style={font=\small, text=popdark},
    year label/.style={font=\normalsize\bfseries, text=popblue, yshift=-1.8cm}
}
% Données : {année, %colonisé, %libre, couleur colonisé, couleur libre}
\foreach \year/\col/\lib/\ccol/\lcol [count=\i from 0] in {1870/10/90/popblue/popgreen, 1885/25/75/poporange/popgreen, 1914/90/10/poppurple/popgreen} {
    \begin{scope}[xshift={\i*5.5cm}]
        % Camembert
        \draw[pie slice, fill=\ccol] (0,0) -- (0:1.5) arc (0:{\col*3.6}:1.5) -- cycle;
        \draw[pie slice, fill=\lcol] (0,0) -- ({\col*3.6}:1.5) arc ({\col*3.6}:360:1.5) -- cycle;
        % Légende année
        \node[year label] at (0,0) {\year};
        % Légendes parts (si place)
        \ifnum\col>15
            \pgfmathsetmacro{\midangle}{\col*3.6/2}
            \node[label, text=\ccol] at (\midangle:0.7) {\col\%};
        \fi
        \ifnum\lib>15
            \pgfmathsetmacro{\midangle}{\col*3.6 + \lib*3.6/2}
            \node[label, text=\lcol] at (\midangle:0.7) {\lib\%};
        \fi
    \end{scope}
}
% Légende globale
\begin{scope}[yshift=-3.5cm, xshift=5.5cm]
    \draw[fill=popblue, draw=white] (-3,0) rectangle (-2.5,0.3) node[right, label, xshift=0.2cm] {Colonisé (1870)};
    \draw[fill=poporange, draw=white] (-3,-0.5) rectangle (-2.5,-0.2) node[right, label, xshift=0.2cm] {Colonisé (1885)};
    \draw[fill=poppurple, draw=white] (-3,-1) rectangle (-2.5,-0.7) node[right, label, xshift=0.2cm] {Colonisé (1914)};
    \draw[fill=popgreen, draw=white] (2,0) rectangle (2.5,0.3) node[right, label, xshift=0.2cm] {Indépendant};
\end{scope}
\end{tikzpicture}
\legende{Évolution de la part de l'Afrique sous domination européenne (camemberts 1870, 1885, 1914). Source : estimations historiographiques classiques.}
\end{popfigure}

\courssub{Exercice résolu : analyser l'acte de la Conférence de Berlin}

\begin{exoresolu}[Analyse d'un extrait de l'Acte général de Berlin (26 février 1885)]
\textbf{Document :} « Les Puissances signataires... déclarent que, dans les territoires... elles veilleront à la conservation des populations indigènes et à l'amélioration de leurs conditions morales et matérielles... Elles s'engagent à protéger et favoriser, sans distinction de nationalité ni de religion, les institutions et entreprises religieuses, scientifiques ou charitables... La navigation sur le Niger... est et restera entièrement libre pour les navires de commerce de toutes les nations. »

\textbf{Questions :}
\begin{enumerate}
    \item Identifiez la nature et la date de ce document.
    \item Relevez deux principes juridiques posés par la Conférence.
    \item Montrez en quoi ce texte masque la réalité de la conquête.
\end{enumerate}

\tcblower

\textbf{Solution détaillée :}
\begin{enumerate}
    \item \textbf{Nature :} Extrait de l'\emph{Acte général de la Conférence de Berlin}, signé le 26 février 1885. C'est un texte diplomatique international (droit international public).
    \item \textbf{Deux principes :}
        \begin{itemize}
            \item \textbf{Libre-échange et liberté de navigation :} « La navigation sur le Niger... est et restera entièrement libre pour les navires de commerce de toutes les nations. » (Art. 1-2 de l'Acte).
            \item \textbf{Obligation de protection humanitaire :} « Elles veilleront à la conservation des populations indigènes... protéger et favoriser... les institutions... religieuses, scientifiques ou charitables. » (Art. 6 : mission civilisatrice, lutte contre l'esclavage).
        \end{itemize}
    \item \textbf{Masquage de la réalité :}
        \begin{itemize}
            \item Le texte proclame la « conservation des populations » et le « bien-être matériel » alors que la conquête s'accompagne de guerres de pacification, de travail forcé, de prélèvements fiscaux lourds et de spoliation foncière.
            \item La « liberté de navigation » et le « libre-échange » servent les intérêts des métropoles (écoulement des matières premières, entrée des produits manufacturés sans droits de douane) et empêchent les Africains de protéger leurs économies locales.
            \item L'absence de signataires africains et l'invocation de la « civilisation » légitiment \emph{a posteriori} le partage territorial décidé entre Européens (principe de l'occupation effective).
        \end{itemize}
\end{enumerate}
\end{exoresolu}

\courssub{Synthèse de la section}

\begin{aretenir}[L'impérialisme en Afrique : les clés de compréhension]
\begin{itemize}
    \item \textbf{Définition :} Politique d'expansion et de domination (politique, économique, militaire) d'États industrialisés sur des territoires africains.
    \item \textbf{Causes :} Économiques (matières premières, débouchés, capitaux), Politiques (rivalités, prestige), Idéologiques (mission civilisatrice, racisme scientifique).
    \item \textbf{Acte fondateur :} Conférence de Berlin (1884-1885) : occupation effective, libre navigation Congo/Niger, libre-échange, interdiction traite. \textbf{Zéro Africain invité.}
    \item \textbf{Chronologie :} 1870 (10\%) $\to$ Berlin (règles) $\to$ 1884-1900 (conquêtes militaires) $\to$ 1914 (90\%+, Éthiopie et Libéria libres).
    \item \textbf{Conséquence immédiate :} L'Afrique est découpée en colonies aux frontières artificielles, administrées par la France, le Royaume-Uni, l'Allemagne, la Belgique, le Portugal, l'Italie, l'Espagne.
\end{itemize}
\end{aretenir}

\noindent Cette mise en place brutale de la domination coloniale ne s'est pas faite sans heurts. La section suivante étudiera comment les puissances européennes — au premier chef la France et le Royaume-Uni — ont conduit leurs conquêtes militaires sur le terrain, avant d'analyser les réactions africaines : résistances armées et collaborations.

\coursec{Les conquêtes françaises}

La section précédente a posé le cadre général de l'impérialisme européen : la rivalité entre puissances industrielles, la conférence de Berlin (1884-1885) et le principe de l'occupation effective. Nous allons maintenant examiner comment la France, forte d'une présence ancienne sur quelques points côtiers, a mis en œuvre ce principe pour construire le deuxième empire colonial africain, en suivant une logique de continuité territoriale qui la mènera du Sénégal au Tchad, et de l'Algérie au Congo.

\courssub{Les points de départ : les comptoirs du Sénégal et la base algérienne}

Avant le partage de l'Afrique, la présence française se limite à deux zones distinctes. Au \textbf{Sénégal}, la France possède depuis le XVIIe siècle deux comptoirs insulaires : \textbf{Saint-Louis} (embouchure du fleuve Sénégal) et \textbf{Gorée} (au large de Dakar). Ces points d'appui servent au commerce de la gomme arabique, puis de l'arachide, et constituent la tête de pont pour la pénétration vers le Soudan (l'actuel Mali) par la vallée du fleuve Sénégal puis du Niger.

Au nord, l'\textbf{Algérie} constitue un cas à part : conquise militairement à partir de \textbf{1830}, elle est administrée comme une extension de la métropole (départements français) et non comme une colonie de peuplement classique. Elle offre une base logistique massive pour la poussée vers le sud, en direction du Sahara et du Soudan central.

\begin{definition}[L'AOF et l'AEF : deux fédérations coloniales]
L'\textbf{Afrique occidentale française (AOF)}, créée par décret en \textbf{1895}, regroupe les colonies du Sénégal, de la Guinée, de la Côte d'Ivoire, du Dahomey (Bénin), du Soudan français (Mali), de la Mauritanie, du Niger et de la Haute-Volta (Burkina Faso). Son gouverneur général réside à Dakar.
L'\textbf{Afrique équatoriale française (AEF)}, constituée en \textbf{1910}, fédère le Gabon, le Moyen-Congo (Congo-Brazzaville), l'Oubangui-Chari (République centrafricaine) et le Tchad. Son gouverneur général réside à Brazzaville.
Ces fédérations visent à rationaliser l'administration, la justice et les finances entre territoires contigus.
\end{definition}

\courssub{La conquête de l'Afrique occidentale française (AOF) : la poussée vers le Niger}

La conquête de l'AOF suit l'axe fluvial \textbf{Sénégal $\rightarrow$ Niger}. Elle est l'œuvre de militaires-administrateurs qui agissent souvent en toute autonomie par rapport à Paris (« conquête par les capitaines »).

\begin{itemize}
    \item \textbf{Louis Faidherbe (gouverneur du Sénégal, 1854-1865)} : Il transforme le comptoir en colonie de mise en valeur. Il construit le chemin de fer Dakar-Saint-Louis, pacifie le Cayor et le Walo, et lance les premières expéditions vers le haut Sénégal (Bakel, Médine). Il pose les bases de l'administration indigène (chefs de canton).
    \item \textbf{Joseph Gallieni (commandant supérieur du Soudan français, 1880-1881 ; gouverneur, 1886-1888)} : Théoricien de la « méthode des taches d'huile » (occupation progressive par postes militaires reliés par des routes), il progresse sur le Niger, fonde \textbf{Bamako} (1883) et signe des traités avec les chefs locaux (Bamako, Kita). Il soumet le royaume du Khasso et repousse les Toucouleur d'Ahmadou Sékou.
    \item \textbf{Louis Archinard (1888-1893)} : Il mène la guerre ouverte contre l'empire toucouleur d'Ahmadou Sékou (fils d'El Hadj Oumar). Il s'empare de \textbf{Ségou} (1890) et de \textbf{Nioro} (1891), brisant la résistance la plus structurée du Soudan occidental.
\end{itemize}

La progression se heurte ensuite à l'empire de \textbf{Samory Touré} (voir section résistances) et aux Touaregs au nord. La « mission Voulet-Chanoine » (1898-1899), d'une violence extrême, relie enfin le Soudan au Tchad, bouclant la jonction AOF-AEF.

\begin{methode}[Analyser une carte de conquête coloniale]
Pour lire une carte montrant la progression coloniale (comme la Figure 1 ci-dessous) :
\begin{enumerate}
    \item \textbf{Identifier les points de départ} : ici Saint-Louis, Gorée, Alger, Brazzaville (cercles rouges ou carrés).
    \item \textbf{Repérer les axes de progression} : flèches suivant les fleuves (Sénégal, Niger, Chari, Oubangui) ou les pistes caravanières.
    \item \textbf{Noter les dates clés} sur les flèches ou les postes fondés (ex. : Bamako 1883, Tombouctou 1894, Zinder 1899).
    \item \textbf{Localiser les obstacles} : zones d'ombres (résistances Samory, Behanzin, Touareg), frontières rivales (britanniques au Niger, allemands au Cameroun/Togo).
    \item \textbf{Visualiser l'objectif stratégique} : la ligne continue Dakar-Djibouti (traits mixtes) vs la réalité 1914.
\end{enumerate}
\end{methode}

\begin{popfigure}
\begin{tikzpicture}[scale=0.9, every node/.style={font=\small}]
    % Coastline simplified (Gulf of Guinea)
    \draw[popblue, thick] (0,0) .. controls (1,1) and (3,2) .. (5,2.5) .. controls (6,3) and (7,2) .. (8,1) .. controls (9,0) and (10,-1) .. (10,-2) .. controls (9,-3) and (7,-3) .. (5,-2.5) .. controls (3,-2) and (1,-1) .. (0,0);
    % Rivers
    \draw[popblue!50, thick] (1.5,0.5) .. controls (2.5,1) and (4,1.5) .. (5.5,1.8); % Senegal River
    \draw[popblue!50, thick] (3.5,-0.5) .. controls (4.5,0) and (5.5,0.5) .. (6.5,1); % Niger upper
    \draw[popblue!50, thick] (6.5,1) .. controls (7,0.5) and (7.5,-0.5) .. (8.5,-1.5); % Niger loop
    \draw[popblue!50, thick] (2.5,-1.5) .. controls (4,-1) and (5,-0.5) .. (6,0); % Chari/Logone approx
    
    % Key Cities / Posts AOF (Red)
    \node[circle, fill=popred, inner sep=1.5pt, label=above left:\textcolor{popdark}{Saint-Louis}] (SL) at (1.2,0.8) {};
    \node[circle, fill=popred, inner sep=1.5pt, label=below:\textcolor{popdark}{Gorée/Dakar}] (GK) at (2.0,0.0) {};
    \node[circle, fill=popred, inner sep=1.5pt, label=above:\textcolor{popdark}{Bakel}] (BK) at (2.8,1.0) {};
    \node[circle, fill=popred, inner sep=1.5pt, label=above:\textcolor{popdark}{Médine}] (MD) at (3.2,1.2) {};
    \node[circle, fill=popred, inner sep=1.5pt, label=above:\textcolor{popdark}{Kita}] (KT) at (4.2,1.5) {};
    \node[circle, fill=popred, inner sep=1.5pt, label=above:\textcolor{popdark}{Bamako}] (BM) at (5.0,1.7) {};
    \node[circle, fill=popred, inner sep=1.5pt, label=above:\textcolor{popdark}{Ségou}] (SG) at (6.0,1.5) {};
    \node[circle, fill=popred, inner sep=1.5pt, label=above:\textcolor{popdark}{Nioro}] (NR) at (6.5,1.2) {};
    \node[circle, fill=popred, inner sep=1.5pt, label=above:\textcolor{popdark}{Tombouctou}] (TM) at (7.5,1.0) {};
    \node[circle, fill=popred, inner sep=1.5pt, label=above:\textcolor{popdark}{Gao}] (GA) at (8.0,0.2) {};
    
    % AEF axis from South (Green)
    \node[circle, fill=popgreen, inner sep=1.5pt, label=left:\textcolor{popdark}{Brazzaville}] (BZ) at (3.5,-2.2) {};
    \node[circle, fill=popgreen, inner sep=1.5pt, label=right:\textcolor{popdark}{Bangui}] (BG) at (5.5,-1.8) {};
    \node[circle, fill=popgreen, inner sep=1.5pt, label=right:\textcolor{popdark}{Fort-Lamy}] (FL) at (7.0,-1.0) {};
    \node[circle, fill=popgreen, inner sep=1.5pt, label=above:\textcolor{popdark}{Zinder}] (ZD) at (8.0,-0.5) {};
    
    % Arrows AOF (Red)
    \draw[popred, thick, ->, >=Stealth] (SL) -- (BK) node[midway, above, sloped, xshift=-2pt] {\tiny 1855-65 Faidherbe};
    \draw[popred, thick, ->, >=Stealth] (BK) -- (KT) node[midway, above, sloped] {\tiny 1880 Gallieni};
    \draw[popred, thick, ->, >=Stealth] (KT) -- (BM) node[midway, above, sloped] {\tiny 1883};
    \draw[popred, thick, ->, >=Stealth] (BM) -- (SG) node[midway, above, sloped] {\tiny 1890 Archinard};
    \draw[popred, thick, ->, >=Stealth] (SG) -- (NR) node[midway, above, sloped] {\tiny 1891};
    \draw[popred, thick, ->, >=Stealth] (NR) -- (TM) node[midway, above, sloped] {\tiny 1894};
    \draw[popred, thick, ->, >=Stealth] (TM) -- (GA) node[midway, above, sloped] {\tiny 1899};
    \draw[popred, thick, ->, >=Stealth] (GA) -- (ZD) node[midway, below, sloped] {\tiny 1899 Voulet-Chanoine};
    
    % Arrows AEF (Green)
    \draw[popgreen, thick, ->, >=Stealth] (BZ) -- (BG) node[midway, left, sloped] {\tiny 1880-90 Brazza};
    \draw[popgreen, thick, ->, >=Stealth] (BG) -- (FL) node[midway, left, sloped] {\tiny 1900 Gentil};
    \draw[popgreen, thick, ->, >=Stealth] (FL) -- (ZD) node[midway, below, sloped] {\tiny 1900 Jonction};
    
    % Resistance Zones
    \draw[poporange, dashed, thick] (5.5,1.8) ellipse (1.2 and 0.8);
    \node[poporange, font=\tiny] at (5.5,2.7) {Zone Samory};
    \draw[poporange, dashed, thick] (2.5,-0.5) ellipse (0.8 and 0.6);
    \node[poporange, font=\tiny] at (2.5,-1.5) {Dahomey Behanzin};
    \draw[poporange, dashed, thick] (8.5,1.5) ellipse (1.5 and 1);
    \node[poporange, font=\tiny] at (8.5,2.6) {Touaregs};
    
    % Legend
    \node[draw, fill=popcream, anchor=south west, inner sep=4pt, align=center] at (0.1,-3.2){
        \begin{tabular}{ll}
            \textcolor{popred}{\rule{1em}{1pt}} Axe Sénégal-Niger (AOF) & \textcolor{popgreen}{\rule{1em}{1pt}} Axe Congo-Tchad (AEF) \\
            \textcolor{popred}{\bullet} Postes français AOF & \textcolor{popgreen}{\bullet} Postes français AEF \\
            \textcolor{poporange}{- - -} Zones de résistance majeure &
        \end{tabular}
    };
    
    % North Arrow
    \node at (9.5, -3.0) {\tikz \draw[popdark, thick, ->] (0,0) -- (0,0.5) node[above] {N};};
\end{tikzpicture}
\legende{Figure 1 : Axes de pénétration française en Afrique occidentale et équatoriale (1850-1900). En rouge, la progression AOF depuis Saint-Louis ; en vert, la progression AEF depuis Brazzaville. Les zones hachurées orange indiquent les foyers de résistance majeure (Samory, Behanzin, Touaregs).}
\end{popfigure}

\courssub{La conquête de l'Afrique équatoriale française (AEF) : l'œuvre de Brazza}

Alors que l'AOF progresse d'ouest en est, l'AEF se construit du sud vers le nord, grâce à l'explorateur \textbf{Pierre Savorgnan de Brazza}. Naturalisé français, il explore l'Ogooué et le Congo (1875-1882) avec une méthode originale : la négociation pacifique et le respect des coutumes, ce qui lui vaut le surnom de « père des Noirs ».

\begin{itemize}
    \item \textbf{Traité de Makoko (1880)} : Brazza signe un traité de protectorat avec le \textbf{roi Makoko des Téké}, sur les rives du fleuve Congo (actuel site de Brazzaville). Ce traité, ratifié par la France en 1882, légitime la présence française face aux prétentions de Stanley (au service de Léopold II de Belgique) sur la rive gauche (Congo belge).
    \item \textbf{Implantation} : Fondation de \textbf{Brazzaville} (1883), puis progression vers le nord : \textbf{Franceville} (Gabon, 1880), \textbf{Bangui} (Oubangui-Chari, 1889), \textbf{Fort-Lamy} (Tchad, 1900 par la mission Gentil).
    \item La jonction AOF-AEF est effective en 1900 à la bataille de \textbf{Kousséri} (mission Lamy, rejoint par Foureau-Lamy et Gentil), où est tué le sultan du Bornou, Rabah.
\end{itemize}

\begin{savaistu}[Pierre Savorgnan de Brazza (1852-1905)]
Explorateur d'origine italienne, naturalisé français en 1874. Contrairement à Stanley (violence, traités inégaux), Brazza privilégie le dialogue, apprend les langues, soigne les populations. Il devient commissaire général du Congo français (1886-1897). Démis de ses fonctions pour avoir dénoncé les exactions des compagnies concessionnaires (rapport Brazza, 1905), il meurt épuisé au retour d'une mission d'enquête. Son corps repose à Brazzaville (mémorial inauguré en 2006). Il incarne le mythe d'une colonisation « douce », mais ses traités ont servi la même spoliation foncière.
\end{savaistu}

\courssub{Les grandes résistances africaines face à la conquête française}

La conquête n'est pas une promenade militaire. Elle s'étale sur plus de 40 ans (1854-1914) et se heurte à des États structurés et des sociétés guerrières.

\begin{attention}[Erreur fréquente : une conquête rapide et pacifique]
\emph{Piège de rédaction} : Ne pas écrire « la France a conquis l'AOF en 1895 » (date de création de la fédération administrative). La conquête militaire commence en 1854 (Faidherbe) et s'achève vers 1914 (pacification de la Mauritanie, du Niger, du Tchad).
\emph{Réalité} : Guerres longues (Samory : 1882-1898), coûts humains élevés (missions Voulet-Chanoine, colonnes Mangin), résistance passive (fuite, refus d'impôt, culture du millet au lieu de l'arachide).
\end{attention}

Trois figures majeures illustrent cette résistance armée :

\begin{enumerate}
    \item \textbf{Samory Touré (vers 1830-1900), l'Almamy de l'empire mandingue (Wassoulou)} :
        Ancien marchand dyula, il construit un État moderne (armée permanente \emph{sofa} équipée de fusils à répétition, administration centrale, ateliers de réparation d'armes). Il pratique la \textbf{terre brûlée} et la mobilité stratégique pour éviter l'encerclement. Il signe des traités (Bissandougou 1887, 1890) pour gagner du temps, mais refuse la soumission. Traqué par les colonnes françaises (Combes, Gouraud), il est capturé à \textbf{Guémou} (Côte d'Ivoire) en \textbf{1898} et déporté au Gabon où il meurt.
    \item \textbf{Béhanzin (1844-1906), roi du Dahomey (actuel Bénin)} :
        Héritier d'un État militarisé (amazones, armée permanente), il refuse le protectorat français sur Porto-Novo (1889). Guerre de 1890-1894 (général Dodds). Malgré le courage des amazones, la supériorité d'artillerie française (canons de 80 et 95 mm) l'emporte. Abomey tombe en 1894. Béhanzin se rend en 1894, déporté en Martinique puis en Algérie.
    \item \textbf{Les Touaregs (Sahara central)} :
        Résistance diffuse, mobile, basée sur la connaissance du désert et le raid (razzia). La « pacification » dure jusqu'en 1906 (bataille de Tit) et 1917 (révolte de Kaocen au Niger). Elle illustre la difficulté de contrôler des espaces vides mais stratégiques.
\end{enumerate}

\begin{exoresolu}[Analyse d'un document : la capture de Samory (1898)]
\textbf{Document} : Extrait du rapport du capitaine Gouraud à son commandant supérieur, 29 septembre 1898. « L'Almamy Samory, abandonné de tous, n'ayant plus que 200 fusils et quelques cartouches, s'est rendu à discrétion... Il a été fait prisonnier avec sa famille et ses trésors. »
\textbf{Question} : En quoi cet extrait illustre-t-il l'échec final de la stratégie de Samory ?
\textbf{Solution détaillée} :
\begin{enumerate}
    \item \textbf{Isolation diplomatique} : « Abandonné de tous ». La politique de terre brûlée a aliéné les populations locales ; les chefs voisins (Bambara, Bobo) ont rallié les Français.
    \item \textbf{Épuisement logistique} : « 200 fusils, quelques cartouches ». L'industrie d'armement artisanale ne suffit plus face à l'approvisionnement industriel français (port de Conakry, chemin de fer).
    \item \textbf{Supériorité tactique française} : Colonnes mobiles, télégraphe, artillerie, renseignement local. La « méthode des taches d'huile » (Gallieni) a quadrillé le territoire, privant Samory de sa mobilité.
    \item \textbf{Fin de la résistance structurée} : La capture du chef charismatique entraîne la dislocation de l'empire Wassoulou, même si des foyers de guérilla persistent.
\end{enumerate}
\end{exoresolu}

\courssub{L'idée de continuité territoriale : le rêve du Dakar-Djibouti et la crise de Fachoda}

L'objectif stratégique suprême de la France est de relier ses possessions : \textbf{Alger $\rightarrow$ Sahara $\rightarrow$ AOF $\rightarrow$ AEF $\rightarrow$ Djibouti (Côte des Somalis)}. C'est le projet de \textbf{transafricain} (chemin de fer ou ligne télégraphique), symbole de la grandeur française face à l'axe britannique \textbf{Le Caire-Le Cap} (Cecil Rhodes).

\begin{itemize}
    \item \textbf{La mission Marchand (1896-1898)} : Partie de Brazzaville, elle remonte l'Oubangui, franchit la ligne de partage des eaux (bassin du Nil) et atteint \textbf{Fachoda} (Sud-Soudan actuel) en juillet 1898.
    \item \textbf{L'affrontement} : Le général britannique \textbf{Kitchener}, victorieux des Mahdistes à Omdurman (septembre 1898), descend le Nil et exige l'évacuation de Fachoda.
    \item \textbf{La reculade française} : Face à la menace de guerre navale (flotte britannique supérieure) et à l'isolement diplomatique (affaire Dreyfus affaiblit la France), le ministre Delcassé ordonne le retrait (novembre 1898).
    \item \textbf{Conséquences} : Traité de 1899 : la France renonce au Nil (sphère d'influence britannique) mais obtient la liberté d'action au Maroc et la reconnaissance de sa sphère au Sahara/Tchad. L'axe Est-Ouest est brisé ; l'axe Nord-Sud (Dakar-Djibouti) reste inachevé (chemin de fer jamais construit).
\end{itemize}

\courssub{La conquête de l'Afrique du Nord : protectorats tunisien et marocain}

Parallèlement à la poussée subsaharienne, la France achève le verrouillage du Maghreb.

\begin{itemize}
    \item \textbf{Tunisie (1881)} : Prétexte : incidents frontaliers algéro-tunisiens (Khroumirs). Expédition militaire (30 000 hommes). \textbf{Traité du Bardo (mai 1881)} : protectorat. \textbf{Convention de La Marsa (1883)} : le bey conserve le trône mais le résident général français dirige l'administration, les finances, la justice. L'Italie, rivale déçue, se rapproche de l'Allemagne (Triplice).
    \item \textbf{Maroc (1912)} : Enjeu stratégique (détroit de Gibraltar). Crises successives (Tanger 1905, Agadir 1911) résolues par conférences internationales (Algésiras 1906). \textbf{Traité de Fès (mars 1912)} : protectorat français sur la majeure partie du territoire. Zone nord (Rif) et zone sud (Tarfaïa) à l'Espagne. Tanger internationalisée. Résistance armée (Abdelkrim al-Khattabi au Rif, 1921-1926) au-delà de 1914.
\end{itemize}

\courssub{Bilan : l'empire colonial français en Afrique en 1914}

À la veille de la Grande Guerre, la France domine un ensemble territorial colossal, deuxième empire colonial après la Grande-Bretagne.

\begin{aretenir}[L'empire colonial français en chiffres (1914)]
\begin{itemize}
    \item \textbf{Superficie} : $\approx$ \textbf{10,5 millions de km²} en Afrique (sur 11,5 millions total empire).
    \item \textbf{Comparaison} : Environ \textbf{20 fois la superficie de la métropole} (551 500 km²).
    \item \textbf{Population} : $\approx$ 30 à 35 millions d'Africains (estimation).
    \item \textbf{Organisation} : 2 fédérations (AOF, AEF), 3 territoires d'Afrique du Nord (Algérie départementée, Tunisie et Maroc protectorats), 1 colonie de peuplement (Algérie), 1 territoire de la Corne de l'Afrique (Côte des Somalis/Djibouti).
    \item \textbf{Frontières} : Tracées au cordeau lors de conférences européennes (Berlin, Bruxelles, accords 1898-1899), ignorant les réalités ethniques et historiques $\rightarrow$ source de conflits post-indépendance.
\end{itemize}
\end{aretenir}

\begin{popfigure}
\begin{tikzpicture}[scale=0.85, every node/.style={font=\small}]
    % Africa outline (very simplified polygon)
    \draw[popline, thick, fill=popcream] 
        (0,0) -- (10,0) -- (11,2) -- (11,5) -- (10,7) -- (9,9) -- (7,10) -- (5,10) -- (3,9) -- (1,7) -- (0,5) -- cycle;
    
    % Mediterranean
    \draw[popblue, thick] (0,10) -- (11,10);
    \node[popblue, rotate=90] at (-0.5, 5) {Méditerranée};
    
    % Equator line
    \draw[dashed, popmuted] (0,4.5) -- (11,4.5);
    \node[popmuted, font=\tiny] at (11.2, 4.5) {Équateur};
    
    % NORTH AFRICA (Algeria, Tunisia, Morocco) - PopBlue fill
    \draw[popblue, thick, fill=popblue!20] 
        (2,8.5) -- (9,8.5) -- (9.5,7) -- (8,6) -- (6,5.5) -- (4,6) -- (2,7) -- cycle;
    \node[popdark, font=\bfseries] at (5.5, 7.2) {AFRIQUE DU NORD FRANÇAISE};
    \node[popdark, font=\tiny] at (3.5, 7.8) {Alger (Dép.)};
    \node[popdark, font=\tiny] at (5.5, 7.8) {Tunis (Prot. 1881)};
    \node[popdark, font=\tiny] at (7.5, 7.8) {Fès/Rabat (Prot. 1912)};
    
    % AOF - PopRed fill (West)
    \draw[popred, thick, fill=popred!15] 
        (2,5.5) -- (6,5.5) -- (6.5,4) -- (7,3) -- (6.5,2) -- (5,1.5) -- (3,1.5) -- (1.5,3) -- (1,4.5) -- cycle;
    \node[popdark, font=\bfseries] at (3.5, 3.5) {AOF (1895)};
    \node[popdark, font=\tiny] at (2.5, 4.5) {Sénégal};
    \node[popdark, font=\tiny] at (3.5, 3.0) {Soudan};
    \node[popdark, font=\tiny] at (5.0, 4.0) {Côte d'Ivoire};
    \node[popdark, font=\tiny] at (5.5, 2.5) {Dahomey};
    \node[popdark, font=\tiny] at (4.0, 5.0) {Mauritanie};
    
    % AEF - PopGreen fill (Central)
    \draw[popgreen, thick, fill=popgreen!15] 
        (6,5.5) -- (9,5.5) -- (9.5,4) -- (9,2.5) -- (7,1.5) -- (5.5,1.5) -- (5,2.5) -- (5.5,3.5) -- cycle;
    \node[popdark, font=\bfseries] at (7.5, 3.5) {AEF (1910)};
    \node[popdark, font=\tiny] at (6.5, 4.5) {Tchad};
    \node[popdark, font=\tiny] at (8.0, 4.0) {Oubangui-Chari};
    \node[popdark, font=\tiny] at (8.5, 2.5) {Moyen-Congo};
    \node[popdark, font=\tiny] at (7.0, 2.0) {Gabon};
    
    % Djibouti - PopPurple
    \draw[poppurple, thick, fill=poppurple!15] (9.5, 4.5) -- (10.5, 4.5) -- (10.5, 4) -- (9.5, 4) -- cycle;
    \node[poppurple, font=\tiny, rotate=90] at (10.7, 4.25) {Côte des Somalis};
    
    % Legend Box
    \node[draw=popdark, fill=white, anchor=south east, inner sep=5pt, rounded corners, align=center] at (10.8, -0.2){
        \begin{tabular}{ll}
            \cellcolor{popblue!20}\textcolor{popblue}{\rule{1em}{1pt}} Afrique du Nord (Algérie, Tunisie, Maroc) \\
            \cellcolor{popred!15}\textcolor{popred}{\rule{1em}{1pt}} Afrique Occidentale Française (AOF) \\
            \cellcolor{popgreen!15}\textcolor{popgreen}{\rule{1em}{1pt}} Afrique Équatoriale Française (AEF) \\
            \cellcolor{poppurple!15}\textcolor{poppurple}{\rule{1em}{1pt}} Côte des Somalis (Djibouti) \\
            \multicolumn{2}{l}{\tiny \textit{Total Afrique : $\approx$ 10,5 millions km²}}
        \end{tabular}
    };
    
    % North Arrow
    \node at (0.5, 9.5) {\tikz \draw[popdark, thick, ->] (0,0) -- (0,0.5) node[above] {N};};
\end{tikzpicture}
\legende{Figure 2 : Croquis schématique de l'empire colonial français en Afrique en 1914. Les couleurs distinguent les quatre ensembles administratifs : Afrique du Nord (bleu), AOF (rouge), AEF (vert), Côte des Somalis (violet). Les frontières intérieures sont indicatives.}
\end{popfigure}

La conquête française achève ainsi le partage du continent. Mais l'empire ainsi constitué porte en lui les germes de sa dislocation : frontières artificielles, économies d'extraction (chemin de fer du Niger, cultures de rente), administration directe ou indirecte qui brise les structures traditionnelles, et naissance d'une élite occidentale (écoles de William Ponty, école de médecine de Dakar) qui réclamera, dès l'entre-deux-guerres, l'égalité des droits et l'indépendance. La section suivante analysera la conquête britannique et celles des puissances secondaires (Allemagne, Belgique, Italie, Portugal), avant d'étudier en détail les réactions africaines : résistances et collaborations.

\coursec{Les conquêtes anglaises et des autres puissances}

Après avoir analysé la stratégie française, fondée sur une expansion est-ouest (Dakar-Djibouti) et la constitution de deux grands ensembles (AOF et AEF), nous nous tournons maintenant vers le principal rival de la France : la Grande-Bretagne. La logique impériale britannique est radicalement différente : elle vise une continuité nord-sud, du Cap au Caire, pour sécuriser la route des Indes et contrôler le Nil. Parallèlement, l'Allemagne, la Belgique, l'Italie, le Portugal et l'Espagne complètent le découpage du continent, aboutissant en 1914 à un partage quasi total de l'Afrique.

\courssub{Les ambitions britanniques : l'axe Le Cap – Le Caire}

\cle{Cecil Rhodes} (1853-1902), magnat des diamants et premier ministre de la colonie du Cap, incarne l'impérialisme britannique le plus agressif. Son rêve : « peindre la carte en rouge » (couleur de l'Empire britannique sur les cartes) du \cle{Cap} (Afrique du Sud) au \cle{Caire} (Égypte). Ce projet repose sur deux piliers : le contrôle du canal de Suez (voie vitale vers l'Inde) et la construction d'un \cle{chemin de fer transafricain} reliant Le Cap au Caire, symbole d'unité impériale et outil d'exploitation économique (projet jamais achevé dans sa totalité).

\begin{popfigure}
\begin{tikzpicture}[scale=0.75]
    % Contour très simplifié de l'Afrique (polygone grossier)
    \draw[popline, thick, fill=popcream!30] 
        (0,0) -- (1,2) -- (2,5) -- (4,7) -- (6,6) -- (8,5) -- (9,3) -- (8,1) -- (6,-1) -- (4,-2) -- (2,-1) -- cycle;
    
    % Axe Britannique : Le Cap -> Le Caire (Vertical, Rouge)
    \draw[popred, very thick, ->] (4.5, -1.5) -- (4.5, 6.5) node[midway, right, popred, font=\bfseries] {Axe Le Cap -- Le Caire};
    % Points clés axe britannique
    \foreach \y/\name in {-1/Le Cap, 1/Kimberley, 2.5/Rhodésies, 4/Lacs, 5.5/Fachoda, 6.5/Le Caire} {
        \fill[popred] (4.5, \y) circle (3pt);
        \node[right, font=\small] at (4.7, \y) {\name};
    }
    
    % Axe Français : Dakar -> Djibouti (Horizontal, Bleu)
    \draw[popblue, very thick, ->] (-1, 2.5) -- (9.5, 2.5) node[midway, above, popblue, font=\bfseries] {Axe Dakar -- Djibouti};
    % Points clés axe français
    \foreach \x/\name in {-0.5/Dakar, 2/Bamako, 4.5/Fachoda, 7/Addis, 9.5/Djibouti} {
        \fill[popblue] (\x, 2.5) circle (3pt);
        \node[below, font=\small] at (\x, 2.3) {\name};
    }
    
    % Fachoda : Point de croisement / Conflit
    \node[draw=poporange, thick, circle, fill=poporange!20, font=\bfseries\small, align=center] at (4.5, 2.5) {FACHODA\\1898};
    
    % Légende
    \node[anchor=south west, fill=white, draw=popline, rounded corners, inner sep=4pt, align=center] at (0.5, -2.5){
        \begin{tabular}{ll}
        \textcolor{popred}{\rule{1cm}{2pt}} & Axe britannique (Nord-Sud) \\
        \textcolor{popblue}{\rule{1cm}{2pt}} & Axe français (Est-Ouest) \\
        \textcolor{poporange}{\large $\bullet$} & Fachoda : crise diplomatique 1898
        \end{tabular}
    };
\end{tikzpicture}
\legende{Schéma des deux axes rivaux de pénétration coloniale en Afrique. Le croisement à Fachoda (Soudan) faillit provoquer une guerre entre la France et la Grande-Bretagne en 1898.}
\end{popfigure}

La réalisation de cet axe se heurte à des obstacles majeurs : les républiques boers (Transvaal, État libre d'Orange) au sud, la présence française au Soudan (Fachoda) au centre, et la nécessité de contrôler l'Égypte au nord.

\textbf{Les étapes de la conquête britannique.}\par
\begin{itemize}
    \item \textbf{Égypte (1882) :} Prétexte : protection des actionnaires européens du canal de Suez et rétablissement de l'ordre après la révolte d'Arabi Pacha. L'armée britannique occupe le pays ; l'Égypte devient un \cle{protectorat} officiel en 1914, mais est de fait contrôlée par Londres dès 1882 (Lord Cromer, consul général).
    \item \textbf{Soudan (1898) :} Après l'échec de Gordon à Khartoum (1885) face au Mahdi, Kitchener reconquiert le pays. La victoire d'\cle{Omdurman} (1898) et l'incident de \cle{Fachoda} (mars 1898), où la mission Marchand (française) et l'armée de Kitchener (britannique) se font face, scellent le sort du Soudan : il devient un condominium anglo-égyptien (1899), sous contrôle effectif britannique.
    \item \textbf{Afrique de l'Ouest :} \cle{Nigeria} (compagnie royale du Niger, protectorat 1900, colonie 1914), \cle{Gold Coast} (Ghana, colonie de peuplement côtière étendue à l'intérieur, 1874-1901), \cle{Sierra Leone} (colonie de la Couronne depuis 1808), \cle{Gambie}.
    \item \textbf{Afrique de l'Est :} \cle{Kenya} (protectorat 1895, colonie 1920), \cle{Ouganda} (protectorat 1894), \cle{Zanzibar} (protectorat 1890).
    \item \textbf{Afrique australe :} \cle{Rhodésie du Sud} (Zimbabwe) et \cle{Rhodésie du Nord} (Zambie) concédées à la British South Africa Company (Rhodes) dans les années 1890. \cle{Afrique du Sud} : guerre des Boers (1899-1902) pour l'or du Witwatersrand.
\end{itemize}

\begin{definition}[La guerre des Boers (1899-1902)]
Conflit opposant la \textbf{Grande-Bretagne} aux deux républiques indépendantes fondées par des fermiers d'origine hollandaise (\cle{Boers} ou Afrikaners) : le \cle{Transvaal} (République sud-africaine) et l'\cle{État libre d'Orange}. 
\textbf{Causes :} découverte d'or au Witwatersrand (1886), afflux de chercheurs britanniques (\emph{Uitlanders}), refus des Boers de leur accorder des droits civiques, volonté britannique de contrôler l'ensemble de l'Afrique australe.
\textbf{Déroulement :} guerre conventionnelle (1899-1900) puis guérilla boer et terre brûlée britannique (camps de concentration pour civils, ~28 000 morts). 
\textbf{Conséquences :} traité de Vereeniging (1902), annexion des deux républiques. Union d'Afrique du Sud (1910) : dominion autonome au sein de l'Empire, dominé par la minorité blanche.
\end{definition}

\begin{methode}[Comparer deux empires coloniaux (ex. France vs Grande-Bretagne)]
Pour analyser et comparer les empires coloniaux français et britannique en Afrique, appliquez ces quatre critères :
\begin{enumerate}
    \item \textbf{Superficie et continuité territoriale :} L'empire français est vaste mais fragmenté (AOF, AEF, Maghreb). L'empire britannique vise la continuité (axe Cap-Caire) mais possède des blocs disjoints (Afrique de l'Ouest, Afrique de l'Est, Afrique australe, Égypte).
    \item \textbf{Population et peuplement :} La France privilégie la \cle{mission civilisatrice} et l'assimilation (peuplement limité hors Algérie). La Grande-Bretagne pratique l'\cle{association} / \cle{indirect rule} (Lugard) et encourage le peuplement de peuplement dans les zones tempérées (Kenya, Rhodésies, Afrique du Sud).
    \item \textbf{Modes d'administration :} \textbf{France :} administration directe, centralisée, cadres français, \cle{Code de l'indigénat}. \textbf{Grande-Bretagne :} \emph{Indirect Rule} (gouvernance par les chefs traditionnels sous contrôle britannique), droit coutumier maintenu, moindre présence administrative \textbf{britannique}.
    \item \textbf{Intérêts économiques stratégiques :} France : matières premières (coton, arachide, minerais), débouchés, prestige. Grande-Bretagne : route des Indes (Suez, Cap), or/diamants (Afrique du Sud), huile de palme, coton (Égypte, Soudan), chemin de fer transafricain.
\end{enumerate}
\end{methode}

\courssub{Les conquêtes des autres puissances européennes}

Si la France et la Grande-Bretagne se taillent la part du lion, la Conférence de Berlin (1884-1885) légitime les prétentions des puissances secondaires.

\textbf{L'Allemagne : un empire colonial tardif mais vaste.}\par
Arrivée tardivement (Bismarck réticent jusqu'en 1884), l'Allemagne obtient quatre territoires lors du « partage de l'Afrique » :
\begin{itemize}
    \item \cle{Togo} et \cle{Cameroun (Kamerun)} (Golfe de Guinée, 1884) : traités de protectorat signés par Gustav Nachtigal avec les rois locaux (Duala au Cameroun).
    \item \cle{Afrique du Sud-Ouest} (Namibie actuelle, 1884) : colonie de peuplement, théâtre du premier génocide du XXe siècle contre les Herero et Nama (1904-1908).
    \item \cle{Afrique orientale allemande} (Tanganyika, Rwanda, Burundi, 1885) : conquête violente (révolte Maji Maji 1905-1907).
\end{itemize}

\begin{attention}[Erreur fréquente : l'appartenance du Cameroun avant 1916]
Ne confondez pas la période coloniale allemande (1884-1916) et le mandat français (1919-1960). \textbf{Avant 1916, le Kamerun est une colonie allemande.} Ce n'est qu'à la faveur de la Première Guerre mondiale (conquête franco-britannique 1914-1916) puis du traité de Versailles (1919) que le territoire est partagé en mandats de la SDN (France : 4/5, Grande-Bretagne : 1/5). Au Probatoire, situez toujours l'acte de colonisation initiale (traité de 1884) sous l'autorité allemande.
\end{attention}

\textbf{La Belgique et l'État indépendant du Congo : la propriété d'un roi.}\par
C'est l'affaire personnelle du roi \cle{Léopold II}. Sous couvert d'humanitaire (Conférence de Berlin, Association internationale africaine), il fait reconnaître sa propriété privée sur le bassin du Congo (1885). L'\cle{État indépendant du Congo} (EIC) n'est pas une colonie belge, mais la propriété du roi.

\begin{exemplebox}[Chiffres clés : l'État indépendant du Congo]
\begin{itemize}
    \item \textbf{Superficie :} ~2 345 000 km$^2$ (environ \textbf{76 fois la taille de la Belgique}).
    \item \textbf{Économie :} extraction forcée du \cle{caoutchouc} (pneus, isolation électrique) et de l'ivoire.
    \item \textbf{Méthodes :} travail forcé, prises d'otages (femmes/enfants), mutilations (mains coupées) pour punir le non-respect des quotas, armée privée (\emph{Force Publique}).
    \item \textbf{Bilan humain :} les historiens estiment la perte démographique entre \textbf{5 et 10 millions de morts} (exactions, famines, maladies) sur 20 ans (1885-1908).
    \item \textbf{Fin :} scandale international (rapport Casement, campagne de Morel et Conan Doyle) $\rightarrow$ annexion par l'État belge en \textbf{1908} : naissance du \cle{Congo belge}.
\end{itemize}
\end{exemplebox}

\textbf{Le Portugal : les vieilles possessions consolidées.}\par
Présent depuis le XVe siècle (comptoirs), le Portugal voit ses revendications intérieures (carte rose : Angola-Mozambique) réduites par l'ultimatum britannique de 1890. Il conserve :
\begin{itemize}
    \item \cle{Angola} (Afrique australe occidentale) : café, diamants, travail forcé (\cle{chibalo}).
    \item \cle{Mozambique} (Afrique australe orientale) : main-d'œuvre pour les mines sud-africaines, chemin de fer vers le Transvaal.
    \item \cle{Guinée-Bissau}, \cle{Cap-Vert}, \cle{Sao Tomé-et-Principe} (îles, cacao).
\end{itemize}
L'administration est directe, centralisée, marquée par le \cle{statut de l'indigénat} et l'assimilation très limitée (\cle{assimilados}).

\textbf{L'Italie : la puissance mineure ambitieuse.}\par
Unifiée tardivement (1861), l'Italie cherche sa « place au soleil ».
\begin{itemize}
    \item \cle{Érythrée} (1882, Massawa) et \cle{Somalie} (1889, Benadir) : possessions côtières arides.
    \item \textbf{Échec d'Adoua (1896) :} tentative de conquête de l'\cle{Éthiopie} (seul État africain chrétien indépendant). L'armée italienne (17 000 hommes) est anéantie par le \cle{Négus Ménélik II} (100 000 hommes, armement moderne). Traité d'Addis-Abeba : reconnaissance de l'indépendance éthiopienne.
    \item \cle{Libye} (Tripolitaine, Cyrénaïque, Fezzan) : arrachée à l'Empire ottoman en 1911-1912 (guerre italo-turque), pacification sanglante (Omar Mukhtar) jusqu'aux années 1930.
\end{itemize}

\begin{savaistu}[Adoua (1896) : symbole panafricain]
La bataille d'\cle{Adoua} (1er mars 1896) est la \textbf{seule victoire décisive d'un État africain sur une puissance européenne} lors du partage colonial. Le Négus Ménélik II a su unifier les seigneurs féodaux éthiopiens, se doter d'armes modernes (fusils russes, canons) et jouer des rivalités européennes. Cette victoire a préservé la souveraineté de l'Éthiopie et inspire encore aujourd'hui les mouvements panafricains et la fierté africaine (drapeau vert-jaune-rouge adopté par plusieurs États à l'indépendance en hommage à l'Éthiopie).
\end{savaistu}

\textbf{L'Espagne : des miettes stratégiques.}\par
Puissance déclinante, l'Espagne obtient des territoires marginaux lors de la Conférence de Berlin et des accords franco-espagnols (1900, 1904, 1912) :
\begin{itemize}
    \item \cle{Rio Muni} (Guinée équatoriale continentale) et \cle{Fernando Poo} (île de Bioko) : cacao, bois.
    \item \cle{Sahara occidental} (Rio de Oro, Saguia el-Hamra) : phosphates découverts plus tard.
    \item \cle{Nord du Maroc} (zone du Rif, protectorat 1912) : zone d'influence réservée par la France à l'Espagne lors du partage du Maroc.
\end{itemize}

\courssub{Bilan cartographique : le partage de l'Afrique en 1914}

À la veille de la Première Guerre mondiale, le continent est presque entièrement découpé. Seuls deux États conservent leur souveraineté internationale : le \cle{Libéria} (fondé par d'anciens esclaves américains, indépendance 1847, sous influence économique US/Firestone) et l'\cle{Éthiopie} (Empire chrétien millénaire, vainqueur d'Adoua).

\begin{popfigure}
\begin{tikzpicture}[scale=0.65]
    % Contour Afrique simplifié
    \draw[popline, thick, fill=popcream!20] 
        (0,0) -- (0.5,2) -- (1.5,5) -- (3.5,7) -- (6,6.5) -- (8,5) -- (9,3) -- (8.5,1) -- (7,-1) -- (5,-2) -- (3,-1.5) -- (1,-0.5) -- cycle;
    
    % Définition des couleurs par puissance (blocs rectangulaires approximatifs)
    % FRANCE (popblue) : AOF (Ouest), AEF (Centre), Maghreb (Nord), Madagascar (Est - hors cadre simplifié mais on met un point)
    \fill[popblue!40] (0.5,1) -- (2,3) -- (3,4) -- (2,5) -- (1,4) -- (0.5,2.5) -- cycle; % AOF approx
    \fill[popblue!40] (3,3) -- (5,4) -- (4.5,5.5) -- (2.5,5) -- cycle; % AEF approx
    \fill[popblue!40] (2,6) -- (4,6.5) -- (4,7) -- (2,6.5) -- cycle; % Maghreb (Algerie/Tunisie)
    
    % GRANDE-BRETAGNE (popred) : Egypte, Soudan, Afrique Est, Afrique Sud, Nigeria, Gold Coast
    \fill[popred!40] (5.5,5.5) -- (7,5.5) -- (7,4) -- (5.5,4) -- cycle; % Egypte/Soudan
    \fill[popred!40] (6,3) -- (7.5,3.5) -- (7.5,2) -- (6,1.5) -- cycle; % Afrique Est (Kenya/Ouganda)
    \fill[popred!40] (5,1) -- (6.5,1.5) -- (6.5,0) -- (5,-0.5) -- cycle; % Afrique Sud (Cap/Rhodesies)
    \fill[popred!40] (1.5,2) -- (3,2.5) -- (3,1.5) -- (1.5,1) -- cycle; % Nigeria/Gold Coast
    
    % ALLEMAGNE (popgray) : Cameroun, Togo, SW Africa, Est Africa
    \fill[popgray!40] (2.5,2) -- (3.5,2.5) -- (3.5,1.5) -- (2.5,1) -- cycle; % Cameroun
    \fill[popgray!40] (1.2,1.8) -- (1.8,2) -- (1.8,1.5) -- (1.2,1.3) -- cycle; % Togo
    \fill[popgray!40] (5,-0.5) -- (6.5,0) -- (6.5,-1) -- (5,-1.5) -- cycle; % SW Africa
    \fill[popgray!40] (6.5,2) -- (8,2.5) -- (8,1) -- (6.5,0.5) -- cycle; % Est Africa
    
    % BELGIQUE (poppurple) : Congo
    \fill[poppurple!40] (3.5,2) -- (5.5,3) -- (5.5,0.5) -- (3.5,0) -- cycle;
    
    % PORTUGAL (popgreen) : Angola, Mozambique
    \fill[popgreen!40] (4,-0.5) -- (5.5,0) -- (5.5,-1.5) -- (4,-2) -- cycle; % Angola
    \fill[popgreen!40] (7,0.5) -- (8.5,1) -- (8.5,-0.5) -- (7,-1) -- cycle; % Mozambique
    
    % ITALIE (poporange) : Libye, Erythree, Somalie
    \fill[poporange!40] (4.5,5.5) -- (6,6) -- (6,5.5) -- (4.5,5) -- cycle; % Libye
    \fill[poporange!40] (6.5,3.5) -- (7.5,4) -- (7.5,3) -- (6.5,2.5) -- cycle; % Erythree/Somalie
    
    % ESPAGNE (poppink) : Maroc Nord, Sahara, Guinee Eq
    \fill[poppink!40] (2,6) -- (2.5,6.5) -- (3,6) -- (2.5,5.5) -- cycle; % Maroc Nord
    \fill[poppink!40] (1.5,3.5) -- (2.5,4) -- (2.5,3) -- (1.5,2.5) -- cycle; % Sahara
    \fill[poppink!40] (1.5,1.2) -- (2,1.5) -- (2,1) -- (1.5,0.7) -- cycle; % Guinee Eq
    
    % LIBERIA / ETHIOPIE (Independants) - Hachurés ou blanc avec texte
    \node[font=\bfseries\small, align=center, text=popdark] at (1.2, 1.5) {LIBÉRIA};
    \node[font=\bfseries\small, align=center, text=popdark] at (6.8, 3.5) {ÉTHIOPIE};
    
    % Légende
    \node[anchor=south west, fill=white, draw=popline, rounded corners, inner sep=4pt, font=\small, align=center] at (9.2, -1.5){
        \begin{tabular}{ll}
        \cellcolor{popblue!40}     & France (AOF, AEF, Maghreb) \\
        \cellcolor{popred!40}      & Grande-Bretagne \\
        \cellcolor{popgray!40}     & Allemagne \\
        \cellcolor{poppurple!40}   & Belgique (Congo) \\
        \cellcolor{popgreen!40}    & Portugal \\
        \cellcolor{poporange!40}   & Italie \\
        \cellcolor{poppink!40}     & Espagne \\
        \textbf{Blanc / Texte} & Libéria, Éthiopie (indép.)
        \end{tabular}
    };
    
    % Flèche Nord
    \draw[->, thick, popdark] (9.5, 6) -- (9.5, 7) node[above] {N};
\end{tikzpicture}
\legende{Carte schématique du partage de l'Afrique en 1914. Les couleurs indiquent la puissance coloniale dominante. Notez l'enchevêtrement des possessions et les deux États indépendants (Libéria, Éthiopie).}
\end{popfigure}

\begin{aretenir}[Synthèse : L'Afrique colonisée en 1914]
\begin{itemize}
    \item \textbf{Deux logiques rivales :} Axe français Est-Ouest (Dakar-Djibouti) vs Axe britannique Nord-Sud (Le Cap-Le Caire). Crise de Fachoda (1898) = apogée de la tension, résolution diplomatique au profit de la GB.
    \item \textbf{Partage quasi total :} Conférence de Berlin (1884-1885) + traités bilatéraux $\rightarrow$ frontières artificielles (lignes droites, méridiens/parallèles) ignorant les réalités ethniques.
    \item \textbf{7 puissances coloniales :} France (1\textsuperscript{re} superficie), Grande-Bretagne (1\textsuperscript{re} population/cohérence), Allemagne, Belgique, Portugal, Italie, Espagne.
    \item \textbf{Deux exceptions :} \cle{Libéria} (indépendance 1847, sphère US) et \cle{Éthiopie} (victoire d'Adoua 1896).
    \item \textbf{Conséquences structurelles :} économies d'exportation (monocultures, minerais), travail forcé/recruté, administration directe (France, Belgique, Allemagne, Portugal) vs indirecte (GB), résistance armée ou collaboration des élites locales.
\end{itemize}
\end{aretenir}

\begin{exoresolu}[Analyse d'un document : La conférence de Berlin (1885), Acte général, Article 34]
\textbf{Énoncé :} « Toute puissance qui prend possession d'un territoire sur les côtes du continent africain... en donne notification aux autres puissances signataires... afin de leur permettre, le cas échéant, de faire valoir leurs éventuelles prétentions. »
\textbf{Question :} En quoi cet article illustre-t-il la logique du « partage » de l'Afrique entre puissances européennes ?

\textbf{Solution détaillée :}
\begin{enumerate}
    \item \textbf{Identification :} Extrait de l'Acte général de la Conférence de Berlin (février 1885), conférence réunissant 14 puissances (dont France, GB, Allemagne, Portugal, Belgique) sous l'égide de Bismarck.
    \item \textbf{Contexte :} « Course au partage » (Scramble for Africa). Risque de conflits entre Européens (ex. Congo : France/Portugal/GB/Léopold II ; Niger : France/GB).
    \item \textbf{Analyse de la phrase :} 
    \begin{itemize}
        \item « Prend possession d'un territoire sur les côtes » : la possession effective (\cle{effectivité}) prime sur les droits historiques (Portugal) ou les traités de protectorat locaux. C'est la loi du plus fort sur le terrain.
        \item « En donne notification aux autres puissances » : obligation de transparence \emph{entre Européens}. Les Africains ne sont pas consultés, ni même informés. La souveraineté africaine est niée.
        \item « Faire valoir leurs éventuelles prétentions » : mécanisme de règlement des différends \emph{européens} par la diplomatie (échange de territoires, compensation) plutôt que par la guerre. Exemple : Heligoland (Allemagne) contre Zanzibar (GB) en 1890 ; Maroc (France) contre concessions au Congo (Allemagne) en 1911.
    \end{itemize}
    \item \textbf{Conclusion :} L'article 34 institutionnalise le \cle{partage pacifique} de l'Afrique par les puissances industrielles, transformant le continent en « réserve de chasse » pour les matières premières et les débouchés, sans aucune légitimité africaine. C'est l'acte de naissance juridique de la colonisation formelle.
\end{enumerate}
\end{exoresolu}

\begin{exercice}[Entraînement : Carte et comparatif]
\begin{enumerate}
    \item Sur une carte muette de l'Afrique, coloriez les possessions des 7 puissances coloniales en 1914 en respectant le code couleur du cours. Placez les capitales coloniales : Dakar, Alger, Brazzaville, Le Caire, Le Cap, Nairobi, Lagos, Windhuk, Dar es Salaam, Luanda, Lourenço Marques, Tripoli, Tétouan.
    \item Complétez le tableau comparatif suivant :
    \begin{center}
    \begin{tabularx}{\linewidth}{|l|X|X|}\hline
    \textbf{Critère} & \textbf{Empire français} & \textbf{Empire britannique} \\ \hline
    Continuité territoriale & & \\ \hline
    Mode d'administration & & \\ \hline
    Objectif stratégique majeur & & \\ \hline
    Exemple de résistance majeure & & \\ \hline
    \end{tabularx}
    \end{center}
    \item Expliquez pourquoi l'Allemagne, arrivée tardivement, a obtenu des territoires vastes mais disjoints.
\end{enumerate}
\end{exercice}

\begin{corrige}[Exercice n°1]
\begin{enumerate}
    \item \textbf{Carte :} Vérifier : France (bleu) : AOF, AEF, Maghreb, Madagascar. GB (rouge) : Égypte, Soudan, Afrique Est, Afrique Australe, Nigeria, Gold Coast. Allemagne (gris) : Cameroun, Togo, SW Africa, Est Africa. Belgique (violet) : Congo. Portugal (vert) : Angola, Mozambique, Guinée, Cap-Vert, São Tomé. Italie (orange) : Libye, Érythrée, Somalie. Espagne (rose) : Maroc Nord (capitale Tétouan), Sahara, Rio Muni. Libéria/Éthiopie : blanc/hachuré.
    \item \textbf{Tableau :}
    \begin{itemize}
        \item Continuité : France = discontinue (blocs séparés par GB/Allemagne). GB = vise continuité N-S (Cap-Caire) mais blocs disjoints (Ouest, Est, Sud, Nord).
        \item Administration : France = Directe (cercles, commandants de cercle, Code indigénat). GB = Indirecte (Indirect Rule, chefs traditionnels, Native Authorities).
        \item Objectif : France = Gloire, mission civilisatrice, matières premières, bloc continental. GB = Route des Indes (Suez/Cap), or/diamants, hégémonie maritime, axe Cap-Caire.
        \item Résistance : France = Samory Touré, Ahmadou Sekou, Behanzin. GB = Mahdi (Soudan), Ndebele/Shona (Rhodésie), Boers (Afrique du Sud), Mau Mau (plus tard, Kenya).
    \end{itemize}
    \item \textbf{Allemagne :} Arrivée après 1884 (Bismarck). Les zones côtières stratégiques déjà prises (France, GB, Portugal). Restent des zones intérieures ou secondaires (Cameroun, Togo, SW Africa, Est Africa) obtenues par traités locaux (Nachtigal, Peters) et accords inter-européens (traité Heligoland-Zanzibar 1890, accords Congo 1911). Territoires disjoints car « restes » du partage.
\end{enumerate}
\end{corrige}

\coursec{Les résistances africaines}

La section précédente a montré comment les puissances européennes, au terme de la conférence de Berlin (1884-1885), se sont lancées dans une course effrénée à l'occupation effective du continent africain. Les conquêtes françaises, britanniques, allemandes, belges, portugaises et italiennes ne se sont toutefois pas déroulées sans heurts. Loin de l'image d'une « pacification » aisée propagée par la propagande coloniale, l'imposition de la domination européenne a pris la forme d'une guerre de conquête longue, coûteuse et sanglante, face à des sociétés africaines qui ont massivement refusé de perdre leur souveraineté. Cette section analyse la diversité, l'ampleur et la portée de ces résistances.

\courssub{Typologie et diversité des résistances africaines}

Face à l'agression coloniale, les réactions africaines ne furent pas monolithiques. Les historiens distinguent généralement plusieurs formes de résistance, souvent combinées au sein d'une même société.

\begin{definition}[Résistance primaire]
On appelle \textbf{résistance primaire} (ou résistance initiale) la réaction armée immédiate des sociétés africaines face à la pénétration coloniale, dès les premiers contacts militaires ou administratifs. Elle s'oppose à la \emph{résistance secondaire}, qui survient plus tard, souvent après l'installation de l'administration coloniale (révoltes fiscales, mouvements syndicaux, nationalistes).
\end{definition}

On peut classer les résistances primaires en quatre grandes catégories, qui se chevauchent fréquemment :

\begin{itemize}
    \item \textbf{Résistances armées :} Confrontations militaires directes (batailles rangées, guérilla, siège). C'est la forme la plus visible, menée par des États structurés (empires, royaumes) ou des confédérations de chefferies.
    \item \textbf{Résistances diplomatiques :} Négociations, traités signés sous contrainte mais interprétés différemment, envois d'ambassades en Europe (ex: roi Léopold II de Belgique, reine Victoria), recours au droit international pour dénoncer la violation des accords de Berlin.
    \item \textbf{Résistances religieuses :} Mobilisation au nom de l'islam (jihad, mahdisme) ou de cultes traditionnels/syncrétiques (maji-maji en Afrique orientale allemande, kimbanguisme plus tard). La religion fournit l'idéologie unificatrice, la légitimité du chef et la motivation du sacrifice.
    \item \textbf{Résistances passives ou « par les pieds » :} Refus de payer l'impôt, refus de fournir des porteurs ou des travailleurs forcés, fuites massives vers la brousse ou vers des territoires voisins non encore conquis (ex: fuite vers le Liberia, la Gold Coast britannique ou le Cameroun français), ralentissement du travail, dissimulation des récoltes.
\end{itemize}

\begin{attention}[Piège historiographique : le mythe de l'« archaïsme »]
\emph{Erreur fréquente :} Présenter les résistances africaines comme des réactions « archaïques », « instinctives » ou « désorganisées » face à la « modernité » européenne.
\textbf{Réalité :} Beaucoup de résistances furent hautement organisées, stratégiquement réfléchies et diplomatiquement habiles. Samory Touré a modernisé son armée (fusils à répétition, ateliers de fabrication de cartouches). Ménélik II a joué les puissances européennes les unes contre les autres pour s'armer. Rudolf Duala Manga Bell a saisi la justice allemande et l'opinion publique internationale. Les qualifier d'archaïques, c'est nier l'agentivité politique et militaire des Africains.
\end{attention}

\courssub{Les grandes résistances armées de longue durée}

\textbf{2.1 Samory Touré et l'empire du Wassoulou (1882-1898) : la guérilla moderne}

Samory Touré (vers 1830-1900) est l'archétype du chef de guerre africain qui adapte ses structures à la menace européenne. Ancien marchand et chef militaire mandingue, il construit un État (le Wassoulou) s'étendant sur l'actuelle Guinée, le Mali, la Côte d'Ivoire et le Burkina Faso.

\begin{exemplebox}[Samory Touré : 16 ans de résistance]
\textbf{Contexte :} À partir de 1881, les colonnes françaises (Borgnis-Desbordes, Gallieni, Combes) progressent depuis le Sénégal et le Soudan français.
\textbf{Moyens :} Une armée permanente (\emph{sofa}) professionnalisée, pouvant atteindre \textbf{30 000 à 35 000 hommes}, équipée de fusils à répétition (Chassepot, Gras, Kropatschek) achetés sur la côte ou capturés. Création d'ateliers de réparation et de fabrication de cartouches (à Bissandougou puis en déplacement).
\textbf{Stratégie :} \textbf{La guerre de mouvement et le repli stratégique (1891-1898).} Face à la supériorité de l'artillerie et de la logistique françaises, Samory évite les batailles rangées définitives. Il applique la \emph{terre brûlée}, déplace sa capitale (Kankan $\to$ Bissandougou $\to$ Dabakala $\to$ Kong $\to$ Bouna), et tente de créer un nouvel empire à l'est (pays Kénédougou, puis pays Gouro) pour maintenir son indépendance.
\textbf{FIN :} Trahi par des alliés locaux, encerclé à Guélemou (Côte d'Ivoire actuelle), il est capturé le 29 septembre 1898 par le capitaine Gouraud. Déporté au Gabon, il y meurt en 1900.
\end{exemplebox}

\begin{exoresolu}[Analyse de l'échec de Samory]
\textbf{Énoncé :} Pourquoi Samory, malgré une armée moderne et une résistance de 16 ans, a-t-il fini par échouer ?
\textbf{Solution détaillée :}
\begin{enumerate}
    \item \textbf{Supériorité technologique et logistique française :} Artillerie (canons de 75 mm, 80 mm, 95 mm), mitrailleuses Maxim, chemin de fer (Dakar-Saint-Louis, puis Niger) et vapeur sur le Niger pour le ravitaillement. Samory n'avait pas d'artillerie lourde ni de voie ferrée.
    \item \textbf{Enclavement et blocus :} La conférence de Berlin (1885) a fermé les frontières. Samory ne peut plus s'approvisionner en armes par la côte (Liberia, Sierra Leone, Côte d'Ivoire) car les puissances européennes (France, Angleterre) coopèrent pour lui interdire l'accès aux fusils et à la poudre.
    \item \textbf{Usure humaine et matérielle :} La stratégie du repli (exode de populations entières, \emph{sofa} et civils) épuise les ressources vivrières et la cohésion sociale. Les ateliers de fabrication ne suffisent pas à compenser l'interdiction d'importation.
    \item \textbf{Divisions africaines :} Samory a dû combattre sur plusieurs fronts : Français au nord/ouest, Anglais au sud (traité de 1887 non respecté par les deux), et rois africains hostiles (Tieba de Sikasso, chefs kénédougou, gouro) qui préfèrent traiter avec les Français ou craignent l'hégémonie mandingue.
    \item \textbf{Manque d'unité nationale :} L'empire du Wassoulou est une construction militaire récente, multiethnique (Mandingues, Sénoufos, Bobos, etc.), tenue par la force et l'islam, mais fragile face à la promesse française de « paix » et de fin des razzias.
\end{enumerate}
\end{exoresolu}

\textbf{2.2 Les résistances royales : Dahomey, Ashanti, Ndébélé}

Dans les zones de forêt et de savane boisée, des royaumes centralisés très anciens opposent une résistance farouche, souvent sous forme de guerres conventionnelles.

\begin{center}
\begin{tabularx}{\linewidth}{|l|X|X|X|}
\hline
\textbf{Royaume / Chef} & \textbf{Période clé} & \textbf{Adversaire} & \textbf{Spécificités} \\
\hline
\textbf{Béhanzin (Dahomey)} & 1890-1894 & France (Dodds) & Armée permanente (\emph{Amazones} / \emph{Agoodjié}), fortifications de Cana et Abomey, diplomatie active (recherche d'armes en Europe, lettres à la reine Victoria). Défait par l'artillerie et la colonne Dodds (nov. 1892 - janv. 1894). Déporté en Martinique puis Algérie. \\
\hline
\textbf{Prempeh I (Ashanti / Asante)} & 1896-1900 (Guerre du Siège d'Or) & Angleterre & Refus du protectorat britannique (1891). Déporté aux Seychelles (1896). Résistance finale menée par la reine mère Yaa Asantewaa (1900) autour du \emph{Siège d'Or} (symbole de l'unité nationale). Annexion formelle en 1901. \\
\hline
\textbf{Lobengula (Ndébélé / Matabele)} & 1893-1894 & Compagnie britannique d'Afrique du Sud (Rhodes / Jameson) & État militaire zoulou implanté au Zimbabwe actuel. Tente la diplomatie (concession Rudd 1888), puis guerre. Défait par la \textbf{mitrailleuse Maxim} à la bataille de la rivière Shangani (1893). Disparait dans la brousse ; son fils couronné plus tard sous autorité britannique. \\
\hline
\end{tabularx}
\end{center}

\begin{aretenir}[Le choc technologique : la mitrailleuse Maxim]
L'arme décisive dans ces trois guerres (et bien d'autres) est la \textbf{mitrailleuse Maxim} (adoptée par l'armée britannique en 1889, française en 1895). Tirant 600 coups/minute, elle neutralise les charges massives d'infanterie africaine (lances, fusils à un coup). À la bataille de Shangani (1893), 50 hommes de la BSAC avec 4 Maxim mettent en déroute 5 000 guerriers ndébélés en quelques minutes. C'est le symbole de l'asymétrie technologique.
\end{aretenir}

\courssub{Les résistances à caractère religieux : le mahdisme, les confréries et le maji-maji}

L'islam et les croyances traditionnelles fournissent souvent le cadre idéologique et organisationnel pour transcender les clivages ethniques et mobiliser contre l'« infidèle » colonisateur.

\textbf{3.1 Le Mahdisme au Soudan (1881-1898) : un État théocratique anti-impérial}

Muhammad Ahmad, proclamé \textbf{Mahdi} (le « Bien guidé ») en 1881, soulève le Soudan anglo-égyptien contre la corruption turco-égyptienne et l'ingérence britannique.
\begin{itemize}
    \item \textbf{Succès initial :} Prise d'El Obeid (1883), anéantissement de la colonne Hicks (1883), siège et prise de Khartoum (janvier 1885), mort du général Gordon. Le Mahdi meurt en 1885.
    \item \textbf{Son successeur, le Khalifa Abdallahi :} Organise un État centralisé (Khartoum puis Omdurman), administre la justice islamique, lève l'impôt, tente d'industrialiser (arsenal).
    \item \textbf{Fin :} Reconquête britannique méthodique (Kitchener, chemin de fer du Nil, canonnières, artillerie, Maxim). \textbf{Bataille d'Omdurman (2 septembre 1898) :} 11 000 Britanniques/Égyptiens (artillerie, Maxim, fusils Lee-Metford) face à 50 000 \emph{ansar} (derviches). Bilan : ~10 000 morts africains, 48 morts britanniques. Fin de l'État mahdiste.
\end{itemize}

\textbf{3.2 Les confréries musulmanes (Tidjanya, Qadiriya, Sanoussya)}

En Afrique de l'Ouest et au Sahara, les confréries structurent la résistance :
\begin{itemize}
    \item \textbf{Tidjanya :} Cheikh Al-Hajj Umar Tall (Toucouleur, 1850-1864) précurseur ; plus tard, Cheikh Malick Sy, Cheikh Ahmadu Bamba Mbacké (Sénégal) : résistance spirituelle et culturelle, refus de l'école française, exil (Bamba au Gabon puis Mauritanie).
    \item \textbf{Sanoussya (Libye, Tchad, Sahara) :} Guerre de guérilla de 20 ans (1911-1931) contre l'Italie fasciste (Omar Al-Mokhtar), appuyée sur les zawiyas (centres religieux) et le réseau commercial transsaharien.
\end{itemize}

\textbf{3.3 La révolte Maji-Maji en Afrique orientale allemande (1905-1907)}

Mouvement syncrétique unifiant une vingtaine d'ethnies (Matumbi, Ngoni, Yao, etc.) autour de l'eau « sacrée » (maji) du prophète Kinjikitile Ngwale, censée rendre invulnérable aux balles allemandes.
\begin{itemize}
    \item \textbf{Déclencheur :} Travail forcé dans les plantations de coton, impôt de capitation, brutalité des \emph{akidas} (chefs indigènes auxiliaires).
    \item \textbf{Déroulement :} Attaques de postes allemands et de plantations (juillet 1905). Guérilla généralisée. Répression allemande impitoyable : terre brûlée, destruction des greniers, famine provoquée.
    \item \textbf{Bilan :} 75 000 à 300 000 morts (famine, combats). Échec militaire mais prise de conscience allemande : réforme de l'administration coloniale (Gouverneur Rechenberg, suppression du travail forcé pour le coton).
\end{itemize}

\courssub{L'exception éthiopienne : la victoire d'Adoua (1896)}

L'Éthiopie (Abyssinie) est le seul État africain à avoir conservé son indépendance par les armes au XIXe siècle (avec le Liberia, non colonisé mais sous hégémonie US).

\begin{exemplebox}[Ménélik II : modernisation et diplomatie]
\textbf{Contexte :} Traité de Wuchale (1889) avec l'Italie. Version italienne : protectorat. Version amharique : simple alliance. Ménélik dénonce le traité (1893), achète des armes (France, Russie, Italie même !), modernise l'armée (officiers russes, français), crée une artillerie, un hôpital mobile, un début de chemin de fer (Djibouti-Addis).
\textbf{Bataille d'Adoua (1er mars 1896) :} 100 000 à 120 000 Éthiopiens (fusils modernes, artillerie, cavalerie) contre 17 000 Italiens et ascari (artillerie supérieure mais mal déployée, cartes erronées). Victoire écrasante éthiopienne : ~7 000 morts italiens, 3 000 prisonniers.
\textbf{Conséquence :} Traité d'Addis-Abeba (octobre 1896) : reconnaissance de l'indépendance totale de l'Éthiopie. Choc mondial : première victoire d'une armée africaine sur une puissance européenne à l'ère coloniale. Inspiration pour le panafricanisme naissant.
\end{exemplebox}

\courssub{Les résistances au Cameroun face à la colonisation allemande (1884-1914)}

Le Cameroun, proclamé « protectorat » allemand en 1884 (traité de Douala), connaît une résistance multiforme, précoce et tardive, côtière et intérieure.

\begin{itemize}
    \item \textbf{Résistances côtières précoces (1884-1890) :} Rois Duala (Bell, Akwa, Deido) contre l'appropriation des terres et le monopole commercial. Tensions aboutissant à l'arrestation et la déportation de certains chefs ; le roi Bell (Ngando a Kwa) signe le traité mais conteste l'interprétation allemande.
    \item \textbf{Résistances de l'intérieur (pacification militaire 1890-1907) :}
    \begin{itemize}
        \item \textbf{Bamiléké (Ouest) :} Guérilla dans les montagnes, villages fortifiés (\emph{laf}), résistance de Fotso (Bandjoun), de Njoya (Bamoun) qui tente l'adaptation (écriture shümom, religion syncrétique) mais subit la déposition (1931).
        \item \textbf{Bassa (Sanaga-Maritime) :} Révolte de 1891-1892 (Ngog Lituba), puis guerre de 1900-1907 (Madola, Edinguele). Terre brûlée allemande, déportations.
        \item \textbf{Bulu / Fang (Sud) :} Résistance farouche (1895-1907), chef charismatique \textbf{Martin Paul Samba} (ex-officier allemand, fusillé 1914).
    \end{itemize}
\end{itemize}

\begin{savaistu}[Rudolf Duala Manga Bell (1873-1914) : le roi juriste]
Éduqué en Allemagne (gymnasium d'Aalen, droit à Bonn), roi des Duala (Bell) en 1908. Face au projet allemand d'exproprier les Duala de leurs terres (décret de 1910 pour créer une « ville européenne » à Douala), il mène une \textbf{résistance juridique et diplomatique} sans précédent :
\begin{enumerate}
    \item Pétitions au Reichstag (parlement allemand) et au chancelier Bethmann-Hollweg.
    \item Recours aux avocats allemands (Dr. Bell).
    \item Appel à la Société des Nations (naissante) et à l'opinion publique internationale.
    \item Alliance avec d'autres chefs (Bassa, Bamiléké) pour un front commun.
\end{enumerate}
\textbf{Issue :} Arrêté en 1914 (début de la guerre), accusé de haute trahison (contacts avec les Anglais), jugé sommairement par un tribunal militaire allemand, \textbf{pendu le 8 août 1914} à Douala avec son collaborateur Ngoso Din. Aujourd'hui, il est célébré comme \textbf{héros national} du Cameroun, symbole de la résistance par le droit et la modernité.
\end{savaistu}

\courssub{Pourquoi les résistances ont-elles (presque) toutes échoué ?}

\begin{propriete}[Causes structurelles de l'échec militaire africain (1880-1914)]
L'échec n'est pas dû à un manque de courage ou d'intelligence, mais à une \textbf{asymétrie systémique} :
\begin{enumerate}
    \item \textbf{Supériorité technologique majeure :} Fusils à répétition (Lebel 8mm, Mauser, Lee-Enfield) vs fusils à un coup (Chassepot, Snider) ou lances/arcs. \textbf{Mitrailleuse Maxim} (portable, cadence 600 c/m). \textbf{Artillerie de campagne} (75 mm français, 77 mm allemand) mobile, précise, à tir rapide. Aviation (reconnaissance, bombardement léger) dès 1911-1914 (Libye, Maroc).
    \item \textbf{Supériorité logistique et médicale :} Chemins de fer, vapeurs fluviaux, routes, télégraphe pour le commandement. Quinine (prophylaxie antipaludéenne) systématique pour les troupes européennes (troupes coloniales, Légion, tirailleurs) : taux de mortalité non-combattante divisé par 10 par rapport aux premières expéditions.
    \item \textbf{Unité de commandement et doctrine :} États-majors professionnels, doctrine de « guerre totale » (terre brûlée, prise d'otages, destruction des greniers), coordination inter-armes (infanterie, artillerie, génie, renseignement).
    \item \textbf{Divisions africaines (facteur décisif) :} Rivalités inter-ethniques, inter-dynastiques, religieuses (musulmans vs non-musulmans), sociales (esclaves vs maîtres). Les Européens exploitent systématiquement la stratégie du \emph{divide et impera} (diviser pour régner) : recrutement de \textbf{tirailleurs sénégalais, spahis, goumiers, askaris} pour combattre d'autres Africains. Ex: Tirailleurs contre Samory ; Askaris allemands contre Bamiléké ; Troupes égyptiennes/britanniques contre Mahdistes.
    \item \textbf{Blocus diplomatique et commercial :} Conférence de Berlin (acte général, art. 5-6) : libre circulation mais interdiction de fait de vendre des armes aux Africains (déclarations de Bruxelles 1890, Saint-Germain 1919). Samory et Ménélik sont les rares à percer le blocus.
\end{enumerate}
\end{propriete}

\courssub{Portée historique : la conquête fut une guerre, pas une « pacification » }

\begin{aretenir}[Bilan et postérité des résistances primaires]
\begin{itemize}
    \item \textbf{Démystification du vocabulaire colonial :} Les termes « pacification », « soumission », « tournée administrative » masquent une réalité militaire : \textbf{conquête par la guerre}. Les archives militaires françaises (SHD Vincennes) et allemandes (Bundesarchiv) comptent des centaines d'« opérations militaires », de « colonnes », d'« expéditions punitives » entre 1880 et 1914 (jusqu'en 1930 pour le Maroc, 1932 pour le Sahara mauritanien).
    \item \textbf{Coût humain considérable :} Des centaines de milliers de morts africains (combattants, civils, famines provoquées par la terre brûlée). Des milliers de morts européens (maladies, combats).
    \item \textbf{Héritage politique :} Ces résistances fondent la \textbf{mémoire nationale} des États africains indépendants (Samory en Guinée/Mali, Béhanzin au Bénin, Yaa Asantewaa au Ghana, Adoua en Éthiopie, Rudolf Duala Manga Bell au Cameroun). Elles nourrissent le nationalisme des années 1940-1960.
    \item \textbf{Continuité :} La résistance primaire se transforme en résistance secondaire (révoltes de 1914-1918, guerre du Kongo-Wara 1928-1931, maquisards UPC au Cameroun 1950-1970, Mau Mau au Kenya 1952-1960).
\end{itemize}
\end{aretenir}

\begin{methode}[Étudier un personnage de résistance : grille d'analyse en 5 étapes]
Pour traiter un sujet d'examen (commentaire de document, dissertation, exposé) sur un résistant africain, appliquez cette méthode :
\begin{enumerate}
    \item \textbf{Contexte :} Situation politique, sociale, économique de la société avant l'arrivée des Européens. Nature de l'État (centralisé, segmentaire, théocratique). Relations préexistantes avec l'extérieur (commerce, islam, traite).
    \item \textbf{Le déclencheur :} Événement précis (traité inégal, incident frontalier, dépôt de borne, exigence d'impôt/travail forcé, violation de traité).
    \item \textbf{Moyens et organisation :} Type d'armée (permanente, levée en masse, guérilla), armement (provenance, modernité), logistique (ravitaillement, ateliers), alliances (locales, extérieures), rôle de l'idéologie (islam, tradition, nationalisme naissant).
    \item \textbf{Stratégie et chronologie :} Phases de la résistance (offensive, défensive, repli, guérilla, diplomatie parallèle). Adaptation face à l'ennemi. Durée.
    \item \textbf{Bilan et postérité :} Causes de l'échec (ou succès pour Ménélik). Sort du chef (mort au combat, capture, exil, exécution, collaboration forcée). Héritage mémoriel (héros national, figure oubliée, instrumentalisation politique).
\end{enumerate}
\end{methode}

\begin{popfigure}
\begin{tikzpicture}[scale=0.85, every node/.style={font=\small}]
% Fond carte Afrique simplifiée
\draw[popmuted, line width=0.5pt] (0,0) rectangle (16,14);
% Contours très simplifiés (polygones)
\draw[popline, line width=1pt, fill=popcream] 
(1,2) -- (3,1) -- (5,1.5) -- (6,0.5) -- (7,2) -- (8,1) -- (9,2.5) -- (10,1.5) -- (11,3) -- (12,2) -- (13,4) -- (13,6) -- (11,7) -- (10,9) -- (9,8) -- (8,10) -- (6,11) -- (4,10) -- (2,9) -- (1,7) -- (0.5,5) -- (1,3) -- cycle;
% Mer
\fill[popblueL!30] (0,0) rectangle (16,14);
\draw[popline, line width=1pt, fill=popcream] 
(1,2) -- (3,1) -- (5,1.5) -- (6,0.5) -- (7,2) -- (8,1) -- (9,2.5) -- (10,1.5) -- (11,3) -- (12,2) -- (13,4) -- (13,6) -- (11,7) -- (10,9) -- (9,8) -- (8,10) -- (6,11) -- (4,10) -- (2,9) -- (1,7) -- (0.5,5) -- (1,3) -- cycle;
% Foyers de résistance : points + labels
% 1. Samory (Guinée/Mali)
\fill[popred] (3.5,3.5) circle (3pt);
\node[above right, popdark, font=\small\bfseries] at (3.5,3.5) {Samory (1882-98)};
\draw[popred, dashed, thick] (3.5,3.5) -- (2.5,4.5) -- (4,5.5) -- (5,4) -- (3.5,3.5); % Zone d'influence approx
% 2. Ashanti (Ghana)
\fill[popred] (4.5,3.0) circle (3pt);
\node[below right, popdark, font=\small\bfseries] at (4.5,3.0) {Ashanti (Prempeh/Yaa Asantewaa)};
% 3. Dahomey (Bénin)
\fill[popred] (5.2,2.8) circle (3pt);
\node[below, popdark, font=\small\bfseries] at (5.2,2.8) {Dahomey (Béhanzin)};
% 4. Mahdi (Soudan)
\fill[popred] (9.5,4.5) circle (3pt);
\node[above left, popdark, font=\small\bfseries] at (9.5,4.5) {Mahdisme (Soudan)};
% 5. Adoua (Éthiopie)
\fill[popgreen] (10.5,5.0) circle (3pt);
\node[above right, popgreen!80!black, font=\small\bfseries] at (10.5,5.0) {Adoua (1896) Victoire};
% 6. Cameroun
\fill[popred] (6.0,2.5) circle (3pt);
\node[below left, popdark, font=\small\bfseries] at (6.0,2.5) {Cameroun (Duala, Bamiléké, Bassa, Bulu)};
% 7. Sanoussya (Libye/Tchad)
\fill[popred] (8.5,3.0) circle (3pt);
\node[below, popdark, font=\small\bfseries] at (8.5,3.0) {Sanoussya (Libye/Tchad)};
% 8. Ndebele (Zimbabwe)
\fill[popred] (8.5,9.5) circle (3pt);
\node[above, popdark, font=\small\bfseries] at (8.5,9.5) {Ndébélé (Lobengula)};
% 9. Maji Maji (Tanzanie)
\fill[poporange] (10,8) circle (3pt);
\node[above, poporange!80!black, font=\small\bfseries] at (10,8) {Maji Maji (1905-07)};
% Légende
\node[anchor=south west, draw=popline, fill=popcream, rounded corners, inner sep=4pt, font=\small, align=center] at (0.2,0.2){
\textbf{Légende :}\\
$\bullet$ \textcolor{popred}{Point rouge} : Résistance armée majeure (échec militaire)\\
$\bullet$ \textcolor{popgreen}{Point vert} : Victoire militaire (Adoua)\\
$\bullet$ \textcolor{poporange}{Point orange} : Résistance religieuse/massive (Maji Maji)\\
Traitillés : Zone d'influence / déplacement (ex: Samory)
};
% Flèche Nord
\draw[->, thick, popdark] (14.5, 12.5) -- (14.5, 13.5) node[above] {N};
\end{tikzpicture}
\legende{Carte schématique des principaux foyers de résistance africaine face à la conquête coloniale (1880-1914). Les frontières actuelles sont indiquées en filigrane pour le repérage.}
\end{popfigure}

\begin{popfigure}
\begin{tikzpicture}[scale=1.2]
% Schéma "Photographie commentée" : Portrait type Ménélik II ou Samory
% Cadre photo
\draw[popdark, line width=2pt, rounded corners=5pt] (0,0) rectangle (6,8);
% Fond sépia
\fill[popgoldL!40] (0,0) rectangle (6,8);
% Silhouette stylisée (buste)
\draw[popdark, line width=1.5pt] (3,7.5) ellipse ({1.5} and {1.2}); % Tête
\draw[popdark, line width=1.5pt] (3,6.3) .. controls (2,5) and (2,3) .. (3,2) .. controls (4,3) and (4,5) .. (3,6.3); % Épaules/buste
% Détails : Couronne / Chapeau (Ménélik) ou Turban (Samory)
\draw[popgold, line width=2pt] (1.5,8) .. controls (3,9.5) .. (4.5,8); % Couronne
\draw[popgold, fill=popgold!50] (3,8.2) circle (0.15); % Joyau
% Médaille / Décoration
\draw[popblue, line width=1pt, fill=popblue!20] (2.2,4.5) rectangle (2.8,5.2);
\draw[popblue, line width=1pt, fill=popblue!20] (3.2,4.5) rectangle (3.8,5.2);
% Barbe (Ménélik)
\draw[popdark, line width=0.8pt] (2.2,6.8) .. controls (1.8,7.2) .. (2.2,7.6);
\draw[popdark, line width=0.8pt] (3.8,6.8) .. controls (4.2,7.2) .. (3.8,7.6);
\draw[popdark, line width=0.8pt] (3,6.6) -- (3,7.4);
% Légendes fléchées (callouts)
\draw[->, popblue, thick] (0.5,7.2) -- (-1.5,7.5) node[left, align=right, popdark, font=\small] {Couronne impériale\\ (symbole de légitimité\\ dynastique salomonide)};
\draw[->, popblue, thick] (0.5,5.0) -- (-1.5,5.0) node[left, align=right, popdark, font=\small] {Décorations\\ européennes/réussites\\ diplomatiques};
\draw[->, poporange, thick] (6.5,6.5) -- (7.5,6.5) node[right, align=left, popdark, font=\small] {Tenue militaire\\ modernisée\\ (officiers russes)};
\draw[->, poporange, thick] (6.5,3.5) -- (7.5,3.5) node[right, align=left, popdark, font=\small] {Regard ferme :\\(diplomate stratège,\\ unificateur)};
% Titre dans l'image
\node[align=center, popdark, font=\bfseries\small] at (3,0.5) {Ménélik II (1844-1913)\\\textit{Architecte de la victoire d'Adoua}};
\end{tikzpicture}
\legende{Portrait schématique commenté de l'empereur Ménélik II. L'analyse de l'iconographie (tenue, attributs, décorations) révèle la stratégie de modernisation sélective et de légitimation dynastique mise en œuvre pour résister à l'impérialisme italien.}
\end{popfigure}

\begin{exercice}[Entraînement au Probatoire : Commentaire de document]
\textbf{Document :} Extrait du discours de Samory Touré à ses chefs (1891), rapporté par le colonel Humbert : 
\begin{quote}
« Les Blancs veulent nous prendre notre pays. Ils ont des canons qui tonnent loin, des fusils qui tirent vite. Mais nous avons le fer, nous avons la poudre, nous avons le courage. Si nous restons unis, ils ne pourront pas nous vaincre. Mais si nous nous querellons entre nous, si les uns font la paix pendant que les autres font la guerre, alors nous serons perdus. »
\end{quote}
\textbf{Questions :}
\begin{enumerate}
    \item Situez le document (auteur, date, contexte historique précis).
    \item Identifiez les deux facteurs de force et les deux facteurs de faiblesse que Samory identifie pour son camp.
    \item Montrez, en vous appuyant sur vos connaissances, comment la stratégie de Samory (1891-1898) a tenté de pallier la faiblesse technologique.
    \item Expliquez pourquoi l'unité africaine (« si nous restons unis ») est restée un vœu pieux face à la conquête européenne.
\end{enumerate}
\end{exercice}

\begin{corrige}[Exercice n°1 : Commentaire de document sur Samory]
\textbf{1. Situation :} Auteur : Samory Touré, almamy de l'empire du Wassoulou. Date : 1891 (début de la phase de repli stratégique après les premiers revers face aux colonnes françaises de Gallieni/Combes). Contexte : Offensive française sur le Niger (prise de Ségou 1890, Nioro 1891), blocus des armes, nécessité de souder l'armée et les populations pour l'exode vers l'est.
\textbf{2. Facteurs identifiés par Samory :}
\begin{itemize}
    \item \textit{Forces :} Disponibilité du « fer » et de la « poudre » (ateliers de Bissandougou, achats résiduels), « courage » (motivation religieuse/patriotique, discipline des \emph{sofa}).
    \item \textit{Faiblesses :} Infériorité technologique implicite (« canons qui tonnent loin, fusils qui tirent vite» = artillerie et fusils à répétition/mitrailleuses), divisions internes (« si nous nous querellons », « si les uns font la paix »).
\end{itemize}
\textbf{3. Stratégie d'adaptation (1891-1898) :}
\begin{itemize}
    \item \textbf{Guerre de mouvement / guérilla :} Refus de la bataille rangée face à l'artillerie ; harcèlement des colonnes, embuscades, coupure des lignes de ravitaillement.
    \item \textbf{Repli stratégique (terre brûlée) :} Déplacement de l'État (capitale, population, ateliers) vers l'est (Kénédougou, pays Gouro) pour gagner du temps, de l'espace et tenter de reconstruire une base hors d'atteinte française.
    \item \textbf{Diplomatie de l'arme :} Tentatives d'achat d'armes via la Sierra Leone (Angleterre) et le Liberia, envoi d'émissaires en Europe.
    \item \textbf{Unification islamique :} Utilisation de l'islam comme ciment idéologique pour transcender les ethnies (Mandingues, Sénoufos, Bobos) et justifier le jihad défensif.
\end{itemize}
\textbf{4. Échec de l'unité africaine :}
\begin{itemize}
    \item \textbf{Rivalités précoloniales :} Royaumes voisins (Sikasso/Tieba, Kong) voient en Samory un conquérant mandingue plus menaçant que les Français lointains. Ils s'allient aux Français.
    \item \textbf{Stratégie européenne du « diviser pour régner » :} Traités séparés, recrutement de tirailleurs parmi les ennemis de Samory, promesses de protection aux chefs ralliés.
    \item \textbf{Hétérogénéité de l'empire du Wassoulou :} Construction récente par la force, populations non islamisées rétives à la contrainte religieuse et aux réquisitions de l'exode.
    \item \textbf{Blocus international :} Accords entre puissances (France/UK) pour interdire la vente d'armes à Samory (Conférence de Berlin, Acte de Bruxelles 1890).
\end{itemize}
\end{corrige}

\coursec{Les collaborations africaines}

Si la résistance fut la réaction majoritaire, la \textbf{collaboration} (ou « ralliement », « alliance ») fut un phénomène massif, structurant et souvent mal compris. Elle ne se réduit pas à la « trahison » : elle relève de logiques politiques, stratégiques et sociales complexes.

\courssub{Typologie des collaborations}

\begin{itemize}
    \item \textbf{Collaboration d'alliance / protectorat :} Chefs signant des traités pour se protéger d'un ennemi africain voisin ou d'une puissance européenne rivale (ex: rois du Fouta-Djalon face à Samory ; Makoko du Teke signant avec Brazza contre les rivaux du Nord).
    \item \textbf{Collaboration d'intermédiation commerciale :} Marchands et chefs côtiers servant d'intermédiaires obligés (compradores) : traite du caoutchouc, de l'ivoire, de l'huile de palme. Ils gagnent armes, prestige, monopoles.
    \item \textbf{Collaboration militaire :} Fourniture d'auxiliaires (porteurs, éclaireurs, tirailleurs). Ex: \textbf{Tirailleurs sénégalais} (créés 1857, massivement recrutés de force ou par engagement dès 1880), \textbf{Spahis} (Algérie, Maroc), \textbf{Goumiers} (Maroc), \textbf{Askaris} (Afrique orientale allemande, britannique), \textbf{Force Publique} (Congo belge).
    \item \textbf{Collaboration administrative :} Chefs de canton/cercle nommés par l'administration (chefs « warrant » au Nigeria britannique, chefs de village au Cameroun allemand/ français) pour percevoir l'impôt, recruter la main-d'œuvre, faire appliquer le code de l'indigénat.
\end{itemize}

\courssub{Motivations et logiques des collaborateurs}

\begin{aretenir}[Grille de lecture : pourquoi collaborer ?]
\begin{enumerate}
    \item \textbf{Survie physique et politique :} Face à la puissance de feu écrasante, le ralliement évite la destruction du lignage, du village, de l'État.
    \item \textbf{Rivalités inter-africaines :} L'ennemi de mon ennemi est mon allié. Samory pousse ses voisins vers les Français ; les Ndébélé poussent les Shona vers les Britanniques ; les Peuls du Macina vers les Français contre les Touareg.
    \item \textbf{Accès aux ressources et au pouvoir :} Armes, munitions, textiles, statut de « chef reconnu », monopoles commerciaux, parts de butin de guerre.
    \item \textbf{Conversion religieuse / idéologique :} Adhésion sincère à l'islam (chefs se ralliant aux Français pour contrer un jihad mahdiste ou tidjane) ou au christianisme (écoles, protection missionnaire).
    \item \textbf{Modernisation sélective :} Certains (Njoya au Cameroun, rois du Buganda en Ouganda) voient dans l'alliance européenne un levier pour moderniser leur État (écriture, administration, armée, agriculture).
\end{enumerate}
\end{aretenir}

\courssub{Figures emblématiques de la collaboration}

\begin{center}
\begin{tabularx}{\linewidth}{|l|X|X|}
\hline
\textbf{Personnage / Groupe} & \textbf{Contexte} & \textbf{Logique / Héritage} \\
\hline
\textbf{Charles Atangana (Cameroun, ~1880-1943)} & Chef Beti (Ewondo), nommé par Allemands (1901), confirmé par Français (1916). & \textbf{Intermédiation pragmatique :} Médiateur, collecteur d'impôts, recruteur. Maintient l'autonomie relative de son peuple. Figure ambiguë : collaborateur assumé, mais protecteur des siens. \\
\hline
\textbf{Apollo Kagwa (Buganda / Ouganda, 1864-1927)} & Premier ministre (Katikkiro) du Kabaka, artisan de l'accord de 1900 (Mailo Land). & \textbf{Alliance stratégique :} Buganda devient province modèle, obtient avantages (terres, administration propre, éducation). Kagwa = architecte de la « collaboration nationaliste ». \\
\hline
\textbf{Tirailleurs sénégalais / Askaris} & Recrutés (souvent de force : \emph{raflage}) dès 1880-1914. & \textbf{Contrainte et carrière :} Soldat = salaire, prestige, exemption d'impôt, mais chair à canon. Mémoire complexe : héros oubliés de Verdun (1916) ou de Bir Hakeim (1942). \\
\hline
\textbf{Chefs « warrant » (Nigeria britannique, Lugard)} & Chefs créés \emph{ex nihilo} dans les sociétés acéphales (Igbo, Tiv) pour percevoir l'impôt. & \textbf{Invention de la tradition :} Imposition d'une autorité étrangère, source de corruption et de révoltes futures (ex: guerre des femmes 1929). \\
\hline
\end{tabularx}
\end{center}

\courssub{Conséquences : une société coloniale fracturée}

\begin{itemize}
    \item \textbf{Fracture sociale :} Émergence d'une \textbf{élite collaboratrice} (chefs, commis, interprètes, premiers lettrés) vs masses paysannes soumises au code de l'indigénat, travail forcé, impôt.
    \item \textbf{Instrumentalisation de la tradition :} L'administration « invente » des chefferies là où il n'y en avait pas (sociétés segmentaires), fige des coutumes fluides, crée des « coutumiers » à sa solde.
    \item \textbf{Héritage postcolonial :} Ces clivages (Nord/Sud, côte/intérieur, ethnies « favorisées » / « marginalisées ») structurent encore les politiques nationales (ex: Cameroun, Nigeria, Rwanda, Côte d'Ivoire).
\end{itemize}

\coursec{Bilan et synthèse : l'Afrique colonisée en 1914}

\courssub{Le partage achevé : la carte de 1914}

En 1914, à la veille de la Première Guerre mondiale, \textbf{90 \% du continent} est sous domination européenne formelle. Seuls l'\textbf{Éthiopie} et le \textbf{Liberia} conservent une indépendance nominale (l'Éthiopie par la victoire militaire, le Liberia par la tutelle américaine).

\begin{popfigure}
\begin{tikzpicture}[scale=0.8, every node/.style={font=\tiny}]
% Carte Afrique 1914 simplifiée : zones colorées par puissance
% Contour Afrique
\draw[popline, line width=1pt, fill=popcream] 
(1,2) -- (3,1) -- (5,1.5) -- (6,0.5) -- (7,2) -- (8,1) -- (9,2.5) -- (10,1.5) -- (11,3) -- (12,2) -- (13,4) -- (13,6) -- (11,7) -- (10,9) -- (9,8) -- (8,10) -- (6,11) -- (4,10) -- (2,9) -- (1,7) -- (0.5,5) -- (1,3) -- cycle;
% Zones coloniales (très approximatives)
% France (popblueL) : AOF, AEF, Maghreb, Madagascar
\fill[popblueL!50] (1,2) -- (3,1) -- (5,1.5) -- (6,2) -- (6,4) -- (5,5) -- (4,5) -- (3,4) -- (2,5) -- (1,4) -- cycle; % AOF approx
\fill[popblueL!50] (4,5) -- (6,4) -- (7,5) -- (7,7) -- (5,8) -- (4,7) -- cycle; % AEF approx
\fill[popblueL!50] (7,2) -- (9,2) -- (9,4) -- (7,4) -- cycle; % Tunisie/Algérie approx
\fill[popblueL!50] (10,7) -- (11,7) -- (11,9) -- (10,9) -- cycle; % Madagascar
% UK (popredL) : Égypte, Soudan, Afrique de l'Est, Afrique du Sud, Nigeria, Gold Coast, Sierra Leone, Gambie
\fill[popredL!50] (9,3) -- (11,3) -- (11,5) -- (9,5) -- cycle; % Égypte/Soudan
\fill[popredL!50] (10,5) -- (12,5) -- (12,8) -- (10,8) -- cycle; % Afrique Est
\fill[popredL!50] (8,9) -- (11,9) -- (11,11) -- (8,11) -- cycle; % Afrique du Sud
\fill[popredL!50] (5,3) -- (6,3) -- (6,4) -- (5,4) -- cycle; % Nigeria
\fill[popredL!50] (4,2.5) -- (5,2.5) -- (5,3) -- (4,3) -- cycle; % Gold Coast
% Allemagne (popgrayL) : Cameroun, Togo, SWA, AOE
\fill[popgrayL!50] (5.5,2.5) -- (6.5,2.5) -- (6.5,3.5) -- (5.5,3.5) -- cycle; % Cameroun
\fill[popgrayL!50] (4.2,2.5) -- (4.8,2.5) -- (4.8,3) -- (4.2,3) -- cycle; % Togo
\fill[popgrayL!50] (8,8) -- (9,8) -- (9,10) -- (8,10) -- cycle; % SWA
\fill[popgrayL!50] (10,6) -- (11.5,6) -- (11.5,8) -- (10,8) -- cycle; % AOE
% Belgique (popgold!18) : Congo
\fill[popgold!18] (7,4) -- (10,4) -- (10,7) -- (7,7) -- cycle;
% Portugal (poppurpleL) : Angola, Mozambique, Guinée
\fill[poppurpleL!50] (7,7) -- (9,7) -- (9,10) -- (7,10) -- cycle; % Angola approx
\fill[poppurpleL!50] (10,8) -- (12,8) -- (12,11) -- (10,11) -- cycle; % Mozambique
% Italie (poporangeL) : Libye, Érythrée, Somalie
\fill[poporangeL!50] (8,3) -- (10,3) -- (10,5) -- (8,5) -- cycle; % Libye
\fill[poporangeL!50] (10.5,4.5) -- (11.5,4.5) -- (11.5,5.5) -- (10.5,5.5) -- cycle; % Corne de l'Afrique
% Espagne (poppinkL) : Maroc espagnol, Rio de Oro, Guinée espagnole
\fill[poppinkL!50] (7,2) -- (7.5,2) -- (7.5,3) -- (7,3) -- cycle;
% Éthiopie (popgreenL) - Indépendant
\fill[popgreenL!50] (10,4.5) -- (11.5,4.5) -- (11.5,6.5) -- (10,6.5) -- cycle;
\node[popgreen!80!black, font=\bfseries\tiny] at (10.7,5.5) {ÉTHIOPIE};
% Liberia (popgreenL) - Indépendant
\fill[popgreenL!50] (3.5,2) -- (4.2,2) -- (4.2,2.8) -- (3.5,2.8) -- cycle;
\node[popgreen!80!black, font=\bfseries\tiny] at (3.8,2.4) {LIBERIA};
% Légende
\node[anchor=south west, draw=popline, fill=popcream, rounded corners, inner sep=3pt, font=\tiny, align=center] at (0.2,0.2){
\textbf{Légende 1914 :}\\
\textcolor{popblueL}{\rule{3mm}{2mm}} France (AOF, AEF, Maghreb, Madagascar)\\
\textcolor{popredL}{\rule{3mm}{2mm}} Royaume-Uni (Égypte, Soudan, Est, Sud, Nigeria, Gold Coast)\\
\textcolor{popgrayL}{\rule{3mm}{2mm}} Allemagne (Cameroun, Togo, SWA, AOE)\\
\textcolor{popgold!18}{\rule{3mm}{2mm}} Belgique (Congo)\\
\textcolor{poppurpleL}{\rule{3mm}{2mm}} Portugal (Angola, Mozambique, Guinée)\\
\textcolor{poporangeL}{\rule{3mm}{2mm}} Italie (Libye, Corne de l'Afrique)\\
\textcolor{poppinkL}{\rule{3mm}{2mm}} Espagne (Maroc, Rio de Oro)\\
\textcolor{popgreenL}{\rule{3mm}{2mm}} Indépendants (Éthiopie, Liberia)
};
% Flèche Nord
\draw[->, thick, popdark] (14.5, 12.5) -- (14.5, 13.5) node[above] {N};
\end{tikzpicture}
\legende{Carte schématique du partage de l'Afrique en 1914. Les zones sont approximatives et ne rendent pas compte des enclaves, condominiums (Soudan anglo-égyptien) ni des frontières exactes.}
\end{popfigure}

\courssub{Tableau récapitulatif du partage colonial (1914)}

\begin{center}
\begin{tabularx}{\linewidth}{|l|X|c|X|}
\hline
\textbf{Puissance} & \textbf{Principaux territoires} & \textbf{Superficie (km²) env.} & \textbf{Population (1914) env.} \\
\hline
France & AOF, AEF, Algérie, Tunisie, Maroc (prot. 1912), Madagascar, Djibouti, Comores & $\sim$ 10,5 millions & $\sim$ 30 millions \\
\hline
Royaume-Uni & Égypte (prot. 1882), Soudan (condominium), Nigeria, Gold Coast, Sierra Leone, Gambie, Kenya, Ouganda, Nyassaland, Rhodésie du Nord/Sud, Afrique du Sud (dominion), Bechuanaland, Basutoland, Swaziland, Zanzibar (prot.) & $\sim$ 9 millions & $\sim$ 40 millions \\
\hline
Allemagne & Cameroun, Togo, Afrique du Sud-Ouest, Afrique orientale allemande & $\sim$ 2,5 millions & $\sim$ 13 millions \\
\hline
Belgique & Congo belge (État indépendant du Congo jusqu'en 1908) & $\sim$ 2,3 millions & $\sim$ 10-15 millions \\
\hline
Portugal & Angola, Mozambique, Guinée portugaise, São Tomé-et-Principe, Cap-Vert & $\sim$ 2 millions & $\sim$ 8 millions \\
\hline
Italie & Libye (conquise 1911-12), Érythrée, Somalie italienne & $\sim$ 1,9 million & $\sim$ 2 millions \\
\hline
Espagne & Maroc espagnol (prot. 1912), Rio de Oro (Sahara occidental), Guinée espagnole & $\sim$ 0,3 million & $\sim$ 0,3 million \\
\hline
\end{tabularx}
\end{center}

\courssub{Les trois piliers du système colonial naissant}

\begin{propriete}[L'architecture de la domination (1914)]
\begin{enumerate}
    \item \textbf{L'administration directe / indirecte :} 
    \begin{itemize}
        \item \emph{Directe (France, Portugal, Belgique, Allemagne) :} Cadres européens à tous les échelons, code de l'indigénat, assimilation théorique (France), exploitation concessionnaire (Belgique, Allemagne).
        \item \emph{Indirecte (Royaume-Uni, Lugard) :} « \emph{Indirect Rule} » : pouvoir confié aux chefs traditionnels (ou créés) sous contrôle du \emph{District Officer}. Moins de cadres, coûts moindres, mais fige les sociétés.
    \end{itemize}
    \item \textbf{L'économie d'extraction :} 
    \begin{itemize}
        \item Chemins de fer « du port à la mine/plantation » (Dakar-Niger, Abidjan-Niger, Matadi-Léopoldville, Benguela, Tanzanie).
        \item Cultures de rente imposées (coton, cacao, café, arachide, huile de palme, caoutchouc, sisal).
        \item Travail forcé (corvée, \emph{prestation}, \emph{chibalo} au Mozambique, \emph{ibuko} au Congo belge) et recrutement migrant (mine d'or du Witwatersrand, cuivre du Katanga).
    \end{itemize}
    \item \textbf La « mission civilisatrice » : 
    \begin{itemize}
        \item École coloniale (très sélective : fils de chefs, interprètes), santé (lutte contre la trypanosomiase, variole, lèpre), « mise en valeur » (barrages, ports).
        \item Mais : ségrégation raciale (quartiers européens / indigènes), justice d'exception, natalité contrôlée, démographie stagnante ou en baisse dans zones de travail forcé (Congo, AEF).
    \end{itemize}
\end{enumerate}
\end{propriete}

\courssub{L'Afrique en 1914 : une poudrière pour la Grande Guerre}

\begin{aretenir}[L'Afrique dans la Première Guerre mondiale (1914-1918)]
\begin{itemize}
    \item \textbf{Théâtre d'opérations majeur :} Campagnes du Togo, Cameroun, Afrique du Sud-Ouest, Afrique orientale allemande (von Lettow-Vorbeck). 2 millions d'Africains mobilisés (combattants + porteurs). 200 000 à 250 000 morts africains (maladies, épuisement, combats).
    \item \textbf{Révélateur des faiblesses coloniales :} Défaite allemande = redistribution des mandats (SDN) : Cameroun (France/UK), Togo (France/UK), SWA (Afrique du Sud), AOE (Tanganyika UK, Ruanda-Urundi Belgique).
    \item \textbf{Accélérateur du nationalisme :} Retour des tirailleurs (promesses non tenues), découverte de la métropole affaiblie, diffusion des idées wilsoniennes (droit des peuples), révoltes de 1916-1918 (Volta-Bani, Dahomey, Libya, Égypte 1919).
\end{itemize}
\end{aretenir}

\begin{methode}[Réussir une question de synthèse / dissertation sur le partage de l'Afrique]
\textbf{Sujet type :} « L'Afrique face à l'impérialisme européen (1880-1914) : conquêtes, résistances, collaborations. »
\begin{enumerate}
    \item \textbf{Analyse du sujet :} Bornes chronologiques (1880 Conférence de Berlin / 1914 Guerre), concepts clés (impérialisme, conquête, résistance, collaboration), espace (Afrique entière).
    \item \textbf{Problématique :} Comment le partage de l'Afrique résulte-t-il de l'interaction entre la violence coloniale, les stratégies africaines (résistance/collaboration) et les rivalités européennes ?
    \item \textbf{Plan en 3 parties (équilibrées) :}
    \begin{enumerate}
        \item \textbf{La conquête : une violence structurée par les rivalités européennes.} (Conférence de Berlin, traités, occupation effective, moyens techniques, piliers admin/éco).
        \item \textbf{Les réponses africaines : un spectre complet de la résistance à la collaboration.} (Typologie, grands exemples : Samory, Béhanzin, Mahdi, Adoua, Maji-Maji, Cameroun ; Collaboration : types, motivations, figures, conséquences).
        \item \textbf{1914 : un continent transformé, une domination fragile.} (Carte 1914, bilan démographique/éco, fractures sociales, entrée en guerre = test de vérité).
    \end{enumerate}
    \item \textbf{Rédaction :} Intro (accroche, définitions, pb, plan). Développement (arguments + exemples précis + chiffres). Conclusion (bilan + ouverture sur 1918-1945).
\end{enumerate}
\end{methode}

\begin{exercice}[Entraînement au Probatoire : Dissertation]
\textbf{Sujet :} « Les sociétés africaines n'ont pas été de simples victimes passives de la colonisation. » Discutez cette affirmation en vous appuyant sur l'étude des résistances et des collaborations (1880-1914).
\textbf{Attendus :} Nuancer « victimes passives » : agency (agentivité) africaine. Montrer la diversité des stratégies (résistance armée, diplomatique, religieuse, passive ; collaboration d'alliance, militaire, administrative). Montrer que la collaboration elle-même est une stratégie active de survie/pouvoir. Conclure sur l'héritage complexe (héroïsation des résistants, stigmatisation des collaborateurs, fractures postcoloniales).
\end{exercice}

\begin{corrige}[Exercice n°2 : Dissertation - Éléments de correction]
\textbf{Introduction :} Accroche (Césaire « chose colonisée » vs agency). Définitions. Problématique : En quoi les Africains ont-ils été acteurs de leur histoire face à l'impérialisme ? Plan en 3 parties.
\textbf{I. La résistance : affirmation de la souveraineté et de l'agentivité.}
\begin{itemize}
    \item A. Diversité des formes (armée, diplomatique, religieuse, passive) et des acteurs (États, chefferies, confréries, masses).
    \item B. Exemples de stratégies élaborées : Samory (guérilla moderne, ateliers), Ménélik (diplomatie, modernisation, victoire), Mahdi (État théocratique), Maji-Maji (unification interethnique).
    \item C. Échec militaire structurel (asymétrie techno/logistique, divisions, blocus) mais succès mémoriel (fondements des nationalismes).
\end{itemize}
\textbf{II. La collaboration : une stratégie active d'adaptation et de pouvoir.}
\begin{itemize}
    \item A. Pas une soumission passive : choix rationnel face au rapport de forces (survie, rivalités, accès ressources).
    \item B. Figures de l'intermédiation : chefs de canton (Atangana), élites lettrées (Duala Manga Bell - résistant par le droit mais issu de l'élite collaboratrice), rois du Buganda (Kagwa).
    \item C. Coût social : fracture élites/masses, invention de la tradition, instrumentalisation ethnique.
\end{itemize}
\textbf{III. 1914 : un continent africain acteur de la reconfiguration mondiale.}
\begin{itemize}
    \item A. Participation massive à la Grande Guerre (tirailleurs, porteurs, ressources) -> prise de conscience politique.
    \item B. Redéploiement colonial (mandats SDN) négocié sans Africains mais sur leurs terres.
    \item C. Émergence de nouvelles formes d'agentivité : syndicalisme, presse, associations, panafricanisme (Congrès 1919).
\end{itemize}
\textbf{Conclusion :} Bilan : Africains = acteurs stratégiques (résistance/collaboration) dans des contraintes structurelles. Héritage : mémoires plurielles, États-nations nés du partage, défis de la réconciliation mémorielle.
\end{corrige}

\coursec{Les collaborations africaines}

Après avoir étudié les multiples formes de résistance opposées par les sociétés africaines à la conquête européenne, il convient d'examiner un phénomène tout aussi massif mais souvent moins visible : la collaboration. Loin d'être un bloc monolithique, l'Afrique a vu émerger des stratégies d'alliance, d'accommodation ou d'appui actif aux puissances coloniales. Comprendre ces logiques, sans anachronisme ni jugement moral hâtif, est indispensable pour saisir la complexité de la conquête et la formation des États coloniaux.

\courssub{Définition et formes de la collaboration}

\begin{definition}[La collaboration]
La \cle{collaboration} désigne l'ensemble des attitudes et des actions par lesquelles des individus, des groupes ou des entités politiques africains apportent un appui — actif ou passif, ponctuel ou durable — à la conquête, à l'installation et au fonctionnement de l'administration coloniale européenne. Elle ne se confond pas avec une adhésion idéologique au projet colonial ; elle relève le plus souvent d'un calcul pragmatique face à un rapport de force militaire et technologique écrasant.
\end{definition}

La collaboration prend des formes variées, qui s'imbriquent souvent :

\begin{itemize}
    \item \textbf{La signature de traités de protectorat ou d'amitié :} des chefs acceptent, parfois sous la contrainte ou par méconnaissance des clauses écrites, de placer leur territoire sous la « protection » d'une puissance européenne (traité germano-douala de 1884, traités français avec le roi Makoko du Teke en 1880).
    \item \textbf{L'alliance militaire et la fourniture d'auxiliaires :} des groupes africains combattent aux côtés des colonnes européennes contre d'autres Africains (peuples alliés des Français contre Samory Touré, auxiliaires britanniques dans les guerres de l'Ashanti).
    \item \textbf{Le recrutement dans les corps militaires coloniaux :} les \cle{tirailleurs sénégalais}, créés en 1857, constituent l'exemple majeur de cette main-d'œuvre guerrière africaine au service de la France.
    \item \textbf{L'appui logistique et renseigné :} porteurs, guides, interprètes, informateurs sont indispensables à la progression des colonnes en terrain inconnu.
    \item \textbf{L'intégration dans l'administration indigène :} chefs de canton, gardes-cercles, commis, interprètes officiels assurent le relais local du pouvoir colonial.
    \item \textbf{L'adhésion culturelle et religieuse :} conversion au christianisme, scolarisation dans les écoles des missions, adoption de la langue et des codes du colonisateur, ouvrant l'accès aux emplois subalternes de l'administration.
\end{itemize}

\begin{popfigure}
\begin{tikzpicture}[scale=0.85, every node/.style={font=\small}]
% Tableau comparatif Résistances / Collaborations
\draw[popline, line width=1pt] (0,0) rectangle (16,10.5);
% En-têtes
\fill[popblue!20] (0,9.5) rectangle (16,10.5);
\node[align=center, font=\bfseries] at (1.75,10) {Critère};
\node[align=center, font=\bfseries] at (6,10) {Résistances};
\node[align=center, font=\bfseries] at (12.25,10) {Collaborations};
% Lignes de séparation horizontales
\foreach \y in {0,1,2,3,4,5,6,7,8,9.5} \draw[popline] (0,\y) -- (16,\y);
% Lignes verticales
\draw[popline] (3.5,0) -- (3.5,10.5);
\draw[popline] (8.5,0) -- (8.5,10.5);

% Contenu
% Ligne 1 : Formes
\node[align=left] at (1.75,9) {\textbf{Formes}};
\node[align=left] at (6,8.7) {Guerres prolongées, guérilla,\\ refus de traiter, soulèvements,\\ jihad, diplomatie anti-européenne};
\node[align=left] at (12.25,8.7) {Signatures de traités, enrôlement\\ (tirailleurs), auxiliaires militaires,\\ guides/porteurs, administration\\ indigène, conversion/scolarisation};

% Ligne 2 : Acteurs principaux
\node[align=left] at (1.75,7.5) {\textbf{Acteurs}};
\node[align=left] at (6,7.2) {Chefs traditionnels (Samory,\\ Béhanzin, Rabah), chefs religieux,\\ populations entières, élites lettrées\\ (ex. Abd el-Kader)};
\node[align=left] at (12.25,7.2) {Chefs rivaux (Makoko, rois Douala),\\ élites commerçantes, convertis,\\ lettrés missionnés, jeunes\\ en quête de promotion sociale};

% Ligne 3 : Motivations
\node[align=left] at (1.75,6) {\textbf{Motivations}};
\node[align=left] at (6,5.7) {Défense de la souveraineté,\\ de l'islam, de l'honneur,\\ refus de l'impôt/travail forcé,\\ survie culturelle};
\node[align=left] at (12.25,5.7) {Rivalités interethnatiques,\\ élimination d'un ennemi local,\\ intérêts commerciaux, calcul\\ politique, survie physique,\\ promotion individuelle};

% Ligne 4 : Exemples majeurs
\node[align=left] at (1.75,4.5) {\textbf{Exemples}};
\node[align=left] at (6,4.2) {Samory Touré (Mandinka),\\ Béhanzin (Dahomey), Rabah\\ (Bornou), Menelik II (Éthiopie),\\ Maji Maji (Afrique orientale)};
\node[align=left] at (12.25,4.2) {Makoko (Teke), rois Bell/Akwa\\ (Douala), tirailleurs sénégalais,\\ auxiliaires de Voulet-Chanoine,\\ chefs de canton en AOF};

% Ligne 5 : Issus / Devenir
\node[align=left] at (1.75,3) {\textbf{Issus}};
\node[align=left] at (6,2.7) {Échec militaire fréquent,\\ déportation, mort au combat,\\ mais mémoire héroïque,\\ fondement des nationalismes};
\node[align=left] at (12.25,2.7) {Intégration dans l'appareil\\ colonial, privilèges relatifs,\\ mais perte d'autonomie réelle,\\ stigmatisation post-indépendance};

% Ligne 6 : Frontières poreuses
\node[align=left] at (1.75,1.5) {\textbf{Porosité}};
\node[align=left] at (6,1.2) {Certains résistants signent\\ un traité puis le rompent\\ (ex. Samory, Béhanzin)};
\node[align=left] at (12.25,1.2) {Certains collaborateurs\\ basculent dans la résistance\\ (ex. chefs de canton,\\ élites lettrées nationalistes)};
\end{tikzpicture}
\legende{Tableau comparatif : résistances et collaborations africaines face à la conquête européenne (Première technique — Histoire).}
\end{popfigure}

\courssub{Les motivations profondes : rivalités, intérêts et calculs politiques}

La décision de collaborer ne relève jamais d'un unique mobile. L'historien doit croiser les facteurs structurels et conjoncturels :

\begin{enumerate}
    \item \textbf{Les rivalités interethnétiques et interétatiques :} c'est le moteur principal. Un chef menacé par un voisin puissant (Samory, Rabah, l'empire Ashanti, le royaume du Dahomey) voit dans l'Européen un allié de circonstance pour renverser la menace. L'ennemi de mon ennemi devient mon ami, fût-il étranger et lointain.
    \item \textbf{Les intérêts commerciaux :} les élites côtières (Douala, Lagos, Freetown, Saint-Louis) tirent profit du commerce légitime (huile de palme, arachide, caoutchouc) et craignent que la résistance n'entrave les flux. Elles négocient leur statut d'intermédiaires indispensables.
    \item \textbf{Le calcul politique des chefs traditionnels :} face à la supériorité militaire évidente (fusils à répétition, artillerie, marine de guerre), certains souverains préfèrent « traiter » pour préserver une autonomie interne, leur dynastie et la vie de leurs sujets. C'est la logique du \emph{moindre mal}.
    \item \textbf{La quête de promotion sociale individuelle :} pour des jeunes sans héritage, l'enrôlement comme tirailleur, porteur ou commis offre un salaire, un statut, une formation technique, une mobilité géographique impossible dans la société traditionnelle.
    \item \textbf{La conversion religieuse et l'accès à l'instruction :} les missions catholiques et protestantes offrent éducation, soins, protection juridique. Les « catéchistes » et premiers lettrés forment une nouvelle élite médiatrice, souvent en tension avec les chefs traditionnels.
\end{enumerate}

\begin{attention}[Le piège de l'anachronisme moral]
\emph{Erreur fréquente :} qualifier les collaborateurs de « traîtres » ou de « vendus » avec les catégories patriotiques d'aujourd'hui (nation, indépendance, panafricanisme).
\emph{Réalité historique :} au XIXe siècle, l'identité nationale n'existe pas ; les loyautés sont lignagères, ethniques, religieuses, villageoises. Un chef bassa qui aide les Allemands contre un chef bamileke ne « trahit pas le Cameroun » (qui n'existe pas encore), il défend son clan contre un rival séculaire. L'historien doit \textbf{contextualiser} : replacer l'acteur dans son horizon d'attente, ses contraintes immédiates, ses sources d'information limitées.
\end{attention}

\courssub{Exemples majeurs : alliances, traités et tirailleurs}

\begin{exemplebox}[Le roi Makoko et le traité de 1880 (Congo français)]
Le roi \textbf{Makoko} (Iloo Ier), souverain des Téké sur le fleuve Congo, signe le 10 septembre 1880 un traité avec \textbf{Pierre Savorgnan de Brazza}. Menacé par les raids des négriers bobangis venus du nord et par l'expansion du royaume voisin du Kongo, Makoko cherche un allié puissant. Il cède la rive droite du Congo à la France en échange de protection militaire et de biens d'échange (tissus, armes, alcool). Ce traité permet à la France de fonder Brazzaville et de verrouiller le bassin du Congo face à Léopold II de Belgique. Makoko conserve son autorité coutumière jusqu'à sa mort (1887), mais son successeur perdra toute marge de manœuvre.
\end{exemplebox}

\begin{exemplebox}[Les rois Bell et Akwa : le traité germano-douala du 12 juillet 1884]
Le \textbf{12 juillet 1884}, les rois \textbf{Bell (Ndumbé Lobé)} et \textbf{Akwa (Dika Mpondo)}, souverains des deux principales chefferies douala sur l'estuaire du Wouri, signent un traité de protectorat avec le consul allemand \textbf{Nachtigal}. Motivations : garantir leur monopole commercial sur l'intérieur (huile de palme, ivoire), obtenir des armes contre les raids de l'intérieur (Bassa, Bakoko), et jouer les Allemands contre les Britanniques (déjà installés à Victoria/Ambas Bay). Le traité stipule le respect des droits coutumiers et du commerce libre. \textbf{Réalité :} l'Allemagne n'en fera qu'à sa tête, imposant bientôt l'administration directe, l'impôt et le travail forcé. Les rois deviennent des « chefs de canton » sous tutelle.
\end{exemplebox}

\begin{savaistu}[Le traité fondateur du Kamerun]
Ce traité du 12 juillet 1884 est l'acte de naissance juridique du Kamerun allemand. Nachtigal, commissaire impérial, utilise une canonnière (\emph{Möwe}) pour imposer sa présence. Les rois Douala, marchands avertis, pensent signer un traité commercial et d'alliance entre égaux ; le texte allemand (\emph{Schutzvertrag}) instaure une suzeraineté juridique. Cette asymétrie informationnelle est la marque de fabrique de la conférence de Berlin (1884-1885) qui suit de peu.
\end{savaistu}

\begin{definition}[Les tirailleurs sénégalais]
Corps de soldats africains recrutés par la France dès \textbf{1857} (décret de Louis Faidherbe, gouverneur du Sénégal), d'abord parmi les anciens captifs affranchis (« tirailleurs » = tireurs d'élite en tiraillement), puis par conscription (loi Mangin 1912). Ils constituent la \textbf{chair à canon} et le fer de lance des conquêtes françaises en Afrique occidentale et équatoriale (guerres contre Samory, Rabah, conquête du Maroc), puis dans les deux guerres mondiales (Verdun, Chemin des Dames, Italie, Provence). Leur statut est ambigu : soldats français à part entière sur le papier, soumis au code de l'indigénat dans les faits, payés moins que les métropolitains, cantonnés dans des bataillons séparés.
\end{definition}

\begin{exemplebox}[Les auxiliaires de la mission Voulet–Chanoine (1898-1899)]
La colonne Voulet–Chanoine, chargée de rallier le Tchad, recrute massivement des \textbf{auxiliaires africains} : tirailleurs réguliers, mais aussi « goumiers » irréguliers, porteurs, femmes de troupe, recrutés au fil de l'avancée (Sénégal, Soudan, Dahomey, Niger). Ces auxiliaires commettent les pires exactions (pillages, viols, massacres de villages) sous les ordres d'officiers français eux-mêmes hors de contrôle. L'épisode illustre la \textbf{violence intrinsèque de la conquête} et le rôle des Africains comme acteurs (et non simples victimes) de cette violence, enrôlés de force ou par appât du butin.
\end{exemplebox}

\courssub{Collaboration et résistance : des stratégies changeantes}

La dichotomie « résistants / collaborateurs » est un leurre commode. La réalité est faite de \textbf{glissements, de réversibilités, de doubles jeux}.

\begin{methode}[Nuancer un jugement historique : la méthode du croisement]
\begin{enumerate}
    \item \textbf{Identifier la source :} qui parle ? (rapport colonial, mémoire orale, correspondance de chef, presse nationaliste ultérieure). Quelle intention ?
    \item \textbf{Distinguer les temporalités :} l'attitude en 1890 (conquête) n'est pas celle de 1920 (administration) ni de 1950 (nationalisme).
    \item \textbf{Repérer les rapports de force locaux :} l'acteur a-t-il le choix ? Subit-il une contrainte militaire, une famine, une épidémie ?
    \item \textbf{Croiser les points de vue :} le même chef peut être « loyal » pour l'administrateur, « traître » pour le chef rival, « sage » pour sa propre lignée.
    \item \textbf{Éviter l'essentialisme :} « les Douala ont collaboré » est faux ; certains clans Douala ont résisté (Rudolf Duala Manga Bell, exécuté en 1914), d'autres ont négocié.
\end{enumerate}
\end{methode}

\begin{exemplebox}[Le cas camerounais : Rudolf Duala Manga Bell, du collaborationnisme à la résistance légale]
Fils du roi Bell, éduqué en Allemagne, \textbf{Rudolf Duala Manga Bell} (1873-1914) incarne l'élite « évoluée » collaboratrice : il parle allemand, porte l'uniforme, siège au conseil de district. Mais face à l'expropriation programmée des terres douala (projet de zone européenne séparée), il bascule. Il utilise les \textbf{armes du droit colonial} : pétitions au Reichstag, recours juridiques, presse allemande (alliés socialistes et libéraux). Il tente une alliance transethnique avec les Bassa et les Bakoko. Arrêté pour « haute trahison » (correspondance avec les Britanniques en 1914), il est pendu le 8 août 1914 avec son allié Ngoso Din. \textbf{Leçon :} la collaboration d'hier (éducation, administration) fournit les outils juridiques et langagiers de la résistance de demain.
\end{exemplebox}

\begin{exemplebox}[Les chefs de canton en AOF : de l'administration indirecte au nationalisme]
Le système de l'\textbf{administration indirecte} (Lugard, puis généralisé en AOF) transforme les chefs traditionnels en \textbf{chefs de canton} : percepteurs d'impôts, recruteurs de main-d'œuvre forcée (prestations), auxiliaires de police. Beaucoup acceptent par contrainte ou intérêt. Mais dans l'entre-deux-guerres, certains (ex. chefs baoulé en Côte d'Ivoire, chefs peuls au Soudan) utilisent leur position pour protéger leurs sujets, freiner les excès, voire soutenir secrètement les premiers syndicats (RDA). À l'indépendance, beaucoup sont écartés par les nouveaux partis uniques, accusés de « féodalisme » et de « collaboration ».
\end{exemplebox}

\courssub{Bilan nuancé : la collaboration comme fait structurel de la conquête}

\begin{aretenir}[Bilan : la collaboration, moteur et limite de l'ordre colonial]
\begin{itemize}
    \item \textbf{Facilitatrice indispensable :} sans guides, porteurs, interprètes, tirailleurs, chefs signataires, la conquête aurait été plus lente, plus coûteuse, peut-être impossible avec les faibles effectifs européens (ratio 1/1000 en AOF en 1914).
    \item \textbf{Logiques locales, non impériales :} les acteurs africains poursuivent \emph{leurs} intérêts (sécurité, commerce, pouvoir, survie), non le projet colonial. L'Européen est une \emph{ressource} parmi d'autres.
    \item \textbf{Fragilité de l'alliance :} le colonisateur trahit systématiquement les clauses protectrices des traités (Makoko, Douala, chefs du Fouta-Djalal). La collaboration n'achète qu'un sursis.
    \item \textbf{Héritage contradictoire :} les élites issues de la collaboration (lettrés, militaires, chefs de canton) forment le vivier des premiers nationalistes (Houphouët-Boigny, Senghor, Ahidjo, Nkrumah étaient fils de chefs ou lettrés missionnés). La \cle{décolonisation} se fait souvent par les « fils de collaborateurs ».
\end{itemize}
\end{aretenir}

\begin{popfigure}
\begin{tikzpicture}[scale=0.75, every node/.style={font=\small}]
% Carte schématique Afrique de l'Ouest : alliances locales lors de la conquête
% Contours très simplifiés
\draw[popline, line width=1.5pt] (0,0) .. controls (2,-1) and (6,-1) .. (8,0) .. controls (10,1) and (10,4) .. (9,6) .. controls (7,8) and (4,8) .. (2,7) .. controls (0,6) and (0,3) .. (0,0);
% Fleuve Niger
\draw[popblue, line width=2pt] (2,1) .. controls (3,2) and (4,3) .. (5,4) .. controls (5.5,5) and (6,5.5) .. (6.5,6);
% Fleuve Sénégal
\draw[popblue, line width=1.5pt] (1,1) .. controls (1.5,2) .. (1.5,3);
% Villes / points clés
\node[circle, fill=popred, inner sep=2pt] (samory) at (4.5,2.5) {};
\node[above] at (samory) {Samory (Wassoulou)};
\node[circle, fill=popgreen, inner sep=2pt] (bamako) at (3.5,3.5) {};
\node[left] at (bamako) {Bamako (Fr)};
\node[circle, fill=poporange, inner sep=2pt] (segou) at (5,3.8) {};
\node[right] at (segou) {Ségou (Toucouleur)};
\node[circle, fill=poppurple, inner sep=2pt] (kankan) at (5.5,2) {};
\node[below] at (kankan) {Kankan (Samory)};
\node[circle, fill=popblue, inner sep=2pt] (kita) at (3,3) {};
\node[left] at (kita) {Kita (Fr)};
\node[circle, fill=poppink, inner sep=2pt] (sikasso) at (5.2,2.8) {};
\node[above right] at (sikasso) {Sikasso (Tata)};
% Zones d'alliances : hachures
\pattern[pattern=north west lines, pattern color=popgreen!40] (2,2) .. controls (3,2.5) and (3.5,3) .. (4,3) .. controls (4.5,2.5) and (4,2) .. (3.5,2) .. controls (3,1.5) .. (2,2);
\node[align=center, font=\tiny, text=popgreen!80!black] at (3,2.3) {Alliés français\\ (Khassonké,\\ Bambara du Kaarta)};
\pattern[pattern=north east lines, pattern color=popred!40] (5,2.5) .. controls (5.5,2.8) and (6,3) .. (6.5,3) .. controls (6,2.5) and (5.5,2) .. (5,2.5);
\node[align=center, font=\tiny, text=popred!80!black] at (5.7,2.7) {Alliés de\\ Samory (Dyula,\\ certains Malinké)};
% Côte : Saint-Louis, Gorée, Dakar
\node[circle, fill=popblue, inner sep=2pt] (stl) at (1.2,2.5) {};
\node[left] at (stl) {Saint-Louis};
\node[circle, fill=popblue, inner sep=2pt] (dakar) at (1.5,1.8) {};
\node[below] at (dakar) {Dakar};
% Flèches de progression française
\draw[popgreen, ->, thick, decorate, decoration={snake, amplitude=1pt, segment length=4pt}] (2.5,3.5) -- (4,3.5);
\draw[popgreen, ->, thick, decorate, decoration={snake, amplitude=1pt, segment length=4pt}] (3,3) -- (4.5,2.8);
% Légende
\begin{scope}[xshift=10cm, yshift=2cm]
\node[align=left, font=\small] at (0,0) {\textbf{Légende}};
\draw[popline] (0,-0.5) rectangle (4,-4.2);
\node[circle, fill=popred, inner sep=2pt] at (0.5,-0.8) {}; \node[right] at (0.7,-0.8) {État de Samory (résistance)};
\node[circle, fill=popgreen, inner sep=2pt] at (0.5,-1.3) {}; \node[right] at (0.7,-1.3) {Postes / axes français};
\node[circle, fill=poporange, inner sep=2pt] at (0.5,-1.8) {}; \node[right] at (0.7,-1.8) {Empire toucouleur (Ahmadou)};
\node[circle, fill=poppurple, inner sep=2pt] at (0.5,-2.3) {}; \node[right] at (0.7,-2.3) {Bases de Samory};
\node[circle, fill=poppink, inner sep=2pt] at (0.5,-2.8) {}; \node[right] at (0.7,-2.8) {Tata de Sikasso (résistance)};
\draw[pattern=north west lines, pattern color=popgreen!40] (0.3,-3.1) rectangle (0.7,-3.3); \node[right] at (0.9,-3.2) {Populations alliées aux Français};
\draw[pattern=north east lines, pattern color=popred!40] (0.3,-3.5) rectangle (0.7,-3.7); \node[right] at (0.9,-3.6) {Populations alliées à Samory};
\draw[popblue, line width=2pt] (0.3,-3.9) -- (0.7,-3.9); \node[right] at (0.9,-3.9) {Fleuve Niger};
\draw[popgreen, ->, thick, decorate, decoration={snake, amplitude=1pt, segment length=4pt}] (0.3,-4.2) -- (0.7,-4.2); \node[right] at (0.9,-4.2) {Progression française};
\end{scope}
% Nord
\node[rotate=0] at (14,7) {\textbf{N}};
\draw[->, thick] (14,6.5) -- (14,7.5);
\end{tikzpicture}
\legende{Carte schématique : alliances locales et fronts de la conquête française au Soudan occidental (1880-1898). Les hachures vertes indiquent les populations (Khassonké, Kaarta) alliées aux Français contre Samory ; les hachures rouges, les groupes ralliés à Samory.}
\end{popfigure}

\begin{exoresolu}[Analyse d'un document : le traité germano-douala de 1884]
\textbf{Document (extrait) :} « Sa Majesté le Roi Bell et Sa Majesté le Roi Akwa, rois du pays des Douala, déclarent placer leurs pays sous la protection de l'Empire allemand... L'administration allemande s'engage à respecter les droits de propriété et les coutumes des indigènes... Les rois conservent leur autorité judiciaire et administrative selon les coutumes locales. »
\textbf{Question :} En quoi ce traité illustre-t-il les ambiguïtés de la collaboration africaine ?
\textbf{Solution détaillée :}
\begin{enumerate}
    \item \textbf{Identification :} Traité de protectorat (protection inégale), signé le 12 juillet 1884 par deux rois rivaux (Bell et Akwa) unis par l'intérêt commun face à l'intérieur et aux Britanniques.
    \item \textbf{Analyse des clauses :} reconnaissance mutuelle (les rois « déclarent placer » — forme volontaire), promesse de respect des coutumes et de l'autorité judiciaire (clause clé pour les chefs), silence sur la souveraineté foncière (piège : l'Allemagne s'arroge le domaine éminent).
    \item \textbf{Contextualisation :} Les Douala sont des intermédiaires commerciaux puissants (huile de palme). Ils craignent l'expansion britannique (Victoria) et les raids de l'intérieur. L'Allemagne (Nachtigal) arrive avec des canons (canonnière \emph{Möwe}) et des traités types.
    \item \textbf{Nuance :} Ce n'est pas une « vente » du pays. Les rois pensent signer un \emph{traité d'alliance entre souverains}. La traduction allemande (Schutzvertrag) implique une subordination juridique que la version orale douala ne dit pas. \textbf{Asymétrie informationnelle} majeure.
    \item \textbf{Devenir :} Dès 1885, l'administration allemande impose des taxes, restreint le commerce des rois, nomme des « chefs de canton » rivaux. La résistance légale de Rudolf Duala Manga Bell (1914) naît de cette trahison du texte signé.
    \item \textbf{Conclusion :} La collaboration ici est un \textbf{choix rationnel sous contrainte}, basé sur une incompréhension juridique délibérément entretenue. Elle permet l'installation allemande mais contient en germe la contestation future.
\end{enumerate}
\end{exoresolu}

\begin{exercice}[Pour approfondir]
\begin{enumerate}
    \item À l'aide du tableau comparatif (figure 1), rédigez un paragraphe argumenté de 10 lignes montrant que la frontière entre résistance et collaboration est poreuse. Appuyez-vous sur deux exemples précis (un d'Afrique de l'Ouest, un du Cameroun).
    \item Un candidat à l'oral du Probatoire affirme : « Les tirailleurs sénégalais étaient des mercenaires qui ont trahi l'Afrique. » En vous appuyant sur la définition encadrée et la méthode « Nuancer un jugement historique », construisez une réponse structurée en trois points pour corriger cette affirmation.
    \item Sur la carte schématique (figure 2), identifiez deux zones où les populations ont collaboré avec les Français contre Samory Touré. Expliquez, pour chacune, la motivation locale probable (rivalité ethnique, religieuse, économique).
\end{enumerate}
\end{exercice}

\begin{corrige}[Exercice n°1 -- Éléments de réponse]
\textbf{1. Porosité résistance/collaboration :} Le chef bassa qui guide la colonne allemande contre un chef bamileke rival (collaboration tactique) peut, cinq ans plus tard, refuser de fournir la quote-part de porteurs exigée (résistance passive). Au Cameroun, les rois Bell et Akwa signent le traité de 1884 (collaboration diplomatique) ; leur héritier Rudolf Duala Manga Bell mène une résistance juridique jusqu'à la potence (1914). En Afrique de l'Ouest, les Khassonké du Kaarta s'allient aux Français pour détruire l'empire de Samory (collaboration militaire, 1890) ; mais dès 1915, leurs descendants refusent la conscription pour la Grande Guerre (résistance à l'ordre colonial). La collaboration n'est pas une identité figée, c'est une \textbf{stratégie situationnelle}.
\textbf{2. Tirailleurs : pas des mercenaires, des conscrits.} (a) Recrutement : d'abord volontaires (affranchis), puis \textbf{conscription légale} (loi Mangin 1912) — obligation citoyenne dans l'empire français, pas contrat mercenaire. (b) Statut : soldats français, soumis au code de justice militaire, décorés (Légion d'honneur), pensionnés — non traités en auxiliaires jetables. (c) Motivation : contrainte légale, mais aussi promotion sociale (solde, grade, formation), devoir perçu vis-à-vis de la « mère patrie » (idéologie républicaine enseignée à l'école). Juger « trahison » suppose une conscience nationale panafricaine inexistante en 1914.
\textbf{3. Zones de collaboration anti-Samory :} (a) \textbf{Kaarta (région de Nioro/Kita)} : royaume bambara rival de Samory (conflit ancien pour le contrôle des routes du sel et de l'or). Alliés français dès 1890. (b) \textbf{Khassonké (région de Bafoulabé/Kayes)} : peuple mandingue non islamisé, en butte aux razzias de Samory (jihad imposé). Fournissent guides, porteurs, auxiliaires. Dans les deux cas : \textbf{survie face à un prédateur local} prime sur l'opposition à l'étranger européen.
\end{corrige}

\coursec{Bilan et synthèse : l'Afrique colonisée en 1914}

Après avoir analysé les résistances armées et les stratégies de collaboration des sociétés africaines face à l'invasion européenne, nous devons mesurer l'ampleur du bouleversement. En 1914, à la veille de la Première Guerre mondiale, le « partage de l'Afrique » est consommé. Le continent a cessé d'être une mosaïque d'États et de sociétés souveraines pour devenir un échiquier colonial. Cette section dresse le bilan territorial, politique, économique et social de cette conquête, avec une attention particulière au cas camerounais, avant d'ouvrir sur les dynamiques qui mèneront aux nationalismes.

\courssub{Un continent partagé : le bilan territorial de 1914}

\emph{Intuition :} Imaginez une carte de l'Afrique où chaque couleur représente un maître européen. En 1880, l'essentiel de l'intérieur du continent échappe encore aux Européens ; il est occupé par des États africains puissants (califat de Sokoto, empire de Samory Touré, royaume zoulou, royaume Merina à Madagascar, empire d'Éthiopie). Trente-quatre ans plus tard, la carte a radicalement changé : deux seules taches échappent aux couleurs européennes.

\begin{popfigure}
\begin{tikzpicture}[scale=0.75]
% Contour très simplifié de l'Afrique (polygone approximatif)
\draw[popdark, thick, fill=popcream]
(0,0) -- (1,-0.5) -- (3,-1) -- (5,-1.5) -- (7,-2) -- (9,-3) -- (10,-5) -- (9.5,-7) -- (8,-9) -- (6,-10) -- (4,-10.5) -- (2,-10) -- (0.5,-9) -- (-0.5,-7) -- (-1,-5) -- (-1,-3) -- (0,-1.5) -- cycle;

% France (AOF, AEF, Maghreb) : bleu
\fill[popblue!40] (-1,-1.5) -- (0,0) -- (3,-1) -- (4,-3) -- (3,-4) -- (1,-4) -- (-0.5,-3) -- cycle;
\fill[popblue!40] (5,-1.5) -- (7,-2) -- (6,-4) -- (4,-3) -- cycle;
% Grande-Bretagne : encre (Nigeria, Afrique de l'Est, Afrique australe, Égypte-Soudan)
\fill[popink!40] (2,-4) -- (3,-4) -- (3,-6) -- (1.5,-6) -- cycle;
\fill[popink!40] (5,-4) -- (7,-3) -- (8,-6) -- (6,-6) -- cycle;
\fill[popink!40] (4,-8) -- (7,-8) -- (6,-10) -- (3,-10) -- cycle;
\fill[popink!40] (-0.5,-5) -- (1,-6) -- (0,-9) -- (-1,-8) -- cycle;
% Allemagne : gris-bleu (Cameroun, Togo, Afrique de l'Est, Sud-Ouest)
\fill[popmuted!70] (3,-4) -- (4,-3) -- (5,-4) -- (4,-5.5) -- cycle;
\fill[popmuted!70] (1.5,-4.5) -- (2.5,-4) -- (3,-5) -- (2,-5.5) -- cycle;
\fill[popmuted!70] (5.5,-5) -- (7,-4) -- (7.5,-7) -- (6,-7.5) -- cycle;
\fill[popmuted!70] (4,-8.5) -- (5.5,-8) -- (5,-9.5) -- (3.5,-9.5) -- cycle;
% Belgique (Congo) : orange
\fill[poporange!50] (3,-5) -- (5,-4) -- (6,-6) -- (5,-8) -- (3,-7) -- cycle;
% Portugal (Angola, Mozambique) : vert
\fill[popgreen!40] (3,-7) -- (4.5,-6.5) -- (5,-8.5) -- (3.5,-9) -- cycle;
\fill[popgreen!40] (6.5,-7) -- (8,-6.5) -- (8.5,-9) -- (7,-9.5) -- cycle;
% Italie (Libye, Corne) : violet
\fill[poppurple!40] (-0.5,-2) -- (1,-1.5) -- (2,-3) -- (0.5,-3.5) -- cycle;
\fill[poppurple!40] (6.5,-4.5) -- (7.5,-4) -- (8,-6) -- (7,-6.5) -- cycle;
% Espagne (Rif, Rio de Oro, Guinée) : rose
\fill[poppink!50] (-0.5,-1) -- (0.5,-0.5) -- (1.5,-1.5) -- (0.5,-2) -- cycle;
\fill[poppink!50] (0.5,-4.5) -- (1.5,-4) -- (2,-5) -- (1,-5.5) -- cycle;

% États indépendants : or
\fill[popgold!60] (5.5,-4.5) -- (6.5,-4) -- (7,-5.5) -- (6,-6) -- cycle;
\node[popdark, font=\bfseries\small] at (6.2,-5) {Éthiopie};
\fill[popgold!60] (1,-4.5) -- (1.8,-4) -- (2,-5) -- (1.2,-5.5) -- cycle;
\node[popdark, font=\bfseries\small] at (1.5,-4.7) {Libéria};

% Mise en évidence du Cameroun
\draw[poporange, very thick] (3,-4) -- (4,-3) -- (5,-4) -- (4,-5.5) -- cycle;
\node[poporange, font=\bfseries\small, rotate=20] at (4.6,-3.4) {\textbullet\ Cameroun};

% Villes repères
\foreach \x/\y/\name in {1/-2/Dakar, 3.5/-5.7/Yaoundé, 6/-5.6/Nairobi, 5/-9.3/Le Cap, -0.5/-3.2/Alger, 6.5/-4.3/Addis-Abeba}
  \node[popdark, font=\tiny] at (\x,\y) {\name};

% Flèche du Nord
\node[popdark] at (-2,1) {\Large $\uparrow$ \textbf{N}};

% Légende manuelle
\begin{scope}[yshift=-12cm, xshift=-1cm]
\node[anchor=west, font=\small] at (0,0) {\colorbox{popblue!40}{\phantom{aa}} France};
\node[anchor=west, font=\small] at (3,0) {\colorbox{popink!40}{\phantom{aa}} Grande-Bretagne};
\node[anchor=west, font=\small] at (6.5,0) {\colorbox{popmuted!70}{\phantom{aa}} Allemagne};
\node[anchor=west, font=\small] at (9.5,0) {\colorbox{poporange!50}{\phantom{aa}} Belgique};
\node[anchor=west, font=\small] at (0,-0.7) {\colorbox{popgreen!40}{\phantom{aa}} Portugal};
\node[anchor=west, font=\small] at (3,-0.7) {\colorbox{poppurple!40}{\phantom{aa}} Italie};
\node[anchor=west, font=\small] at (6.5,-0.7) {\colorbox{poppink!50}{\phantom{aa}} Espagne};
\node[anchor=west, font=\small, text=popdark] at (9.5,-0.7) {\colorbox{popgold!60}{\phantom{aa}} Indépendants};
\end{scope}
\end{tikzpicture}
\legende{Carte schématique de l'Afrique en 1914 (croquis simplifié, sans précision cartographique). Les couleurs indiquent la puissance coloniale. Le Cameroun (entouré d'orange) est alors une colonie allemande. Seuls l'Éthiopie et le Libéria conservent leur souveraineté.}
\end{popfigure}

\begin{aretenir}[Bilan territorial de 1914]
À la veille de la Grande Guerre, environ \textbf{90 \% du territoire africain} est sous domination européenne. Seuls deux États africains conservent une indépendance reconnue :
\begin{itemize}
\item L'\textbf{Éthiopie} (Abyssinie), qui a vaincu l'Italie à la bataille d'\textbf{Adoua} (1\ier{} mars 1896) et modernisé son armée ;
\item Le \textbf{Libéria}, fondé au XIX\ieme{} siècle par d'anciens esclaves affranchis venus des États-Unis, et resté sous influence américaine.
\end{itemize}
Les puissances coloniales principales sont la \textbf{France} (empire le plus étendu : AOF, AEF, Maghreb, Madagascar), la \textbf{Grande-Bretagne} (empire le plus peuplé et le plus continu, du Caire au Cap : Égypte, Soudan, Nigeria, Afrique de l'Est, Afrique australe), l'\textbf{Allemagne} (quatre colonies : Togo, Cameroun, Afrique orientale allemande, Sud-Ouest africain), la \textbf{Belgique} (Congo), le \textbf{Portugal} (Angola, Mozambique, Guinée-Bissau), l'\textbf{Italie} (Libye, Érythrée, Somalie) et l'\textbf{Espagne} (Rif, Rio de Oro, Guinée équatoriale).
\end{aretenir}

\begin{popfigure}
\begin{tikzpicture}
\begin{axis}[
    ybar,
    bar width=12pt,
    width=\linewidth,
    height=8cm,
    enlargelimits=0.15,
    ylabel={Superficie (millions de km$^2$)},
    symbolic x coords={France, Grande-Bretagne, Allemagne, Belgique, Portugal, Italie, Espagne},
    xtick=data,
    nodes near coords,
    nodes near coords align={vertical},
    ymin=0,
    ymax=12,
    ymajorgrids=true,
    grid style=dashed,
    title style={font=\bfseries},
    title={Superficies des empires coloniaux en Afrique (vers 1914)},
    x tick label style={rotate=30, anchor=east, font=\small},
]
\addplot[fill=popblue!60, draw=popdark] coordinates {(France, 10.7) (Grande-Bretagne, 9.3) (Allemagne, 2.6) (Belgique, 2.3) (Portugal, 2.0) (Italie, 1.9) (Espagne, 0.3)};
\end{axis}
\end{tikzpicture}
\legende{Diagramme en barres des superficies coloniales en Afrique vers 1914 (valeurs approchées, en millions de km$^2$). La France domine par l'étendue (le Sahara représente une large part de son empire), la Grande-Bretagne par la population et la continuité territoriale. Source : estimations historiques usuelles.}
\end{popfigure}

\courssub{Conséquences politiques : des frontières artificielles}

\emph{Intuition :} Avant 1885, les limites entre États africains sont des « marches » floues, des zones d'influence entre royaumes, souvent définies par des fleuves, des chaînes de montagnes ou des lignes de partage des eaux. La conférence de Berlin (1884-1885) impose le principe de l'\emph{occupation effective} et le tracé de frontières géométriques sur la carte, en Europe, sans connaissance du terrain ni consultation des populations.

\begin{definition}[Frontières artificielles]
Les frontières coloniales sont des lignes tracées \emph{à la règle} (méridiens, parallèles) ou suivant des cours d'eau, lors de conférences diplomatiques européennes (Berlin, 1884-1885, puis accords bilatéraux franco-britanniques, franco-allemands, etc.). Elles ignorent les réalités ethniques, linguistiques, religieuses et historiques des populations africaines.
\end{definition}

\begin{exemplebox}[Des peuples divisés par les frontières]
Le tracé des frontières a éclaté de nombreux groupes humains entre deux, trois, voire quatre administrations coloniales distinctes.
\begin{itemize}
\item Les peuples \textbf{Béti-Pahuin} (dont les Ewondo) : partagés entre le \textbf{Cameroun}, le \textbf{Gabon} (AEF française) et la \textbf{Guinée équatoriale} espagnole ;
\item Les \textbf{Somali} : divisés entre la \textbf{Somalie italienne}, le \textbf{Somaliland britannique}, l'\textbf{Éthiopie} (Ogaden), la \textbf{Côte française des Somalis} (future Djibouti) et le \textbf{Kenya} britannique ;
\item Les \textbf{Ewe} : partagés entre le \textbf{Togo} allemand (puis français et britannique) et la \textbf{Côte-de-l'Or} britannique (futur Ghana) ;
\item Les \textbf{Haoussa} : partagés entre le \textbf{Nigeria} britannique et le \textbf{Niger} français ;
\item Les \textbf{Bakongo} : partagés entre le \textbf{Congo belge}, le \textbf{Congo français} et l'\textbf{Angola} portugais.
\end{itemize}
Cette division a des conséquences durables : difficultés de construction nationale après les indépendances, revendications frontalières, problèmes de gestion transfrontalière des ressources.
\end{exemplebox}

\begin{savaistu}[Des frontières à la règle]
Une part importante des frontières africaines actuelles est constituée de lignes droites (méridiens ou parallèles) tracées sur des cartes en Europe, sans repère géographique naturel sur le terrain : c'est le cas, par exemple, d'une grande partie des frontières de la Libye, de l'Égypte, du Soudan ou de la Namibie. Ce phénomène est sans équivalent en Europe, où les frontières suivent le plus souvent des éléments naturels ou des héritages historiques.
\end{savaistu}

\emph{Raisonnement :} Cette géométrie coloniale a détruit la légitimité des États précoloniaux (royaumes du Dahomey, ashanti, du Wadaï, du Buganda, califat de Sokoto, empire de Samory Touré). Les administrateurs coloniaux ont soit démantelé ces structures (arrestation et déportation des souverains : Béhanzin au Dahomey, Samory Touré, Prempeh I\ier{} chez les Ashanti), soit les ont instrumentalisées au moyen de l'\emph{indirect rule} (voir plus loin), créant une fracture durable entre l'État moderne hérité du colonisateur et les sociétés traditionnelles.

\courssub{Conséquences économiques : l'économie de traite et l'exploitation}

\emph{Intuition :} La colonisation n'est pas seulement politique ; elle restructure en profondeur les économies africaines pour les insérer de force dans le capitalisme mondial, comme fournisseurs de matières premières et comme marchés captifs pour les produits européens.

\begin{propriete}[Caractéristiques de l'économie coloniale d'exploitation]
\begin{enumerate}
\item \textbf{Spécialisation forcée (monoculture ou mono-extraction) :} chaque territoire est assigné à une production d'exportation : arachide (Sénégal), cacao et café (Côte d'Ivoire, Cameroun, Gold Coast), coton (Soudan français, Ouganda), cuivre (Katanga, Rhodésie du Nord), or (Afrique du Sud, Gold Coast), caoutchouc (Congo, Cameroun).
\item \textbf{Travail forcé et contrainte :} réquisition de main-d'œuvre pour les travaux publics (chemins de fer, ports), les plantations concédées aux compagnies et les mines. Exemples : le chantier du chemin de fer Congo-Océan (1921-1934), qui coûta la vie à plusieurs milliers de travailleurs forcés (estimations usuelles : 15 000 à 20 000 morts) ; le \emph{chibalo} au Mozambique ; le \textbf{code de l'indigénat} en AOF et en AEF, qui légalise les réquisitions et les prestations.
\item \textbf{Infrastructures « en éventail » :} les chemins de fer et les ports relient l'intérieur producteur à la côte exportatrice, mais ne relient pas les territoires entre eux (pas de liaison Est-Ouest). Exemples : lignes Dakar-Niger, Abidjan-Ouagadougou, Douala-Yaoundé puis Ngaoundéré.
\item \textbf{Démantèlement de l'artisanat local :} l'importation massive de produits manufacturés européens (textiles, alcool, quincaillerie) ruine les productions locales (tissage, forge, poterie).
\item \textbf{Monétarisation et fiscalité :} l'impôt de capitation, payable en monnaie, oblige les paysans à vendre leur force de travail ou leurs récoltes pour payer l'impôt, au détriment des cultures vivrières.
\end{enumerate}
\end{propriete}

\begin{attention}[Erreur fréquente : traite négrière et économie coloniale]
Ne pas confondre l'\emph{économie de la traite négrière atlantique} (XVI\ieme{}-XIX\ieme{} siècles) et l'\emph{économie coloniale} (fin XIX\ieme{} siècle-1960). La traite négrière exportait des \textbf{êtres humains}. L'économie coloniale exporte des \textbf{matières premières} (produits tropicaux, minerais) en exploitant une main-d'œuvre africaine contrainte \emph{sur place}. C'est une mutation structurelle majeure qu'il faut savoir expliquer au Probatoire.
\end{attention}

\courssub{Conséquences sociales et culturelles : bouleversements et naissance d'élites}

\emph{Intuition :} La conquête s'accompagne d'une révolution sociale et culturelle. Les missions chrétiennes (catholiques et protestantes) deviennent des auxiliaires majeurs de la colonisation « pacifique » : éducation, santé, encadrement social.

\begin{itemize}
\item \textbf{Bilan démographique contrasté :} chute brutale de la population dans les zones de conquête violente, de travail forcé intensif (Congo, Cameroun, Afrique équatoriale) et d'épidémies (la trypanosomiase, ou maladie du sommeil, est favorisée par les déplacements forcés de populations). Reprise démographique à partir des années 1920-1930 avec les campagnes de vaccination, la fin des guerres de conquête et la « paix coloniale ».
\item \textbf{Christianisation massive :} les chrétiens sont très minoritaires en Afrique subsaharienne vers 1900 ; ils représentent une part importante des populations vers 1960. Les missions construisent les premiers hôpitaux, dispensaires et écoles.
\item \textbf{Scolarisation limitée et sélective :} l'école coloniale forme des auxiliaires (commis, interprètes, infirmiers, catéchistes, instituteurs). Les taux de scolarisation restent très faibles (moins de 10 \% des enfants en AOF en 1945) ; le secondaire et le supérieur sont quasi inexistants (la première université d'AOF, à Dakar, n'est créée qu'en 1950, puis institut officiel en 1957). Les programmes sont calqués sur ceux de la métropole, déconnectés des réalités locales.
\item \textbf{Naissance d'une élite occidentalisée (les \emph{évolués}) :} petits fonctionnaires, enseignants, pasteurs, avocats, médecins formés en Afrique ou en métropole (Dakar, Paris, Londres). Ils maîtrisent les codes du colonisateur (langue, droit, écriture) et deviendront les artisans des nationalismes : Houphouët-Boigny, Senghor, Nkrumah, Nyerere, et au Cameroun Ruben Um Nyobè et Félix-Roland Moumié.
\end{itemize}

\courssub{Le cas camerounais : du Kamerun allemand au partage franco-britannique}

Le Cameroun illustre de manière exemplaire l'arbitraire du partage et la superposition de deux modèles coloniaux.

\begin{enumerate}
\item \textbf{Le Kamerun allemand (1884-1916) :} traité de protectorat signé avec les chefs Duala (12 juillet 1884). Administration centrale forte (gouverneur), « pacification » brutale (campagnes militaires contre les résistances de l'intérieur), économie de plantation (cacao, banane, caoutchouc, huile de palme) reposant sur le travail forcé, premières infrastructures (chemins de fer du Nord et du Centre, port de Douala).
\item \textbf{La conquête alliée (1914-1916) :} pendant la Première Guerre mondiale, double invasion franco-britannique. Chute de Yaoundé (janvier 1916) ; le dernier bastion allemand de Mora capitule en février 1916.
\item \textbf{Le partage (1916-1922) :} la ligne de partage \textbf{Milner-Simon} (1916) est confirmée par le traité de Versailles (1919) puis par les mandats de la Société des Nations (1922).
\begin{itemize}
\item \textbf{Cameroun français} (environ les 4/5 du territoire) : administration directe, langue française, droit français, mise en valeur par les plantations (cacao, café, banane), travail forcé (prestations), code de l'indigénat.
\item \textbf{Cameroun britannique} (deux zones : \emph{Northern} et \emph{Southern Cameroons}) : administré depuis le Nigeria selon l'\emph{indirect rule}, langue anglaise, droit coutumier et \emph{common law}, investissements limités, économie tournée vers le commerce (palmistes, cacao).
\end{itemize}
\item \textbf{Héritage : le dualisme structurel.} Deux systèmes éducatifs (francophone et anglophone), deux traditions juridiques (droit civil français et \emph{common law}), deux langues officielles, deux cultures administratives. Cette fracture, que la réunification de 1961 n'a pas effacée, est au cœur des tensions anglophones contemporaines.
\end{enumerate}

\begin{attention}[Piège à éviter : le changement de maître en 1916]
\emph{Erreur fréquente :} présenter l'histoire coloniale du Cameroun comme une continuité française depuis 1884.
\emph{Réalité :} le Cameroun a changé de puissance coloniale en \textbf{1916}. La période allemande (1884-1916) a structuré le territoire (frontières, axes ferroviaires, introduction du cacao, administration centralisée). La période du double mandat (1916-1960/1961) a superposé deux logiques antagonistes. Comprendre le Cameroun contemporain (bilinguisme, bijuridisme, question anglophone) impose de maîtriser cette \textbf{double tutelle}.
\end{attention}

\courssub{Modèles coloniaux : exploitation ou peuplement ; administration directe ou indirecte}

Pour le Probatoire, il est indispensable de distinguer deux logiques coloniales qui structurent différemment les sociétés africaines.

\begin{center}
\begin{tabularx}{\linewidth}{|l|X|X|}
\hline
\rowcolor{popblueL} \textbf{Critère} & \textbf{Colonie d'exploitation (AOF, AEF, Cameroun français, Congo belge)} & \textbf{Colonie de peuplement (Algérie, Kenya, Rhodésie du Sud, Afrique du Sud)} \\
\hline
Objectif principal & Extraction de richesses (matières premières) & Installation durable de populations européennes (colons) \\
\hline
Population européenne & Faible (administrateurs, militaires, missionnaires, commerçants) & Importante (fermiers, mineurs, fonctionnaires) \\
\hline
Accès à la terre & Terres déclarées « vacantes et sans maître » concédées aux compagnies ; paysannerie maintenue sous contrainte & \textbf{Spoliation massive} des meilleures terres (hauts plateaux du Kenya, terres algériennes) au profit des colons \\
\hline
Main-d'œuvre & Contrainte sur place (impôt, réquisitions, travail forcé) & Contrainte et \textbf{migrations forcées} (réserves, \emph{pass laws}) \\
\hline
Droits politiques & Quasi nuls (statut de sujet, code de l'indigénat) & Nuls pour les Africains ; droits civiques réservés aux colons \\
\hline
Exemples & Majeure partie de l'Afrique de l'Ouest, centrale et orientale & Algérie, Afrique du Sud, Rhodésie du Sud, Kenya, Angola, Mozambique \\
\hline
\end{tabularx}
\end{center}

\begin{definition}[Indirect rule (administration indirecte) — modèle britannique]
Théorisée par Lord \textbf{Lugard} (Nigeria), l'\emph{indirect rule} est un système d'administration coloniale dans lequel la puissance européenne gouverne \emph{à travers} les autorités traditionnelles (chefs, émirs, rois) reconnues ou créées par elle.
\begin{itemize}
\item \textbf{Mécanisme :} le chef traditionnel perçoit l'impôt, recrute la main-d'œuvre, applique le droit coutumier codifié, juge dans les tribunaux indigènes ; le \emph{District Officer} britannique supervise.
\item \textbf{Opposition avec l'administration directe (modèle français) :} la France nomme ses propres agents (commandants de cercle) et des chefs de canton souvent \emph{créés} par l'administration ; elle impose le droit et la langue français. Le chef de canton français est un \emph{auxiliaire de l'administration}, non un chef traditionnel légitime.
\item \textbf{Conséquences :} l'\emph{indirect rule} a figé les chefferies et renforcé le particularisme ethnique, tandis que l'administration directe a favorisé une élite unifiée par la langue et le droit français, mais déracinée des sociétés locales.
\end{itemize}
\end{definition}

\courssub{Ouverture : de la conquête à la mise en valeur, puis aux nationalismes}

Le bilan de 1914 n'est pas une fin, mais le début d'une nouvelle phase : la \textbf{mise en valeur} (années 1920-1930). Les puissances coloniales, affaiblies par la Grande Guerre, cherchent à rentabiliser leurs empires (plans d'équipement, grands travaux, exposition coloniale de 1931). C'est dans ce cadre que naissent, dans les années 1930-1940, les premiers mouvements nationalistes organisés, portés par les élites issues de l'école coloniale, qui retourneront les principes de la « mission civilisatrice » (droits de l'homme, citoyenneté, droit des peuples à disposer d'eux-mêmes) contre le système colonial. Les leçons suivantes analyseront cette transition.

\courssub{Méthode : rédiger une conclusion de dissertation d'histoire}

\begin{methode}[Structure type d'une conclusion de dissertation (Probatoire / Bac technique)]
Une conclusion réussie se construit en \textbf{trois temps obligatoires} :
\begin{enumerate}
\item \textbf{Rappel de la problématique et de la thèse :} reprenez la question posée en introduction et réaffirmez votre réponse principale en une phrase forte. \emph{Exemple : « Au terme de cette étude, il apparaît que le partage de l'Afrique (1880-1914) n'a pas été une simple occupation militaire, mais une restructuration violente des sociétés africaines au profit du capitalisme européen. »}
\item \textbf{Synthèse des grandes parties (sans idée nouvelle) :} résumez en deux ou trois phrases l'articulation de votre démonstration, avec des connecteurs logiques. \emph{Exemple : « Si la conquête s'est heurtée à des résistances variées (I), elle a imposé des modèles administratifs et économiques distincts mais convergents dans l'exploitation (II), dont les frontières artificielles et les dualismes — illustrés par le cas camerounais — constituent l'héritage le plus lourd pour les États postcoloniaux (III). »}
\item \textbf{Ouverture maîtrisée :} élargissez la réflexion sans sortir du sujet. Trois types d'ouverture sont valorisés :
\begin{itemize}
\item \textbf{chronologique :} vers la période suivante (« Cette colonisation achevée en 1914 entre en crise dès l'entre-deux-guerres avec la montée des nationalismes... ») ;
\item \textbf{comparative :} vers un autre espace (« Ce modèle colonial trouve des parallèles en Asie du Sud-Est... ») ;
\item \textbf{actuelle :} vers un débat contemporain (« Les débats sur la restitution des œuvres d'art montrent que ce bilan reste vif... »).
\end{itemize}
\end{enumerate}
\emph{Astuce :} rédigez votre conclusion au brouillon avant l'introduction, pour être sûr de tenir votre thèse.
\end{methode}

\begin{exoresolu}[Entraînement : une conclusion type]
\textbf{Sujet :} « Les conséquences du partage de l'Afrique (1885-1914) sur les sociétés africaines. »

\textbf{Problématique :} en quoi le partage de l'Afrique a-t-il constitué une rupture structurelle dont les effets dépassent la période coloniale ?
\tcblower
\textbf{Conclusion rédigée :}

\emph{En définitive, le partage de l'Afrique (1885-1914) a opéré une rupture majeure dans l'histoire du continent. \textbf{(Rappel de la thèse)} La conquête militaire, facilitée par la supériorité technique européenne et par les divisions africaines, a imposé des frontières artificielles qui ont éclaté les entités politiques et humaines précoloniales. L'économie coloniale, qu'elle soit d'exploitation ou de peuplement, a inséré le continent dans la division internationale du travail comme pourvoyeur de matières premières, au prix d'une contrainte démographique et sociale sans précédent (travail forcé, fiscalité, spoliation foncière). \textbf{(Synthèse des parties)} Toutefois, cette domination a engendré ses propres contestataires : l'école coloniale et l'administration ont formé une élite occidentalisée qui, dès l'entre-deux-guerres, retournera les principes de la « mission civilisatrice » pour réclamer l'égalité des droits, puis l'indépendance, ouvrant la période des nationalismes. \textbf{(Ouverture chronologique)}}
\end{exoresolu}

\begin{exercice}[Entraînement personnel]
Rédigez une conclusion structurée en trois temps pour le sujet suivant : \emph{« L'administration coloniale en Afrique : entre administration directe et indirecte, quelles conséquences politiques pour les sociétés africaines ? »} Vous utiliserez les notions d'\emph{indirect rule}, d'administration directe, de chefferie et d'élites.
\end{exercice}

\begin{corrige}[Exercice : conclusion sur l'administration coloniale]
\emph{Au total, l'opposition entre administration directe française et \emph{indirect rule} britannique masque une même finalité : le contrôle politique des populations à moindre coût. \textbf{(Thèse)} La France a affaibli les chefferies traditionnelles pour imposer une administration centralisée, créant une élite assimilée mais déracinée ; la Grande-Bretagne a instrumentalisé les chefs coutumiers, figeant les identités ethniques et freinant l'émergence d'une conscience nationale unifiée. \textbf{(Synthèse)} Ces deux héritages éclairent les difficultés post-indépendance : crise de légitimité de l'État dans nombre de pays francophones, fédéralisme et conflits ethniques dans plusieurs États anglophones (Nigeria) ; au Cameroun, la réunification de 1961 n'a pas effacé le dualisme juridique et administratif hérité du double mandat. \textbf{(Ouverture actuelle)}}
\end{corrige}

\newpage

\coursec{Activité d'intégration}

\coursec{L'impérialisme en Afrique}

\courssub{Objectifs de l'activité}
\begin{itemize}
    \item \textbf{Savoirs :} Maîtriser la géographie du partage colonial de l'Afrique en 1914 (puissances coloniales, étendues territoriales, exceptions : Libéria, Éthiopie) ; identifier les grands foyers de résistance armée et les zones de collaboration ; comprendre les arguments idéologiques (discours de Jules Ferry) et juridiques (traité Germano-Douala) de la colonisation.
    \item \textbf{Savoir-faire :} Construire un croquis cartographique historique rigoureux (choix du fond de carte, code couleur, légende hiérarchisée, titre, échelle, orientation) ; exploiter un dossier documentaire hétérogène (carte, texte, photographie, tableau statistique) en croisant les informations ; rédiger une réponse argumentée structurée (problématique, thèse, arguments appuyés sur des documents, conclusion) en mobilisant le vocabulaire spécifique (\cle{impérialisme}, \cle{conférence de Berlin}, \cle{résistance}, \cle{collaboration}, \cle{pacification}).
    \item \textbf{Attitudes :} Développer l'esprit critique face aux justifications officielles de la conquête ; travailler en équipe pour mutualiser les analyses documentaires ; respecter la rigueur méthodologique de la discipline historique.
\end{itemize}

\courssub{Situation}
\begin{popfigure}
\begin{tikzpicture}[scale=0.55, every node/.style={font=\footnotesize}]
% Fond de carte Afrique très simplifié (polygone approximatif)
\draw[popline, line width=1.2pt] 
(0,0) -- (10,0) -- (12,2) -- (13,5) -- (12,8) -- (11,11) -- (9,13) -- (7,14) -- (5,13.5) -- (3,12) -- (1,10) -- (0,7) -- (0,3) -- cycle;
% Méridiens/Parallèles repères
\draw[popmuted, dashed] (2,0) -- (2,14); \draw[popmuted, dashed] (6,0) -- (6,14); \draw[popmuted, dashed] (10,0) -- (10,14);
\draw[popmuted, dashed] (0,3) -- (13,3); \draw[popmuted, dashed] (0,7) -- (13,7); \draw[popmuted, dashed] (0,11) -- (12,11);
% Villes repères (points)
\foreach \x/\y/\name in {2.5/1.5/Lagos, 5.5/2.5/Kinshasa, 8/1.5/Luanda, 3/5/Douala, 7/5/Nairobi, 10/5/Mogadiscio, 4/9/LeCaire, 9/9/Khartoum, 2/11/Tunis, 11/11/Djibouti, 5/12.5/Alger}
    \fill[popdark] (\x,\y) circle (2pt) node[above right] {\name};
% Flèche Nord
\draw[->, popblue, line width=1.5pt] (11.5, 13.5) -- (11.5, 14.5) node[above] {N};
% Échelle indicative
\draw[popdark, line width=1.5pt] (0.5, -0.5) -- (3.5, -0.5) node[midway, below] {1000 km};
\end{tikzpicture}
\legende{Fond de carte vierge de l'Afrique (simplifié) fourni aux élèves pour le croquis. Les repères de grille, les capitales actuelles et l'échelle sont indiqués pour faciliter le report.}
\end{popfigure}

Dans le cadre de la préparation de l'exposition « \emph{Cameroun 1884-1914 : de la traite à l'administration coloniale} » au Musée National de Yaoundé, votre classe de Première Technique est chargée de réaliser le panneau central : \emph{« L'Afrique partagée en 1914 : conquêtes, résistances et collaborations »}. Le conservateur vous remet un \textbf{dossier documentaire} composé de cinq pièces :
\begin{enumerate}
    \item \textbf{Doc 1 :} Deux cartes muettes de l'Afrique (1880 et 1914) à l'échelle identique.
    \item \textbf{Doc 2 :} Extrait du \textbf{Traité d'amitié et de commerce signé le 12 juillet 1884 entre le Roi Bell (Douala) et l'Empire allemand} (article 3 : « Les rois et chefs s'engagent à ne céder aucun terrain à une autre puissance... » ; article 5 : « L'Empire allemand assure sa protection aux rois et chefs... »).
    \item \textbf{Doc 3 :} \textbf{Photographie de Samory Touré} (vers 1898), chef mandingue, capturé par le capitaine Gouraud, exilé au Gabon où il meurt en 1900.
    \item \textbf{Doc 4 :} \textbf{Tableau comparatif des superficies des empires coloniaux en Afrique en 1914} (France : 10,5 millions km² ; Royaume-Uni : 9,2 millions km² ; Allemagne : 2,5 millions km² ; Belgique : 2,3 millions km² ; Portugal : 2,1 millions km² ; Italie : 1,9 million km² ; Espagne : 0,3 million km² ; Total Afrique : env. 30 millions km²).
    \item \textbf{Doc 5 :} Extrait du \textbf{discours de Jules Ferry à la Chambre des députés (28 juillet 1885)} : « \emph{Il faut dire ouvertement que les races supérieures ont un droit vis-à-vis des races inférieures... parce qu'il y a un devoir pour elles de les civiliser.} »
\end{enumerate}
Votre mission : produire le panneau d'exposition (carte légendée + commentaire argumenté) pour l'inauguration dans trois séances.

\courssub{Consignes}
\begin{enumerate}
    \item \textbf{Construction du croquis cartographique (Travail de groupe, 1h30) :}
    \begin{enumerate}
        \item À partir du fond de carte vierge (Doc 1, carte 1914), reporter les limites des \textbf{sept empires coloniaux} (France, Royaume-Uni, Allemagne, Belgique, Portugal, Italie, Espagne) et les deux États indépendants (Libéria, Éthiopie) en utilisant un \textbf{code couleur normalisé} (ex. : France = bleu, UK = rouge, Allemagne = gris, Belgique = orange, Portugal = vert, Italie = brun, Espagne = jaune, Indépendants = blanc hachuré).
        \item Réaliser une \textbf{légende organisée} (titre, échelles, symboles de surface, symboles ponctuels pour les capitales coloniales : Dakar, Le Caire, Nairobi, Windhoek, etc., flèches des principales routes de pénétration).
        \item Localiser par des \textbf{symboles distincts} (étoile rouge pour résistance, main serrée verte pour collaboration) : \textbf{trois foyers de résistance majeure} (ex. : Samory en Afrique de l'Ouest, Behanzin au Dahomey, Menelik II en Éthiopie, Mekatilili au Kenya, résistance Herero/Nama au Sud-Ouest africain) et \textbf{deux zones de collaboration significative} (ex. : Côte des Douala au Cameroun, Royaumes du Buganda en Ouganda, Chefferies du Sénégal des « originaires »).
    \end{enumerate}
    \item \textbf{Rédaction du commentaire argumenté (Travail individuel, 1h) :}
    \begin{enumerate}
        \item Répondre à la problématique : \textbf{« La conquête de l'Afrique fut-elle une pacification ou une guerre ? »}
        \item Rédiger \textbf{environ 20 lignes} (180-220 mots).
        \item \textbf{Mobiliser obligatoirement au moins trois documents} du dossier (les citer explicitement : « D'après le Doc 2... », « Le Doc 3 montre... », « Le Doc 4 révèle... »).
        \item Structurer la réponse : Introduction (accroche, définition de la \cle{conférence de Berlin}, énoncé de la thèse nuancée), Développement en deux arguments équilibrés (argument 1 : violence de la conquête/guerres de résistance ; argument 2 : stratégies de collaboration/« pacification » administrative), Conclusion (bilan : conquête par les armes masquée par rhétorique civilisatrice).
        \item Justifier votre démarche en marge : « J'ai croisé le Doc X (source) et le Doc Y (donnée chiffrée) pour démontrer Z. »
    \end{enumerate}
\end{enumerate}

\courssub{Ressources à mobiliser}
\begin{methode}[Méthode du croquis cartographique historique (3 étapes)]
    \begin{enumerate}
        \item \textbf{Analyse du sujet \& Choix de la base :} Identifier la date de référence (1914), le thème (partage colonial), l'échelle. Choisir un fond de carte précis (hydrographie, relief, frontières actuelles en grisé).
        \item \textbf{Hiérarchisation de l'information \& Code graphique :} Classer les phénomènes par ordre d'importance (Empires = surfaces/coloration ; Capitales = points ; Résistances = étoiles ; Collaborations = pictogrammes ; Flux = flèches). Respecter la sémiologie graphique (Bertin) : teintes claires pour grandes surfaces, teintes foncées pour petites ; hachures pour exceptions.
        \item \textbf{Réalisation \& Légende :} Tracer au crayon puis encre. Rédiger une légende \textbf{thématique et hiérarchisée} (I. Empires coloniaux ; II. États souverains ; III. Foyers de résistance ; IV. Zones de collaboration ; V. Axes de pénétration). Indiquer Titre, Auteur, Date, Source, Échelle, Orientation (TADSEO).
    \end{enumerate}
\end{methode}

\begin{aretenir}[Repères chronologiques et conceptuels clés]
    \begin{itemize}
        \item \textbf{1884-1885 :} Conférence de Berlin (règles du partage, libre navigation Congo/Niger, occupation effective).
        \item \textbf{1884 :} Traité Germano-Douala (naissance du Kamerun) ; Traité de protectorat français sur le Gabon/Congo.
        \item \textbf{1890-1898 :} Guerres de résistance majeures (Samory, Behanzin, Rabah, Menelik II/Adoua 1896).
        \item \textbf{1914 :} Quasi-totalité de l'Afrique colonisée (sauf Libéria, Éthiopie).
        \item \textbf{Concepts :} \cle{Impérialisme} (expansion capitaliste, Lénine/Hobson), \cle{Mission civilisatrice} (idéologie), \cle{Collaboration} (alliance tactique chefferies/colonisateur), \cle{Résistance} (primaire/secondaire, armée/religieuse).
    \end{itemize}
\end{aretenir}

\begin{attention}[Pièges à éviter]
\begin{enumerate}
    \item Ne pas confondre \textbf{1880} (pré-colonial, comptoirs côtiers) et \textbf{1914} (partage achevé).
    \item Ne pas colorier l'Éthiopie et le Libéria comme colonies (hachures blanches obligatoires).
    \item Ne pas écrire « pacification » sans guillemets ni analyse critique : c'est un terme de propagande coloniale (\cle{euphémisme}).
    \item Oublier l'échelle ou le nord dans le croquis = pénalité.
    \item Réponse hors-sujet : raconter l'histoire de Samory sans répondre à la problématique « pacification ou guerre ».
\end{enumerate}
\end{attention}

\courssub{Démarche et corrigé commenté}
\begin{exoresolu}[Corrigé type de l'activité d'intégration]
\textbf{1. Construction du croquis : « L'Afrique partagée en 1914 »}
\tcblower
\begin{popfigure}
\begin{tikzpicture}[scale=0.55, every node/.style={font=\tiny}]
% --- FOND DE CARTE AFRIQUE (Contour simplifié) ---
\draw[popline, line width=1pt] 
(0,0) -- (10,0) -- (12,2) -- (13,5) -- (12,8) -- (11,11) -- (9,13) -- (7,14) -- (5,13.5) -- (3,12) -- (1,10) -- (0,7) -- (0,3) -- cycle;
% Grille discrète (repères)
\draw[popmuted!30, dashed] (2,0) -- (2,14); \draw[popmuted!30, dashed] (6,0) -- (6,14); \draw[popmuted!30, dashed] (10,0) -- (10,14);
\draw[popmuted!30, dashed] (0,3) -- (13,3); \draw[popmuted!30, dashed] (0,7) -- (13,7); \draw[popmuted!30, dashed] (0,11) -- (12,11);
% --- ZONES COLONIALES (Bloc rectangulaires schématiques non chevauchants) ---
% Couleurs : France=popblue, UK=popred, Allemagne=popmuted, Belgique=poporange, Portugal=popgreen, Italie=poppurple, Espagne=popgold
% FRANCE (Bleu) - AOF, AEF, Maghreb, Madagascar, Djibouti
\fill[popblue!20] (0,7) -- (3,7) -- (3,10) -- (1,10) -- (0,7) -- cycle; % AOF Ouest (Mauritanie, Sénégal, Soudan fr.)
\fill[popblue!20] (3,7) -- (6,7) -- (6,10) -- (3,10) -- cycle; % AOF Est / AEF Nord (Tchad, Oubangui)
\fill[popblue!20] (0,10) -- (4,10) -- (4,12) -- (1,12) -- cycle; % Maghreb Ouest (Maroc, Algérie, Tunisie)
\fill[popblue!20] (2,2) -- (4,2) -- (4,4) -- (2,4) -- cycle; % AEF Sud (Gabon, Moyen-Congo) - *Corrigé: ne chevauche pas Cameroun*
\fill[popblue!20] (10.5, -0.5) rectangle (12.5, 1.5); % Madagascar (décalé hors cadre principal pour lisibilité)
\fill[popblue!20] (11, 10.5) circle (0.3); % Djibouti (Côte des Somalis)
% ROYAUME-UNI (Rouge) - Égypte, Soudan, Afrique Est, Afrique Australe, Afrique Ouest (Nigeria, Gold Coast), Afrique Centrale (Rhodesies)
\fill[popred!20] (0,1) -- (3,1) -- (3,3) -- (0,3) -- cycle; % AOF Britannique (Gambie, Sierra Leone, Gold Coast, Nigeria Sud)
\fill[popred!20] (6,1) -- (9,1) -- (9,4) -- (6,4) -- cycle; % Afrique Australe (Afrique du Sud, Bechuanaland, Rhodésie Sud)
\fill[popred!20] (6,4) -- (9,4) -- (9,7) -- (6,7) -- cycle; % Afrique Centrale Britannique (Rhodésie Nord, Nyassaland)
\fill[popred!20] (7,7) -- (10,7) -- (10,10) -- (7,10) -- cycle; % Afrique Est (Kenya, Ouganda)
\fill[popred!20] (7,8) -- (10,8) -- (10,11) -- (7,11) -- cycle; % Soudan Anglo-Égyptien / Égypte (chevauchement volontaire Soudan/Égypte)
% ALLEMAGNE (Gris) - Cameroun, Togo, Tanganyika, Sud-Ouest Africain
\fill[popmuted!40] (3,2) -- (5,2) -- (5,4) -- (3,4) -- cycle; % Kamerun (zone Douala/Intérieur)
\fill[popmuted!40] (2.5,1.5) -- (3,1.5) -- (3,2.5) -- (2.5,2.5) -- cycle; % Togo
\fill[popmuted!40] (8,3) -- (10,3) -- (10,6) -- (8,6) -- cycle; % Afrique Est Allemande (Tanganyika)
\fill[popmuted!40] (5,0.5) -- (7,0.5) -- (7,2) -- (5,2) -- cycle; % Sud-Ouest Africain (Namibie)
% BELGIQUE (Orange) - Congo Belge
\fill[poporange!30] (4,1) -- (7,1) -- (7,4) -- (4,4) -- cycle; % Congo Belge (au sud du Cameroun, à l'est de l'AEF)
% PORTUGAL (Vert) - Angola, Mozambique, Guinée, Cap Vert, São Tomé
\fill[popgreen!20] (4,0) -- (6,0) -- (6,2.5) -- (4,2.5) -- cycle; % Angola (au nord de la Namibie, à l'ouest du Congo)
\fill[popgreen!20] (9,1) -- (11,1) -- (11,4) -- (9,4) -- cycle; % Mozambique (au sud du Tanganyika, à l'est de la Rhodésie)
% ESPAGNE (Jaune) - Maroc Nord, Sahara Occidental, Guinée Espagnole
\fill[popgold!40] (1,9) -- (2,9) -- (2,10) -- (1,10) -- cycle; % Maroc Nord (Rif)
\fill[popgold!40] (1,5) -- (2,5) -- (2,6) -- (1,6) -- cycle; % Río de Oro / Sahara Occidental
\fill[popgold!40] (3.2,1.8) -- (3.7,1.8) -- (3.7,2.3) -- (3.2,2.3) -- cycle; % Guinée Espagnole (Rio Muni)
% ITALIE (Brun/Violet) - Libye, Érythrée, Somalie
\fill[poppurple!20] (5,9) -- (8,9) -- (8,12) -- (5,12) -- cycle; % Libye (Tripolitaine, Cyrénaïque, Fezzan)
\fill[poppurple!20] (10.5, 6) rectangle (11.5, 8); % Érythrée / Somalie (Corne de l'Afrique)
% INDÉPENDANTS (Blanc hachuré) - Libéria, Éthiopie
\fill[pattern=north east lines, pattern color=popdark] (1.5, 1.5) rectangle (3, 3); % Libéria (dans golfe de Guinée)
\fill[pattern=north east lines, pattern color=popdark] (9, 7) rectangle (11, 9); % Éthiopie (Corne de l'Afrique, bordée par UK, IT, FR)
% --- SYMBOLES PONCTUELS : Capitales / Villes clés ---
\foreach \x/\y/\name/\col in {
    1.5/2.5/Bathurst/red, 2.5/2/Freetown/red, 3/2.5/Accra/red, 3.5/2/Lagos/red, % UK Afrique Ouest
    3.5/2.5/Douala/gray, % All. Cameroun
    2.7/2/Lomé/gray, % All. Togo
    5.5/1.5/Luanda/green, % Pt. Angola
    5.5/2.5/Kinshasa/orange, % Bel. Congo
    4/3.5/Brazzaville/blue, 3.5/3/Libreville/blue, % Fr. AEF
    1.5/5/Dakar/blue, 3/5/Abidjan/blue, 4/5/Ouagadougou/blue, % Fr. AOF
    2.5/9.5/Alger/blue, 2/10.5/Tunis/blue, 1.5/10/Rabat/blue, % Fr. Maghreb
    4/9.5/LeCaire/red, 8.5/9.5/Khartoum/red, % UK Égypte/Soudan
    7.5/5.5/Nairobi/red, 8.5/6.5/Kampala/red, % UK Afrique Est
    10.5/5.5/Mogadiscio/purple, % It. Somalie
    11/10.5/Djibouti/blue, % Fr. Djibouti
    9.5/8/AddisAbeba/popdark, % Éthiopie
    5.5/0.5/Windhoek/gray, % All. Namibie
    7.5/2.5/Keetmanshoop/gray, % All. Namibie
    9/2.5/DarEsSalaam/gray, % All. Tanganyika
    10.5/1.5/LourencoMarques/green, % Pt. Mozambique
    1.5/3.5/Monrovia/popdark} % Libéria
    \fill[\col] (\x,\y) circle (1.5pt) node[above right, font=\tiny] {\name};
% --- RÉSISTANCES (Étoiles rouges) ---
\foreach \x/\y/\name in {2.5/2.5/Samory, 3/1.5/Behanzin, 9.5/8/MenelikII, 8/5.5/Mekatilili, 5.5/0.5/HereroNama}
    \node[star, star points=5, star point ratio=2.25, fill=poppink, draw=popdark, inner sep=1.5pt] at (\x,\y) {} 
    node[below, font=\tiny, text=popdark, align=center] {\name};
% --- COLLABORATIONS (Cercles verts / Poignée de main symbolique) ---
\foreach \x/\y/\name in {3.5/2.5/Douala, 8.5/6.5/Buganda, 1.5/3.5/Senegal}
    \node[circle, draw=popgreen, thick, inner sep=2pt] at (\x,\y) {} 
    node[above, font=\tiny, text=popgreen, align=center] {\name};
% --- FLÈCHES PÉNÉTRATION (exemples majeurs) ---
\draw[->, popline, thick] (1.5,2) -- (3,4); % Sénégal -> Soudan/Niger
\draw[->, popline, thick] (3.5,2) -- (4,3.5); % Côte ivoire/Ghana -> Burkina/Niger
\draw[->, popline, thick] (3.5,2.5) -- (4,3.5); % Douala -> Intérieur Cameroun
\draw[->, popline, thick] (5.5,1.5) -- (5.5,3.5); % Luanda -> Intérieur Angola
\draw[->, popline, thick] (9.5,2) -- (9.5,4.5); % Mozambique -> Intérieur
\draw[->, popline, thick] (10,5.5) -- (9.5,7); % Côte Somalie/Kenya -> Lacs
\draw[->, popline, thick] (7,9) -- (7,11); % Égypte -> Soudan
% Nord & Échelle
\draw[->, popblue, line width=1pt] (11.5, 13.5) -- (11.5, 14.5) node[above, font=\small] {N};
\draw[popdark, line width=1pt] (0.5, -0.8) -- (3.5, -0.8) node[midway, below, font=\tiny] {1000 km};
\end{tikzpicture}
\legende{Croquis corrigé : Le partage de l'Afrique en 1914. \textbf{Légende hiérarchisée :} \textbf{I. Empires coloniaux (Surfaces colorées)} : \textcolor{popblue}{■} France (10,5 M km² : AOF, AEF, Maghreb, Madagascar, Djibouti), \textcolor{popred}{■} Royaume-Uni (9,2 M km² : Égypte, Soudan, Afrique de l'Est, Australe, Occidentale, Centrale), \textcolor{popmuted}{■} Allemagne (2,5 M km² : Kamerun, Togo, Tanganyika, Namibie), \textcolor{poporange}{■} Belgique (2,3 M km² : Congo), \textcolor{popgreen}{■} Portugal (2,1 M km² : Angola, Mozambique, Guinée), \textcolor{poppurple}{■} Italie (1,9 M km² : Libye, Érythrée, Somalie), \textcolor{popgold}{■} Espagne (0,3 M km² : Maroc N., Sahara, Guinée Eq.). \textbf{II. États souverains} : \textcolor{popdark}{▨} Libéria, Éthiopie. \textbf{III. ★ Foyers de résistance armée} : Samory (Afrique O.), Behanzin (Dahomey), Menelik II (Éthiopie/Adoua 1896), Mekatilili (Kenya), Herero/Nama (Namibie). \textbf{IV. ◎ Zones de collaboration} : Douala (Cameroun), Buganda (Ouganda), Sénégal (Quatre Communes). \textbf{V. ➤ Axes de pénétration coloniale} (fleuves, chemins de fer, pistes militaires). \textbf{TADSEO} : Titre, Auteur, Date, Source (Doc 1, Doc 4, atlas), Échelle, Orientation.}
\end{popfigure}

\textbf{2. Commentaire argumenté : « La conquête de l'Afrique fut-elle une pacification ou une guerre ? »}

\textit{Introduction.} La Conférence de Berlin (1884-1885) acte le partage de l'Afrique entre puissances européennes sous couvert de « mission civilisatrice » et de libre-échange. Pourtant, l'observation de la carte de 1914 (Doc 1) et des archives révèle que cette \cle{pacification} proclamée fut en réalité une \textbf{conquête militaire brutale}, ponctuée de résistances acharnées, même si des collaborations tactiques ont facilité l'entreprise.

\textit{Argument 1 : Une conquête par les armes, masquée par l'idéologie civilisatrice.} Le discours de Jules Ferry (Doc 5) expose la doctrine raciste fondatrice : « les races supérieures ont un droit vis-à-vis des races inférieures ». Cette justification idéologique légitime l'usage de la force. Le tableau des superficies (Doc 4) montre que la France et le Royaume-Uni se partagent près de 20 millions de km², soit les deux tiers du continent, conquis en moins de trente ans. Cette rapidité n'a été possible que par des expéditions militaires massives (colonnes Voulet-Chanoine, Marchand, Kitchener). La photographie de Samory Touré (Doc 3), dernier grand résistant mandingue capturé en 1898 après 16 ans de guerre, incarne la violence de cette « pacification » : son empire fut détruit, les populations déportées, lui-même exilé. De même, la résistance Herero et Nama (1904-1908) en Afrique du Sud-Ouest allemande tourne au premier génocide du XXe siècle. La \cle{pacification} est donc un \textbf{euphémisme} désignant la guerre de conquête.

\textit{Argument 2 : La collaboration, stratégie africaine et levier colonial, ne saurait effacer la violence structurelle.} Le traité Germano-Douala de 1884 (Doc 2) illustre la \cle{collaboration} : les rois Bell et Akwa acceptent le protectorat allemand (art. 5 : « protection ») en échange du maintien de leurs privilèges commerciaux (art. 3 : monopole foncier). Au Buganda (Ouganda), les chefs protestants s'allient aux Britanniques contre les rivaux musulmans et le roi Mwanga. Mais ces alliances sont inégales : elles transforment les chefs en \cle{auxiliaires de l'administration} (chefs de canton, percepteurs d'impôts), brisant la souveraineté africaine. La collaboration est souvent une \textbf{stratégie de survie} face à la supériorité technique (fusils à répétition, artillerie, Maxim) ou une rivalité inter-africaine exploitée par l'Européen (diviser pour régner).

\textit{Conclusion.} En 1914, l'Afrique est quadrillée (carte Doc 1), ses ressources pillées, ses sociétés bouleversées. La « pacification » fut le \textbf{nom de code de la guerre coloniale}. Les résistances (Samory, Behanzin, Menelik II, Adoua 1896, Maji Maji 1905) prouvent que les Africains n'ont pas accepté passivement le partage. Les collaborations (Douala, Buganda, Sénégal) furent des accommodements tactiques dans un rapport de force imposé par la violence. L'historien doit donc refuser le vocabulaire du vainqueur : ce fut une \textbf{guerre de conquête impérialiste}.

\textbf{Justification de la démarche (marge) :} J'ai croisé le \textbf{Doc 5} (source idéologique, discours officiel) et le \textbf{Doc 3} (source iconographique, réalité du terrain : chef vaincu, exilé) pour démontrer le décalage entre rhétorique civilisatrice et violence effective. J'ai utilisé le \textbf{Doc 4} (données chiffrées, surfaces) pour quantifier l'ampleur territoriale de cette conquête militaire. Le \textbf{Doc 2} (traité) a permis de nuancer : la collaboration existe mais sous contrainte de l'occupation effective (Conférence de Berlin).
\end{exoresolu}

\courssub{Bilan de l'activité}
Cette activité d'intégration a permis de valider les acquis de la leçon \emph{L'impérialisme en Afrique} par une mise en situation professionnelle (médiation culturelle muséale).
\begin{itemize}
    \item \textbf{Sur le plan méthodologique :} Les élèves ont appliqué la \cle{méthode du croquis historique} (choix de l'information, sémiologie graphique, légende hiérarchisée TADSEO). Ils ont compris qu'une carte n'est pas une illustration mais une \textbf{analyse spatiale} : le code couleur révèle l'hégémonie franco-britannique, les hachures l'exception éthiopienne/libérienne, les symboles la géographie de la résistance (sahélienne, forestière, est-africaine) vs collaboration (côtière, lacustre).
    \item \textbf{Sur le plan historique :} Le croisement des documents a fait émerger la \textbf{complexité du fait colonial} : ni « don de la civilisation » (thèse coloniale), ni « nuit sans aube » (thèse nationaliste simpliste), mais un affrontement asymétrique où la violence militaire (Doc 3, 4) est le fondement, l'idéologie raciste (Doc 5) le masque, et la collaboration (Doc 2) un rouage de l'administration indirecte. La problématique « pacification ou guerre » a forcé la \textbf{nuance argumentée} : la « pacification » est le terme administratif pour désigner la guerre de conquête et la répression des résistances.
    \item \textbf{Compétences transférables :} Analyse de sources hétérogènes (texte juridique, discours politique, photographie, statistique, carte), rédaction argumentée contrainte (nombre de lignes, citation de docs), travail collaboratif (croquis) puis individuel (rédaction). Ces compétences sont directement réinvestissables en ECE (Épreuve Comptée pour l'Examen) et pour le Grand Oral du Baccalauréat Technique.
\end{itemize}

\begin{savaistu}[Le Cameroun, cas d'école du partage]
Le Cameroun illustre parfaitement la leçon : \textbf{1884}, traité Germano-Douala (Doc 2) $\rightarrow$ Kamerun allemand. \textbf{1911}, Traité de Fès (Agadir) $\rightarrow$ France reçoit le « Neukamerun » (Est/Sud) contre Maroc. \textbf{1916}, Conquête franco-britannique $\rightarrow$ Partage mandat SDN (1919) : 4/5 France, 1/5 UK. \textbf{1960-61}, Indépendances et Réunification. Le croquis de 1914 montre le Kamerun allemand entier (gris) ; la carte de 1920 montrerait la ligne de partage franco-britannique traversant les pays Bamiléké et Bassa. C'est l'exemple type de la \textbf{frontière artificielle} issue de la Conférence de Berlin.
\end{savaistu}

\newpage

\coursec{Méthodes et exercices résolus}

\courssub{Méthodes et exercices résolus}

\begin{methode}[Analyser une carte historique du partage de l'Afrique (1885-1914)]
\textbf{Objectif :} Appliquer les notions de l'unité (conquêtes, partages, résistances) à la lecture d'un document cartographique type Bac.
\begin{enumerate}
    \item \textbf{Identifier le document :} Nature (carte), Titre, Date, Échelle, Légende, Auteur/Source.
    \item \textbf{Analyser la légende :} Repérer les codes couleurs (puissances coloniales), les symboles (capitales, chemins de fer, zones de résistance), les dates clés (Conférence de Berlin 1884-1885, Conférence de Bruxelles 1890, Fachoda 1898, Accords de 1904-1911).
    \item \textbf{Décrire la répartition :} Citer les \cle{sept puissances} coloniales (France, Royaume-Uni, Allemagne, Belgique, Portugal, Italie, Espagne) et leurs blocs territoriaux principaux (AOF, AEF, Afrique du Sud, Afrique de l'Est, Congo, etc.).
    \item \textbf{Expliquer les logiques :} Relier les possessions aux stratégies (française : \emph{est-ouest} / britannique : \emph{nord-sud} / allemande : \emph{morcellement} / belge : \emph{bassin du Congo}).
    \item \textbf{Conclure :} Résumer le bilan (90\% du continent colonisé en 1914, deux États indépendants : Libéria, Éthiopie) et ouvrir sur les conséquences (tracés artificiels, économies d'extraction).
\end{enumerate}
\textbf{Reflexe Bac :} Toujours citer des \textbf{dates précises} et des \textbf{noms de traités/conférences} pour appuyer la description. Ne pas se contenter de lister les couleurs.
\end{methode}

\begin{exoresolu}[Commentaire de carte : Le partage de l'Afrique en 1914]
\emph{Énoncé :} À l'aide de la carte ci-dessous et de vos connaissances, présentez le partage de l'Afrique entre les puissances européennes en 1914 et expliquez les logiques d'expansion de la France et du Royaume-Uni.
\begin{center}
\begin{popfigure}
\begin{tikzpicture}[scale=0.75]
% Contour Afrique simplifié
\draw[popdark, thick] (0,0) .. controls (2,-1) and (5,-2) .. (7,0) .. controls (9,2) and (10,5) .. (8,8) .. controls (6,10) and (3,10) .. (1,8) .. controls (-1,6) and (-1,3) .. (0,0);
% Méditerranée
\draw[popblue, thick] (-1,8) -- (10,8);
% Équateur
\draw[popmuted, dashed] (-1,4) -- (10,4);
% Zones colorées (simplifiées)
% France (AOF + AEF + Maghreb + Madagascar)
\fill[popblue!30] (0,0) .. controls (2,-1) and (5,-2) .. (7,0) .. controls (6,2) and (4,3) .. (2,4) .. controls (0,3) and (-1,2) .. (0,0); % AOF
\fill[popblue!30] (2,4) .. controls (3,5) and (4,5) .. (5,6) .. controls (4,5.5) and (2,5) .. (1,4.5) .. controls (1.5,4) .. (2,4); % AEF part
\fill[popblue!30] (5,6) .. controls (6,6.5) and (7,7) .. (6,8) .. controls (4,7.5) and (3,6.5) .. (2,6) .. controls (2.5,5.5) .. (5,6); % Tchad/Centrafrique
\fill[popblue!30] (7,7.5) rectangle (9.5,8.5); % Tunisie/Algérie (nord)
\node[popblue, font=\bfseries\small] at (1.5,1.5) {France (AOF)};
\node[popblue, font=\bfseries\small] at (3.5,5.5) {France (AEF)};
\node[popblue, font=\bfseries\small] at (8.2,8) {France (Maghreb)};
% UK (Égypte-Soudan, Afrique du Sud, Nigeria, Kenya/Ouganda)
\fill[popred!30] (7,0) .. controls (8,1) and (8,3) .. (7,4) .. controls (6.5,3.5) and (6,2) .. (7,0); % Nigeria
\fill[popred!30] (6,4) .. controls (7,5) and (7.5,6) .. (6.5,7) .. controls (5.5,6.5) and (5,5.5) .. (6,4); % Soudan
\fill[popred!30] (7.5,7) rectangle (9.5,8); % Égypte
\fill[popred!30] (1,1) .. controls (2,1.5) and (2,0.5) .. (1,0) .. controls (0.5,0.5) .. (1,1); % Sierra Leone/Gambie (petit)
\fill[popred!30] (5,1) rectangle (6.5,3); % Afrique du Sud (simplifié bas)
\node[popred, font=\bfseries\small] at (7.5,1.5) {R.-U. (Nigeria)};
\node[popred, font=\bfseries\small] at (6.5,5.5) {R.-U. (Soudan/Égypte)};
\node[popred, font=\bfseries\small] at (5.5,2) {R.-U. (Af. du Sud)};
% Allemagne (Togo, Cameroun, Namibie, Tanzanie, Rwanda/Burundi)
\fill[popgray!40] (3.5,1.5) rectangle (4.5,2.5); % Togo
\fill[popgray!40] (3.5,2.5) rectangle (5,4); % Cameroun
\fill[popgray!40] (5.5,0.5) rectangle (7,1.5); % Namibie (SW)
\fill[popgray!40] (6,4.5) rectangle (7.5,6.5); % Afrique Est
\node[popdark, font=\bfseries\small] at (4,2) {All. (Togo)};
\node[popdark, font=\bfseries\small] at (4.2,3.2) {All. (Kamerun)};
\node[popdark, font=\bfseries\small] at (6.2,1) {All. (SW Af)};
\node[popdark, font=\bfseries\small] at (6.7,5.5) {All. (A.O. Allemande)};
% Belgique (Congo)
\fill[popgold!50] (4,3.5) .. controls (5,4) and (5,5) .. (4,5.5) .. controls (3,5) and (3,4) .. (4,3.5);
\node[popdark, font=\bfseries\small] at (4,4.5) {Belgique (Congo)};
% Portugal (Angola, Mozambique, Guinée)
\fill[popgreen!30] (3,0.5) rectangle (4,2); % Angola (part)
\fill[popgreen!30] (6.5,1.5) rectangle (8,3.5); % Mozambique (part)
\node[popgreen!60!black, font=\bfseries\small] at (3.5,1.2) {Portugal};
\node[popgreen!60!black, font=\bfseries\small] at (7.2,2.5) {Portugal};
% Italie (Libye, Erythrée, Somalie)
\fill[poporange!40] (8,6.5) rectangle (9.5,7.5); % Libye
\fill[poporange!40] (7.5,4) rectangle (8.5,5); % Corne de l'Afrique
\node[poporange!80!black, font=\bfseries\small] at (8.7,7) {Italie};
% Espagne (Rio de Oro, Guinée Eq, Maroc N)
\fill[poppurple!30] (1.5,7.5) rectangle (2.5,8.5); % Maroc N
\node[poppurple, font=\bfseries\small] at (2,8) {Espagne};
% Indépendants
\node[popgreen!50!black, font=\bfseries\small] at (1.5,3.5) {Libéria};
\node[popgreen!50!black, font=\bfseries\small] at (7,3.5) {Éthiopie};
% Flèches logiques
\draw[popblue, thick, ->] (1,2) -- (5,2) node[midway, above, font=\tiny] {Logique Est-Ouest (France)};
\draw[popred, thick, ->] (7,6) -- (7,1) node[midway, right, font=\tiny] {Logique Nord-Sud (R.-U.)};
% Échelle / Nord
\node[anchor=south east, font=\tiny] at (9.5,0.5) {\textbf{N}};
\draw[|-|] (8.5,0.5) -- (9.5,0.5) node[midway, below, font=\tiny] {2000 km};
\end{tikzpicture}
\legende{Carte schématique du partage de l'Afrique en 1914 (simplifiée pour l'exercice). Les couleurs correspondent aux puissances coloniales. Les flèches indiquent les axes d'expansion prioritaires.}
\end{popfigure}
\end{center}
\tcblower
\textbf{Solution détaillée :}

\textbf{1. Identification et présentation du document}
Ce document est une \textbf{carte schématique} représentant le \textbf{partage de l'Afrique en 1914}, à la veille de la Première Guerre mondiale. Elle illustre l'aboutissement de la \cle{conférence de Berlin (1884-1885)} et de la \cle{« course au partage »} (\emph{Scramble for Africa}). La légende (simulée ici par les couleurs) distingue sept puissances coloniales : la France (bleu), le Royaume-Uni (rouge), l'Allemagne (gris), la Belgique (or), le Portugal (vert), l'Italie (orange), l'Espagne (violet). Deux États échappent à la colonisation : le \cle{Libéria} (fondé d'anciens esclaves US) et l'\cle{Éthiopie} (victoire d'Adoua, 1896).

\textbf{2. Description structurée de la répartition territoriale}
\begin{itemize}
    \item \textbf{La France (1er empire colonial africain par la surface) :} Domine l'\cle{Afrique Occidentale Française (AOF)} (Sénégal à Niger) et l'\cle{Afrique Équatoriale Française (AEF)} (Gabon à Oubangui-Chari/Tchad), le \cle{Maghreb} (Algérie dès 1830, Tunisie 1881, Maroc 1912) et \cle{Madagascar} (1896). \textbf{Logique :} Expansion \textbf{Est-Ouest} (Dakar-Djibouti), visant la continuité territoriale (échec de Fachoda 1898).
    \item \textbf{Le Royaume-Uni (1er par la population/richesses) :} Contrôle l'\cle{Égypte} (1882) et le \cle{Soudan anglo-égyptien} (1898), le \cle{Nigeria}, l'\cle{Afrique du Sud} (Cap, Natal, Transvaal, Rhodésies), l'\cle{Afrique de l'Est} (Kenya, Ouganda). \textbf{Logique :} Projet \textbf{Nord-Sud} (\emph{Cape to Cairo}, Cecil Rhodes), contrôle de la route des Indes (Suez) et des ressources minières (or, diamants).
    \item \textbf{L'Allemagne (3e puissance, arrivée tardive) :} Quatre territoires disjoints : \cle{Togo}, \cle{Kamerun}, \cle{Afrique du Sud-Ouest}, \cle{Afrique Orientale Allemande}. \textbf{Logique :} \emph{Weltpolitik} (Bismarck puis Guillaume II), zones d'influence commerciale, morcellement imposé par les accords (Héligoland-Zanzibar 1890).
    \item \textbf{La Belgique :} \cle{État Indépendant du Congo} (propriété privée de Léopold II jusqu'en 1908, puis colonie belge). Cœur du bassin congolais.
    \item \textbf{Autres :} Portugal (Angola, Mozambique, Guinée - présence ancienne, effective tardive), Italie (Libye 1911, Erythrée, Somalie - défaite d'Adoua 1896), Espagne (Rio de Oro, Maroc nord, Guinée équatoriale).
\end{itemize}

\textbf{3. Explication des logiques d'expansion (France vs Royaume-Uni)}
\begin{itemize}
    \item \textbf{France :} \emph{Mission civilisatrice} (Jules Ferry), pression de l'opinion publique, intérêts économiques (matières premières : arachide, caoutchouc), stratégie militaire (marine, légion étrangère), rivalité avec l'Allemagne (revanche 1870) et le R.-U. (Fachoda). Méthode : \cle{conquête militaire} (colonnes Voulet-Chanoine, Gentil, Foureau-Lamy) + \cle{traité de protectorat} (souvent inégaux).
    \item \textbf{Royaume-Uni :} \emph{Indirect Rule} (Lugard), primauté du commerce (Royal Niger Company, BSAC), contrôle des points stratégiques (Suez, Cap, sources du Nil), régulation des conflits boers (Guerre du Transvaal 1899-1902). Moins de troupes métropolitaines, plus de troupes indigènes (King's African Rifles).
\end{itemize}

\textbf{4. Bilan et Conclusion}
En 1914, \textbf{90\% du continent est sous domination européenne}. Les frontières, tracées à la règle géométrique (méridiens/parallèles) ou suivant des cours d'eau, ignorent les réalités ethniques et linguistiques (\cle{balkanisation}). L'économie est réorientée vers l'exportation (infrastructures : chemins de fer « du port à la mine »). Cette situation contient les germes des nationalismes et des guerres de libération post-1945.

\textbf{Phrase de conclusion type Bac :} « La carte de 1914 fige le triomphe de l'impérialisme européen en Afrique, résultat de logiques concurrentielles (France/R.-U., Allemagne/autres) qui ont imposé un ordre territorial artificiel, fondement des États africains contemporains. »
\end{exoresolu}

\begin{methode}[Rédiger une dissertation structurée : Causes et facteurs de l'impérialisme]
\textbf{Objectif :} Rédiger et justifier une démarche argumentée (introduction, développement en parties équilibrées, conclusion) sur une question de synthèse.
\begin{enumerate}
    \item \textbf{Analyser le sujet :} Définir les termes clés (Impérialisme, Partage, 1870-1914), délimiter l'espace (Afrique) et le temps. Identifier la problématique : \emph{Pourquoi et comment l'Afrique est-elle passée de la traite à la colonisation formelle en quelques décennies ?}
    \item \textbf{Construire le plan (3 parties, 2-3 sous-parties chacune) :}
        \begin{itemize}
            \item \textbf{I. Les causes profondes (Le « Pourquoi ») :} Économiques (capitalisme monopoliste, Lénine/Hobson : surplus de capitaux, matières premières, marchés), Politiques (rivalités nationales, prestige, stratégie), Idéologiques (darwinisme social, mission civilisatrice, racisme scientifique).
            \item \textbf{II. Les facteurs déclenchants et le cadre légal (Le « Comment ») :} Explorations (Livingstone, Stanley, Brazza, Gallieni), Conférence de Berlin (1884-85 : principe d'effectivité, libre-échange, navigation), Traité de partage (bilatéraux).
            \item \textbf{III. Les modalités de la conquête (Le « Par quels moyens ») :} Conquête militaire (supériorité technique : fusil à répétition, mitrailleuse, quinine), Diplomatie inégale (traité de protectorat), Administration (Directe/Indirecte/Assimilation/Association).
        \end{itemize}
    \item \textbf{Rédiger l'introduction :} Accroche (chiffre : 10\% colonisé en 1870 $\to$ 90\% en 1914), Définitions, Problématique, Annonce du plan.
    \item \textbf{Développer :} Un argument = une idée force + \textbf{exemple précis daté} (ex: « Surplus de capitaux : investissements dans le chemin de fer Congo-Océan » ; « Mission civilisatrice : discours de Ferry 1885 »). Liaisons logiques (\emph{Par ailleurs, En outre, Cependant}).
    \item \textbf{Conclure :} Réponse synthétique à la problématique, Bilan nuancé (résistances, collaborations), Ouverture (héritage colonial, néocolonialisme).
\end{enumerate}
\textbf{Formule à retenir :} \textbf{1 Idée directrice + 1 Exemple daté/chiffré + 1 Lien avec la problématique} par paragraphe.
\end{methode}

\begin{exoresolu}[Dissertation : « Les causes de la conquête coloniale de l'Afrique (1870-1914) »]
\emph{Énoncé :} En vous appuyant sur des exemples précis, montrez que la conquête coloniale de l'Afrique résulte de la convergence de facteurs économiques, politiques et idéologiques.
\tcblower
\textbf{Solution détaillée :}

\textbf{INTRODUCTION}
\textbf{[Accroche]} En 1870, l'Europe ne contrôle que 10\% du littoral africain (comptoirs). En 1914, 90\% du continent est partagé entre sept puissances. Ce basculement radical en moins d'un demi-siècle est l'œuvre de l'\cle{impérialisme}.
\textbf{[Définitions]} L'impérialisme est la politique d'expansion d'une puissance étendant sa domination (politique, économique, culturelle) sur des territoires extérieurs. Le « partage de l'Afrique » (1881-1914) en est l'illustration majeure.
\textbf{[Problématique]} Comment expliquer cet emballement soudain ? La conquête répond-elle uniquement à des appétits économiques ou résulte-t-elle d'un faisceau de causes conjuguées ?
\textbf{[Annonce du plan]} Nous analyserons d'abord les causes structurelles économiques (I), puis les déterminants politiques et stratégiques (II), enfin les justifications idéologiques et le cadre juridique qui ont légitimé l'entreprise (III).

\textbf{DÉVELOPPEMENT}

\textbf{I. L'impératif économique : le moteur principal du capitalisme monopoliste}
\textbf{A. La recherche de débouchés et de matières premières.} La seconde révolution industrielle (acier, chimie, électricité) crée une demande insatiable : caoutchouc (Congo, AEF), huile de palme (Nigeria, Côte d'Ivoire), coton (Égypte, Soudan), minerais (or/diamants Af. du Sud, cuivre Katanga). \textbf{Exemple :} La \cle{Royal Niger Company} (R.-U.) et la \cle{Compagnie du Congo} (France) monopolisent le commerce.
\textbf{B. L'exportation des capitaux (Thèse de Lénine/Hobson).} La saturation des marchés métropolitains pousse à investir à l'étranger (chemins de fer, mines). \textbf{Exemple :} Le chemin de fer \cle{Cap-Caire} (Rhodes, inachevé) ou \cle{Dakar-Niger} (France) financés par la Bourse de Paris/Londres.
\textbf{C. La main-d'œuvre bon marché.} Le code de l'indigénat et le travail forcé (corvées, impôt de capitation) assurent la rentabilité.

\textbf{II. Les enjeux politiques et stratégiques : la rivalité des puissances}
\textbf{A. L'équilibre européen et la \cle{Weltpolitik}.} L'Allemagne unifiée (1871) réclame sa « place au soleil » (Guillaume II), bousculant le duo France/R.-U. La colonisation devient un indicateur de puissance (\emph{Great Game}).
\textbf{B. La stratégie militaire et les points de passage.} Contrôle du \cle{canal de Suez} (1869, R.-U. 1882), du \cle{détroit de Gibraltar}, de la \cle{route du Cap}. \textbf{Exemple :} Fachoda (1898) : affrontement franco-britannique pour le contrôle du Nil.
\textbf{C. La diplomatie préventive.} La \cle{Conférence de Berlin (1884-1885)} (Bismarck) pose les règles : \cle{principe d'effectivité} (occupation effective > simple découverte), libre navigation (Congo, Niger), abolition de la traite. Elle institutionnalise le partage sans guerre européenne directe.

\textbf{III. L'idéologie et le cadre juridique : légitimer l'inacceptable}
\textbf{A. Le darwinisme social et le racisme scientifique.} Application erronée de la sélection naturelle aux sociétés : « races supérieures » (blanches) destinées à dominer « races inférieures » (noires). \textbf{Exemple :} Gobineau, \emph{Essai sur l'inégalité des races humaines} (1853-55).
\textbf{B. La « Mission civilisatrice » (France) / « White Man's Burden » (Kipling, R.-U.).} Devoir moral d'apporter la paix, la médecine, l'école, le christianisme. \textbf{Exemple :} Discours de \cle{Jules Ferry} (Chambre des députés, 1885) : « Les races supérieures ont un devoir vis-à-vis des races inférieures. »
\textbf{C. Le rôle des explorateurs et missionnaires.} « Avant-garde » de la colonisation (Livingstone, Stanley, Brazza). Ils signent des traités (souvent incompris par les chefs africains) qui servent de titres juridiques.

\textbf{CONCLUSION}
\textbf{[Bilan]} La conquête africaine n'est pas monocausale. Elle naît de la \textbf{convergence} : le capitalisme fournit le \emph{moteur} (profit), les États-nations fournissent la \emph{volonté politique} (puissance, sécurité), l'idéologie fournit le \emph{verni moral} (civilisation).
\textbf{[Nuance]} Cette conquête ne s'est pas faite sans heurts : elle a provoqué des \cle{résistances africaines} (Samory, Behanzin, Menelik) et des \cle{collaborations} complexes.
\textbf{[Ouverture]} Les frontières issues de Berlin (1885) structurent encore l'Afrique actuelle, source de conflits frontaliers et de défis d'intégration (UA, ZLECAf).
\end{exoresolu}

\begin{methode}[Comparer les modes d'administration coloniale : Assimilation, Association, Indirect Rule]
\textbf{Objectif :} Appliquer les notions de l'unité à une situation concrète : différencier les modèles français et britannique pour un tableau comparatif ou une question de cours.
\begin{enumerate}
    \item \textbf{Identifier le modèle français :} \cle{Assimilation} (Théorie : l'indigène devient citoyen français s'il adopte la culture/lois françaises. Pratique : \cle{Quatre Communes} du Sénégal - Saint-Louis, Gorée, Dakar, Rufisque - élues députés, ex: Blaise Diagne 1914). \cle{Association} (Pratique dominante après 1900 : respecter les « coutumes », gouverner par les chefs traditionnels \emph{instrumentalisés}, Code de l'indigénat, travail forcé, pas de citoyenneté). Distinction : \cle{Citoyens} (droits pleins) vs \cle{Sujets} (devoirs, pas de droits politiques).
    \item \textbf{Identifier le modèle britannique :} \cle{Indirect Rule} (Lugard, Nigeria du Nord). Théorie : « Dual Mandate » (développement pour l'Afrique et pour le monde). Pratique : Gouvernement \emph{à travers} les \cle{chefs traditionnels} (Emirs, Obas) maintenus comme auxiliaires de l'administration (Native Authorities), tribunaux coutumiers (Native Courts), trésorerie locale. Moins de fonctionnaires britanniques, pas d'assimilation culturelle.
    \item \textbf{Comparer sur critères fixes :} Statut juridique (Citoyen/Sujet vs Sujet protégé), Rôle des chefs (Remplacés/Intégrés), Droit (Code civil/Coutumier), École (Française/Langue vernaculaire puis anglais), Fiscalité (Impôt unique/Impôts divers).
    \item \textbf{Nuancer :} La France pratique l'Association dans les zones « non assimilées » (majorité). Le R.-U. impose le \cle{Direct Rule} au Kenya (colons blancs) ou en Rhodésie du Sud. L'Allemagne pratique un \cle{Direct Rule} militarisé et brutal (génocide Herero/Nama 1904-1908).
    \item \textbf{Conclure :} Deux logiques : \emph{Uniformisation centralisée} (France) vs \emph{Gestion décentralisée à moindre coût} (R.-U.). Conséquences post-indépendance : centralisme jacobin vs fédéralisme/chefferies persistantes.
\end{enumerate}
\textbf{Reflexe Bac :} Ne pas opposer bêtement « Assimilation = France » et « Indirect Rule = R.-U. ». Préciser : « La France \emph{théorise} l'assimilation mais \emph{pratique} massivement l'association. Le R.-U. \emph{théorise} l'Indirect Rule mais l'adapte (Direct Rule pour colons). »
\end{methode}

\begin{exoresolu}[Analyse comparative : Administration française vs britannique en AOF et au Nigeria]
\emph{Énoncé :} À l'aide d'un tableau comparatif, montrez les différences et similitudes entre l'administration coloniale française en AOF et l'administration britannique au Nigeria (début XXe s.).
\tcblower
\textbf{Solution détaillée :}

\begin{center}
\begin{tabularx}{\linewidth}{|l|X|X|}
\hline
\rowcolor{popblueL}
\textbf{Critère de comparaison} & \textbf{France (AOF) : « Association » (Pratique dominante)} & \textbf{Royaume-Uni (Nigeria) : « Indirect Rule » (Lugard)} \\
\hline
\textbf{Fondement théorique} & \emph{Mission civilisatrice} (Assimilation idéale) mais pragmatisme : respect apparent des « coutumes » pour gouverner à moindre frais. & \emph{Dual Mandate} (Lugard) : Développer l'Afrique pour l'Afrique et pour l'humanité. Gouverner \emph{par} les institutions indigènes. \\
\hline
\textbf{Statut juridique des populations} & \textbf{Dualité stricte :} \cle{Citoyens} (Originaires des 4 Communes : droits civils/politiques, Code civil) vs \cle{Sujets} (Immense majorité : Code de l'indigénat, justice d'exception, travail forcé, pas de droits politiques). & \textbf{Unité de statut :} Tous \cle{Sujets britanniques protégés} (British Protected Persons). Pas de citoyenneté. Différenciation par la « tribu »/ethnie via les \cle{Native Authorities}. \\
\hline
\textbf{Rôle des chefs traditionnels} & \cle{Chefs de canton/cercle} : Nommé\emph{s} et révoqué\emph{s} par l'administrateur français. Agents subalternes de l'État (perception impôt, recrutement main-d'œuvre, police). Perte de légitimité traditionnelle. & \cle{Emirs (Nord) / Obas (Sud-Ouest) / Warrant Chiefs (Sud-Est)} : Maintien de la hiérarchie traditionnelle. \textbf{Native Authorities} : Pouvoir exécutif, judiciaire (Native Courts), financier (Native Treasuries). Légitimité traditionnelle \emph{instrumentalisée}. \\
\hline
\textbf{Organisation administrative} & \textbf{Pyramide centralisée :} Gouverneur Général (Dakar) $\to$ Lieutenants Gouverneurs (Cercles) $\to$ Commandants de Cercle $\to$ Chefs de Canton. Corps des \cle{Administrateurs de la France d'Outre-Mer}. & \textbf{Fédéralisme déconcentré :} Gouverneur (Lagos) $\to$ Lieutenant-Gouverneurs (Nord, Sud, Lagos) $\to$ \cle{Residents / District Officers} (britanniques, peu nombreux) $\to$ \cle{Native Authorities}. \\
\hline
\textbf{Justice et Droit} & Double système : Tribunaux français (citoyens) / \cle{Justice indigène} (sujets : administrateur juge seul, peine de prison, travail forcé). Code de l'indigénat (1887, 1912). & \cle{Tribunaux coutumiers (Native Courts)} : Jugés par les chefs selon la coutume (reconnaissance légale). Cour suprême britannique pour appel. Pas de code pénal unique imposé. \\
\hline
\textbf{Éducation et Culture} & \textbf{École laïque, française, centralisée.} But : créer une élite assimilée (évolués). Missionnaires tolérés mais contrôlés. Langue : Français seul. & \textbf{École missionnaire subventionnée.} But : formation pratique, alphabétisation en langues vernaculaires puis anglais. \cle{Indirect Rule} préserve cultures/langues locales. \\
\hline
\textbf{Économie / Main d'œuvre} & \cle{Travail forcé} (prestations, portage, travaux publics) légalisé (décret 1912). Impôt de capitation (contrainte au salariat). Grandes concessions (CFAO, etc.). & \cle{Travail forcé} limité (abolition officielle 1909 mais persiste sous forme de corvées communautaires). Impôt direct (\cle{Hut Tax}, \cle{Poll Tax}) monétarise l'économie paysanne. Compagnie : \cle{Royal Niger Company} puis État. \\
\hline
\end{tabularx}
\end{center}

\textbf{Synthèse argumentée (Réponse à la problématique) :}
Bien que toutes deux soient des \cle{dominations coloniales} (exploitation économique, racisme structurel, violence), les \textbf{logiques administratives divergent} :
\begin{itemize}
    \item \textbf{France (Centralisation jacobine) :} L'État colonial est \emph{omniprésent} et \emph{uniformisateur}. Il casse les chefferies traditionnelles pour les remplacer par des auxiliaires nommés. L'objectif final (théorique) est l'intégration à la République (Assimilation), créant une fracture juridique violente (Citoyens/Sujets).
    \item \textbf{R.-U. (Décentralisation pragmatique) :} L'État colonial est \emph{léger} et \emph{médiatisé} par les chefferies. Il « gèle » les structures traditionnelles (figement des coutumes, invention de la « tribu »). Pas de perspective d'intégration citoyenne.
\end{itemize}
\textbf{Conséquence historique :} À l'indépendance, les États francophones héritent d'un \textbf{État fort, centralisé, unitaire} (souvent autoritaire). Les États anglophones héritent d'un \textbf{État fédéral/fragmenté}, avec des chefferies rivales du pouvoir central (ex: Nigeria, Ghana).
\end{exoresolu}

\begin{methode}[Commenter un document sur les résistances africaines (Typologie, Acteurs, Échec/Héritage)]
\textbf{Objectif :} Rédiger et justifier une démarche d'analyse de document (texte, image, discours) sur les résistances.
\begin{enumerate}
    \item \textbf{Identification :} Nature (discours, lettre, rapport militaire, oralité transcrite), Auteur (Chef résistant, officier colonial, missionnaire), Date, Contexte (conquête, pacification, rébellion), Destinataire.
    \item \textbf{Classification de la résistance (Typologie) :}
        \begin{itemize}
            \item \textbf{Primaire / Immédiate :} Au contact (ex: \cle{Lat Dior} au Cayor, \cle{Samory} début, \cle{Behanzin} Dahomey).
            \item \textbf{Secondaire / Différée :} Après installation administration (ex: \cle{Baoulé} 1910, \cle{Kongo-Wara} 1928-31 - hors programme 1re mais bon à savoir pour contexte).
            \item \textbf{Religieuse / Mahdiste :} \cle{Mohamed Ahmed} (Soudan), \cle{Cheikh Amadou Bamba} (Mouridisme, résistance spirituelle), \cle{Simon Kimbangu} (plus tard).
            \item \textbf{Diplomatique / Juridique :} \cle{Menelik II} (Éthiopie, traité de Wuchalé, victoire d'Adoua 1896), \cle{Ranavalona III} (pétitions).
        \end{itemize}
    \item \textbf{Analyse du contenu :} Thèses défendues (souveraineté, islam, refus de l'impôt/travail forcé), Arguments (légitimité ancestrale, droit international, foi), Moyens (guérilla, armée moderne, diplomatie, refus de collaboration).
    \item \textbf{Mise en perspective (Connaissances) :} Pourquoi l'échec militaire ? (Supériorité technique : \cle{mitrailleuse}, \cle{quinine}, \cle{artillerie}, divisions africaines, traite négrière affaiblie). Pourquoi l'héritage ? (Symboles nationaux, pères de l'indépendance).
    \item \textbf{Conclusion :} Répondre à la question (ex: « En quoi ce document illustre-t-il la diversité des réactions africaines ? »).
\end{enumerate}
\textbf{Attention :} Ne pas confondre \emph{résistance} (lutte armée/politique contre la conquête) et \emph{révolte} (souvent locale, fiscale, post-conquête). Ne pas présenter les Africains comme passifs.
\end{methode}

\begin{exoresolu}[Commentaire de texte : L'appel à la résistance de Samory Touré (vers 1890)]
\emph{Énoncé :} Vous commenterez le texte suivant (extrait fictif représentatif des propos rapportés de Samory) : « Les Blancs viennent non pour commercer, mais pour voler notre terre. Ils ont des fusils qui tirent vite, mais nous avons le courage de nos ancêtres et la foi en un seul Dieu. Je ne signerai pas de papier qui fait de moi un esclave. Nous combattrons jusqu'à ce qu'Allah décide. »
\tcblower
\textbf{Solution détaillée :}

\textbf{1. Identification et Contexte}
\textbf{Nature :} Extrait de discours / propos rapportés (source orale/écrite coloniale).
\textbf{Auteur :} \cle{Samory Touré} (v. 1830-1900), almamy de l'\cle{Empire wassoulou} (Mandinka), résistant majeur à la conquête française en Afrique de l'Ouest (Guinée, Mali, Côte d'Ivoire actuels).
\textbf{Date :} Vers 1890-1898 (période de la « guerre de Samory » contre les colonnes françaises : Combes, Gouraud, Gentil).
\textbf{Contexte :} Expansion française en AOF (Colonnes Voulet-Chanoine, Foureau-Lamy, Gentil). Samory a déjà signé des traités (Bissandougou 1887, Kénédougou 1889) qu'il considère violés. Il déplace son empire vers l'Est (Kong, Bondoukou) pour échapper à l'étau.
\textbf{Problématique implicite :} Comment ce texte illustre-t-il la nature politique, militaire et spirituelle de la résistance africaine face à la supériorité technique européenne ?

\textbf{2. Analyse structurée du document}
\textbf{A. Une analyse lucide de l'intention coloniale (Versets 1-2) :} « Les Blancs viennent non pour commercer, mais pour voler notre terre. »
\begin{itemize}
    \item \textbf{Justesse politique :} Samory perçoit la rupture avec l'époque de la \cle{traite} et du \cle{commerce légitime} (arachide, or). La \cle{Conférence de Berlin} a transformé le commerce en \cle{conquête territoriale} (principe d'effectivité).
    \item \textbf{Refus du leurre juridique :} « Je ne signerai pas de papier qui fait de moi un esclave. » Référence aux \cle{traité de protectorat} inégaux (souvent traduits oralement de manière trompeuse, signature par marque/croix). Samory, lettré (arabe), en connaît la portée juridique : perte de souveraineté.
\end{itemize}

\textbf{B. La conscience de l'asymétrie militaire et la réponse stratégique (Verset 2) :} « Ils ont des fusils qui tirent vite, mais nous avons le courage... »
\begin{itemize}
    \item \textbf{Reconnaissance technique :} « Fusils qui tirent vite » = \cle{Fusil Gras Mle 1874} (monocoup mais rechargement rapide) et surtout \cle{Mitrailleuse Reffye / Hotchkiss} (colonnes françaises). Samory a tenté de s'adapter : achat de fusils \cle{Chassepot/Remington} chez les Anglais (Sierra Leone/Libéria), fabrication locale de poudre (\cle{ateliers de Bissandougou/Kankan}), tactique de \cle{guérilla} et \cle{terre brûlée} (déménagement de l'empire).
    \item \textbf{Facteur moral :} « Courage de nos ancêtres » = mobilisation de l'\cle{honorabilité guerrière mandinka} (sofa, tirailleurs d'élite).
\end{itemize}

\textbf{C. La dimension spirituelle comme ciment identitaire (Verset 3) :} « ...la foi en un seul Dieu... jusqu'à ce qu'Allah décide. »
\begin{itemize}
    \item \textbf{Islam comme idéologie de résistance :} Samory se présente comme \cle{Almamy} (Commandeur des croyants), chef d'un \cle{État théocratique} (Wassoulou). L'islam structure l'armée (discipline, unité), la justice (charia), la diplomatie (alliance avec Tieba de Sikasso, Ahmadu Seku).
    \item \textbf{Fatalisme actif :} « Jusqu'à ce qu'Allah décide » $\neq$ passivité. C'est l'acceptation du martyre possible (\cle{chahada}) qui libère de la peur de la mort, face à un ennemi technologiquement supérieur.
\end{itemize}

\textbf{3. Mise en perspective historique (Connaissances)}
\begin{itemize}
    \item \textbf{Typologie :} Résistance \textbf{primaire, étatique, islamique, diplomatique et militaire}. Durée exceptionnelle (1881-1898, 17 ans).
    \item \textbf{Moyens modernes :} Armée permanente (\cle{sofa} $\approx$ 30-40 000 hommes), administration centrale, ateliers d'armement, diplomatie (jouer Franco-Anglais).
    \item \textbf{Fin :} Capture à \cle{Guélémou} (1898) par le capitaine \cle{Gourais}. Déportation au Gabon (Ndjolé) puis Tombouctou. Mort 1900.
    \item \textbf{Héritage :} Figure panafricaine majeure. Cité par \cle{Nkrumah}, \cle{Sékou Touré} (parenté revendiquée), \cle{Mandela}. Symbole du \cle{refus de la soumission}.
\end{itemize}

\textbf{4. Conclusion}
Ce texte résume la \textbf{tragédie lucide} de la résistance africaine : une conscience politique aiguë de la nature prédatrice du colonialisme, une adaptation technique réelle mais insuffisante face à la révolution industrielle militaire, et une mobilisation identitaire (islam, tradition) qui permet une résistance d'une durée inédite. Samory incarne la \cle{résistance d'État} la plus aboutie de l'Afrique précoloniale.
\end{exoresolu}

\begin{methode}[Analyser les collaborations africaines : Motivations, Formes, Ambiguïtés]
\textbf{Objectif :} Appliquer les notions pour expliquer pourquoi et comment des Africains ont collaboré, sans anachronisme ni jugement moralisateur.
\begin{enumerate}
    \item \textbf{Définir la collaboration :} Acceptation (contrainte ou choisie) de la présence coloniale, participation à l'administration, à l'économie ou à l'armée coloniale. \cle{Distinction :} Collaboration \emph{de contrainte} (survie, otages) vs \emph{de choix} (intérêt dynastique, moderne, religieux).
    \item \textbf{Identifier les profils types (Acteurs) :}
        \begin{itemize}
            \item \textbf{Chefs traditionnels ralliés :} \cle{Lat Dior} (résistance puis ralliement temporaire, Cayor), \cle{Tofâ} (Dahomey, frère de Béhanzin, installé par Français), \cle{Chefs de canton AOF} (nommés).
            \item \textbf{Évolués / Lettrés / « Assimilés » :} \cle{Blaise Diagne} (Député Sénégal, recrutement 1918), \cle{Galandou Diouf}, \cle{William Ponty} (élèves école normale). Motivations : citoyenneté, modernité, abolition code indigénat.
            \item \textbf{Militaires / Tirailleurs :} \cle{Tirailleurs sénégalais} (création 1857, Faidherbe). Motivations : solde, statut social, évasion condition paysanne.
            \item \textbf{Intermédiaires économiques :} \cle{Commerçants côtiers} (Lagos, Freetown, Grand-Bassam, Douala), \cle{Courtiers} (Dioula, Haoussa).
            \item \textbf{Religieux :} Certains marabouts (Tidjanes, Qadiri) acceptant le protectorat pour protéger leurs zawiyas (ex: \cle{Cheikh Sidia Baba} Mauritanie vs \cle{Ma Ba Diakhou} résistance).
        \end{itemize}
    \item \textbf{Analyser les motivations (Grille de lecture) :}
        \begin{itemize}
            \item \textbf{Politique :} Sauver son trône/dynastie face à un rival interne soutenu par l'ennemi (ex: \cle{Maba Diakhou} vs Lat Dior ; \cle{Tofâ} vs Béhanzin).
            \item \textbf{Économique :} Monopole du commerce, accès aux armes, suppression de l'impôt/travail forcé pour soi.
            \item \textbf{Sociale/Identitaire :} Accès à l'école, à la médecine, statut de « citoyen » ou « notable », conversion au christianisme (missionnaires).
        \end{itemize}
    \item \textbf{Évaluer le bilan :} Court terme : pouvoir partagé, protection. Long terme : \cle{perte de légitimité traditionnelle} (chef devient fonctionnaire), \cle{dépendance} vis-à-vis de l'administration, \cle{instrumentalisation} (recrutement forcé 14-18), \cle{naissance du nationalisme} par les évolués déçus (Diagne $\to$ Lamine Guèye $\to$ Senghor).
    \item \textbf{Rédiger :} Utiliser des termes précis : \cle{Alliance}, \cle{Ralliement}, \cle{Collaborationnisme}, \cle{Intermédiation}. Éviter « traitre » (subjectif) ou « héros » (réducteur).
\end{enumerate}
\end{methode}

\begin{exoresolu}[Étude de cas : La collaboration au Dahomey (Béhanzin vs Tofâ / Agoli-Agbo)]
\emph{Énoncé :} En vous appuyant sur l'exemple du Dahomey (1892-1900), expliquez les logiques de la collaboration africaine face à la conquête française.
\tcblower
\textbf{Solution détaillée :}

\textbf{1. Contexte : La conquête du Dahomey (1892-1894)}
Le royaume du \cle{Dahomey} (Roi \cle{Béhanzin}, 1889-1894) est une puissance militaire centralisée (amazones, armée permanente), négrière puis palmiériste. La France (Général \cle{Dodds}) veut contrôler le littoral (Cotonou, Porto-Novo) et l'intérieur (chemin de fer). Guerre de 1892-94 : \cle{Bataille de Dogba}, \cle{Poguessa}, \cle{Atchoukpa}, prise d'\cle{Abomey} (nov. 1892). Béhanzin refuse le protectorat, mène guérilla, se rend (1894), déporté Martinique puis Algérie.

\textbf{2. La mise en place de la collaboration : Le « Roi » français}
\begin{itemize}
    \item \textbf{Choix français :} Refus de la monarchie forte (Béhanzin). Installation d'un roi « docile » : \cle{Tofâ} (frère cadet de Béhanzin, converti au christianisme, pro-français), puis \cle{Agoli-Agbo} (1894-1900, dernier roi, déposé pour « indiscipline »).
    \item \textbf{Statut :} \cle{Chefs de province} / \cle{Chefs de canton} nommés par l'administrateur. Le roi devient un \cle{fonctionnaire colonial} (solde, prérogatives réduites : justice coutumière surveillée, perception impôt).
\end{itemize}

\textbf{3. Analyse des logiques de collaboration (Grille multi-acteurs)}

\begin{center}
\begin{tabularx}{\linewidth}{|l|X|X|}
\hline
\rowcolor{popblueL}
\textbf{Acteur / Groupe} & \textbf{Motivation principale} & \textbf{Forme de collaboration / Conséquence} \\
\hline
\textbf{Tofâ / Agoli-Agbo (Dynastie)} & \textbf{Survie dynastique \& Pouvoir personnel.} Sauver la monarchie face à l'annexion pure. Rivalité successorale (Béhanzin vs frères). & Acceptation du \cle{protectorat} (traité 1894). Perte de l'armée (amazones dissoutes), de la justice souveraine, du contrôle du commerce (palmistes). Déposition finale (1900) : la collaboration n'a pas sauvé l'institution royale. \\
\hline
\textbf{Princes / Chefs de province (ex: Goho, Yémadjé)} & \textbf{Intérêt foncier \& statut.} Conserver leurs terres, esclaves, prérogatives locales. & Ralliement individuel. Intégration dans l'administration de \cle{cercle} comme chefs de canton. Perte d'autonomie militaire. \\
\hline
\textbf{Commerçants côtiers (Brésiléens/Agoudas, Ouidah, Porto-Novo)} & \textbf{Sécurisation du commerce.} Fin des guerres, des raids, instabilité monétaire. Accès au marché français. & Appui logistique (porteurs, renseignements) à l'expédition Dodds. Émergence d'une bourgeoisie compradore (figure type ultérieure : \cle{Félix Houphouët-Boigny} en Côte d'Ivoire). \\
\hline
\textbf{Population / Anciens esclaves (Donkomè)} & \textbf{Libération / Mobilité.} Fin de la traite interne, des sacrifices humains (coutumes annuelles), recrutement forcé. & Fuite vers le protectorat français (Cotonou, Porto-Novo) = \cle{« vote par les pieds »}. Main-d'œuvre pour colons/chemins de fer. \\
\hline
\textbf{Évolués / Clercs (Premiers lettrés)} & \textbf{Ascension sociale.} École des frères (Ouidah), administration. & Premiers interprètes, commis, instituteurs. Base future de l'élite politique (Sourou-Migan Apithy, Hubert Maga). \\
\hline
\end{tabularx}
\end{center}

\textbf{4. L'ambiguïté fondamentale : Collaboration vs Résistance (Le cas Behanzin vs Tofâ)}
\begin{itemize}
    \item \textbf{Béhanzin} incarne la \cle{résistance souveraine} : refus de signer, guerre, exil. Il devient \textbf{héros national} (Bénin actuel : statue, place, billet).
    \item \textbf{Tofâ/Agoli-Agbo} incarnent la \cle{realpolitik} : sauver le peuple et le trône par le compromis. Ils sont \textbf{oubliés ou critiqués} par l'historiographie nationale (collabos), mais ont permis la survie physique de la lignée et de la cour.
    \item \textbf{Nuance :} Béhanzin lui-même a négocié (traité 1890 avec Bayol) avant de le rompre. La frontière est poreuse.
\end{itemize}

\textbf{5. Conclusion}
Au Dahomey, la collaboration a été \textbf{imposée par la défaite militaire»} (supériorité feu français, amazones décimées) et \textbf{choisie par des élites} (dynastiques, commerciales) pour préserver des intérêts vitaux. Elle a abouti à la \textbf{colonisation « douce »} (protectorat puis colonie 1900) mais a vidé la royauté de sa substance politique, la transformant en rouage de l'\cle{Indirect Rule à la française} (Association). L'héritage est double : une administration locale structurée par les chefferies, et un nationalisme postérieur (années 1940-50) né de la frustration des \cle{évolués} (Diagne, Apithy) exclus du pouvoir réel malgré leur collaboration.
\end{exoresolu}

\begin{methode}[Réaliser une synthèse chronologique et thématique : L'Afrique en 1914 (Bilan du partage)]
\textbf{Objectif :} Rédiger une conclusion de cours / réponse à une question de synthèse (type « Bilan » ou « Situation en 1914 »).
\begin{enumerate}
    \item \textbf{Frise chronologique repères (À mémoriser) :}
        \begin{itemize}
            \item 1876 : Conférence de Bruxelles (Association internationale africaine, Léopold II).
            \item 1880-1885 : Course aux traités (Brazza, Stanley, Peters, Gallieni).
            \item \textbf{1884-1885 : Conférence de Berlin (Acte général).} Règles du jeu.
            \item 1885-1890 : Partage « sur le papier » (traité bilatéraux : France/R.-U., Allemagne/R.-U., etc.).
            \item 1890-1898 : Conquête effective (pacification) : Dahomey, Soudan français, Madagascar, Soudan anglo-égyptien, Matabéléland, Afrique Orientale Allemande.
            \item \textbf{1898 : Fachoda (crise franco-britannique).} Apogée de la rivalité.
            \item 1904-1911 : Accords définitifs (Entente cordiale 1904, Agadir 1911, Maroc).
            \item \textbf{1914 : Afrique partagée (sauf Libéria, Éthiopie).}
        \end{itemize}
    \item \textbf{Structurer le bilan en 3 dimensions :}
        \begin{itemize}
            \item \textbf{Politique :} 7 puissances. Frontières artificielles (lignes droites, fleuves). 2 États indépendants (Libéria, Éthiopie). Administration : Directe/Indirecte/Association. Code de l'indigénat / Travail forcé.
            \item \textbf{Économique :} \cle{Économie d'enclave} (Port-Mine/Plantation-Port). Infrastructures : Chemins de fer « pendulaires » (Dakar-Niger, Abidjan-Ouaga, Cap-Caire inachevé). Monocultures d'exportation (Arachide, Cacao, Café, Coton, Caoutchouc, Minerais). Intégration forcée au capitalisme mondial (prix imposés).
            \item \textbf{Sociale / Culturelle :} École coloniale (élite francophone/anglophone). Médecine (lutte épidémies, mais inégalitaire). Démographie : début transition (paix coloniale, vaccination). \cle{Résistances} (armées, culturelles, religieuses) persistantes. Naissance du \cle{nationalisme} (Diagne, Casely-Hayford, ANC 1912).
        \end{itemize}
    \item \textbf{Rédiger la synthèse :} Introduction (chiffres clés), 3 paragraphes thématiques, Conclusion (héritage contradictoire : État-nation artificiel vs unité panafricaine).
\end{enumerate}
\textbf{Formule Bac :} « En 1914, l'Afrique est un \textbf{continent colonisé} (90\%), \textbf{morcelé} (frontières de Berlin), \textbf{exploité} (économie extravertie), mais \textbf{en mouvement} (résistances, naissances nationales). »
\end{methode}

\begin{exoresolu}[Question de synthèse : « L'Afrique en 1914 : un continent transformé »]
\emph{Énoncé :} Faites le bilan de la situation de l'Afrique à la veille de la Première Guerre mondiale. Montrez en quoi la colonisation a profondément transformé les structures politiques, économiques et sociales du continent.
\tcblower
\textbf{Solution détaillée :}

\textbf{INTRODUCTION}
En 1914, l'Afrique a cessé d'être la « terre incognite » des explorateurs pour devenir un \textbf{échiquier colonial} figé par les traités européens. Seuls le \cle{Libéria} (indépendance 1847, protectorat US de fait) et l'\cle{Éthiopie} (victoire d'\cle{Adoua} 1896 sur l'Italie) échappent au partage. Cette situation résulte de 30 ans de conquête (1885-1914) qui ont restructuré le continent en profondeur.

\textbf{DÉVELOPPEMENT}

\textbf{I. Une transformation politique radicale : l'invention de l'État colonial}
\begin{itemize}
    \item \textbf{Le morcellement territorial :} Le continent est découpé en \textbf{40 entités coloniales} (contre ~10 000 entités politiques précoloniales). Les frontières (80\% lignes géométriques : méridiens/parallèles) ignorent les peuples (\cle{Partition des Somaliens, Yorubas, Haoussas, Bakongo, Luo}). \textbf{Exemple :} Le peuple \cle{Somali} partagé entre R.-U. (Somaliland), France (Côte française des Somalis), Italie (Somalie), Éthiopie (Ogaden).
    \item \textbf{L'instauration de l'administration coloniale :} Remplacement de la souveraineté africaine par la \cle{souveraineté déléguée} (Gouverneurs responsables devant la Métropole). Deux modèles dominants : \cle{Administration directe/Association} (France, Belgique, Allemagne, Portugal) vs \cle{Indirect Rule} (R.-U.). Partout : \cle{monopole de la violence légitime} (troupes coloniales, police indigène), \cle{justice d'exception} (Code de l'indigénat, travail forcé légal).
    \item \textbf{La fin de la traite et de l'esclavage interne :} \cle{Acte de Berlin (1885), Acte de Bruxelles (1890)}. La colonisation impose l'abolition légale (mais le travail forcé la remplace souvent).
\end{itemize}

\textbf{II. Une restructuration économique violente : l'économie extravertie}
\begin{itemize}
    \item \textbf{Spécialisation contrainte :} Chaque colonie assignée à une fonction : \cle{Grenier} (arachide Sénégal, cacao Côte d'Ivoire/Ghana, coton Égypte/Soudan), \cle{Mine} (Or/Diamants Af. Sud, Cuivre Katanga/Zambie, Étain Nigeria), \cle{Plantation} (Caoutchouc Congo/LEO, Palme Dahomey/Cameroun, Sisal Tanzanie/Kenya).
    \item \textbf{Infrastructures « pendulaires » :} Chemins de fer reliant \textbf{l'intérieur (production) au port (exportation)}. Pas de liaison inter-africaine (Dakar-Niger ne va pas à Abidjan). \textbf{Exemple :} Chemin de fer \cle{Congo-Océan} (1921-34, coût humain énorme), \cle{Tazara} (plus tard).
    \item \textbf{Monétarisation et fiscalisation :} Impôt de capitation (\cle{Hut Tax}, \cle{Capitation}) payable en numéraire $\to$ contrainte à la vente de force de travail ou de cultures d'exportation. Insertion forcée dans le \cle{marché capitaliste mondial} (prix fixés à Londres/Paris/Anvers).
\end{itemize}

\textbf{III. Des mutations sociales et culturelles profondes : l'émergence de nouveaux acteurs}
\begin{itemize}
    \item \textbf{Démographie :} Fin des guerres inter-africaines et razzias $\to$ début de croissance démographique (paix coloniale, vaccination variole, sommeil). Mais travail forcé, migrations forcées (plantations, mines) et famines (réquisitions 14-18) freinent.
    \item \textbf{L'école coloniale :} Outil de formation d'auxiliaires (clercs, infirmiers, instituteurs, interprètes). \cle{École des otages} (Saint-Louis, 1855), \cle{École William Ponty} (Gorée, 1903), \cle{Achimoto College} (Gold Coast), \cle{Lovedale} (Af. Sud). Création d'une \cle{élite lettrée} (évolués) bilingue, médiatrice, mais frustrée (plafond de verre, code de l'indigénat).
    \item \textbf{Nouvelles identités et premiers nationalismes :} \cle{Christianisme} (églises indépendantes : Kimbanguisme 1921, Églises éthiopiennes), \cle{Islam} (réformisme, confréries). Premières organisations politiques : \cle{ANC} (Af. Sud, 1912), \cle{NCBWA} (Af. Occ. Britannique, 1920), \cle{Blaise Diagne} élu député (1914). \textbf{Paradoxe :} La colonisation crée les outils (langue, administration, territoire) de sa propre contestation future.
\end{itemize}

\textbf{CONCLUSION}
\textbf{[Bilan]} En 1914, l'Afrique est \textbf{politiquement morcelée}, \textbf{économiquement dépendante} (extraversion), \textbf{socialement bouleversée} (urbanisation naissante, nouvelles élites, migrations).
\textbf{[Héritage]} Les frontières de 1885 deviennent les frontières des États indépendants (1960). L'économie d'enclave persiste (néocolonialisme). Les élites formées à l'école coloniale (Senghor, Houphouët, Nkrumah, Kenyatta, Nyerere, Mandela) mèneront les luttes de libération.
\textbf{[Ouverture]} La \cle{Première Guerre mondiale} (1914-1918) va agir comme un révélateur : participation massive des tirailleurs (200 000+), promesses non tenues, prise de conscience du droit des peuples (Wilson, Lénine), accélération du nationalisme.
\end{exoresolu}

\begin{methode}[Exploiter une source statistique / graphique (Évolution du commerce, chemins de fer, population)]
\textbf{Objectif :} Appliquer les notions à des données chiffrées (tableau, courbe, diagramme) pour étayer une démonstration historique.
\begin{enumerate}
    \item \textbf{Lire le document :} Titre, Source, Unités (tonnes, francs, km, \%), Dates, Échelle.
    \item \textbf{Calculer les indicateurs clés :}
        \begin{itemize}
            \item \textbf{Taux de variation :} $\frac{V_{final} - V_{initial}}{V_{initial}} \times 100$.
            \item \textbf{Part en \% :} $\frac{Valeur\_partie}{Total} \times 100$.
            \item \textbf{Rythme moyen annuel :} $\left(\frac{V_{final}}{V_{initial}}\right)^{\frac{1}{n}} - 1$ (ou approximation linéaire si demandé).
        \end{itemize}
    \item \textbf{Interpréter historiquement :} Relier les chiffres aux événements (Conférence de Berlin, conquête, crise 1929, WWI). \emph{Ex: Hausse exportations cacao 1900-1914 = installation planteurs + chemin de fer + travail forcé.}
    \item \textbf{Critiquer la source :} Données officielles coloniales (sous-estimation production vivrière, contournement douanes, main-d'œuvre invisible). Biais de la « statistique coloniale ».
    \item \textbf{Intégrer dans l'argumentation :} « Le graphique 1 montre que la part de l'Afrique dans le commerce mondial passe de X\% à Y\%, confirmant l'intégration forcée... »
\end{enumerate}
\textbf{Reflexe Bac :} Toujours \textbf{calculer»} au moins un taux de variation ou une part dans la copie. Ne pas dire « ça augmente », dire « ça augmente de +150\% entre 1890 et 1914 ».
\end{methode}

\begin{exoresolu}[Analyse de données : Évolution du réseau ferroviaire en AOF (1900-1930)]
\emph{Énoncé :} Le tableau ci-dessous présente le développement du réseau ferroviaire en Afrique Occidentale Française (AOF). Analysez ces données pour caractériser l'économie coloniale.
\begin{center}
\begin{tabularx}{\linewidth}{|c|X|X|X|X|}
\hline
\rowcolor{popblueL}
\textbf{Ligne (Année mise en service)} & \textbf{Longueur (km)} & \textbf{Port de départ} & \textbf{Terminus / Objectif} & \textbf{Produit principal visé} \\
\hline
Dakar - Saint-Louis (1885) & 265 & Dakar & Saint-Louis & Arachide (Bassin) \\
\hline
Dakar - Kayes - Bamako (1904-1924) & 1288 & Dakar & Bamako (Soudan) & Arachide, Or (Bambouk), Bétail \\
\hline
Abidjan - Ouagadougou (1904-1954) & 1155 & Abidjan & Ouagadougou (Haute-Volta) & Cacao, Coton, Karité \\
\hline
Conakry - Kankan (1900-1914) & 662 & Conakry & Kankan (Niger) & Banane, Palmier, Bauxite (plus tard) \\
\hline
Cotonou - Parakou (1906-1936) & 438 & Cotonou & Parakou (Dahomey) & Palmier, Coton \\
\hline
\textbf{Total AOF (vers 1930)} & \textbf{$\approx$ 3800 km} & & & \\
\hline
\end{tabularx}
\end{center}
\tcblower
\textbf{Solution détaillée :}

\textbf{1. Lecture et calculs}
\begin{itemize}
    \item \textbf{Total km :} $265 + 1288 + 1155 + 662 + 438 = \mathbf{3808 \text{ km}}$ (vers 1930, achèvement progressif).
    \item \textbf{Durée de construction :} Échelonnée sur 45 ans (1885-1930). Ligne la plus longue : Dakar-Niger (1288 km, 20 ans).
    \item \textbf{Coût humain (connaissance) :} Dakar-Niger : $\approx$ 20 000 travailleurs forcés (corvée, 1900-1924), mortalité élevée (paludisme, accidents). Congo-Océan (AEF) : 17 000 morts officiels.
\end{itemize}

\textbf{2. Analyse structurelle : L'économie d'enclave (Port-Mine/Planteur-Port)}
\begin{itemize}
    \item \textbf{Architecture en « étoile » (radiale) :} Toutes les lignes partent d'un \textbf{port côtier} (Dakar, Abidjan, Conakry, Cotonou) vers l'\textbf{intérieur}. \textbf{Aucune liaison transversale} (Est-Ouest) entre lignes (ex: pas de liaison Bamako-Ouagadougou, ni Kankan-Bamako).
    \item \textbf{Logique extractive :} Le terminus correspond à une \textbf{zone de production} (Bassin arachidier Kayes/Bamako, Boucle du cacao Abidjan, Bauxite Kankan). Le chemin de fer ne dessert pas les populations (villes contournées) mais la ressource.
    \item \textbf{Spécialisation par colonie :} Sénégal = Arachide (Dakar-Niger). Côte d'Ivoire = Cacao/Café (Abidjan-Ouaga). Guinée = Fruits/Bauxite (Conakry-Kankan). Dahomey = Palmier/Coton (Cotonou-Parakou). Soudan (Mali) = Grenier + Or.
\end{itemize}

\textbf{3. Mise en perspective historique}
\begin{itemize}
    \item \textbf{Chronologie :} Début précoce (1885 Dakar-Saint-Louis) = ancien comptoir. Accélération \textbf{post-Conférence de Berlin (1885)} et \textbf{loi sur les colonies (1900)} : investissements massifs métropole (emprunts coloniaux) pour \cle{effectivité} (Berlin) et rentabilité.
    \item \textbf{Financement :} Emprunts d'État (garantis par la France) + Budget local (impôt capitation, travail forcé). Rentabilité garantie par \cle{tarifs préférentiels} (pacte colonial).
    \item \textbf{Conséquences spatiales :} Naissance de \cle{villes-gares} (Kayes, Kati, Bouaké, Ferkessédougou, Kankan, Parakou) au détriment des anciennes capitales (Ségou, Kong, Abomey). \cle{Dakar} devient capitale fédérale (1902) grâce au rail.
\end{itemize}

\textbf{4. Conclusion synthétique}
Ces données chiffrées prouvent que le chemin de fer colonial n'est \textbf{pas un outil d'aménagement du territoire} (désenclavement, intégration régionale) mais un \textbf{outil d'extraction} (mise en valeur « mise en valeur » coloniale). Il structure l'espace AOF en \cle{bassins de production spécialisés} tournés vers l'Europe, ignorant les complémentarités africaines (bétail Sahel $\leftrightarrow$ céréales Soudan $\leftrightarrow$ cacao Côte d'Ivoire). Cet héritage explique la faiblesse actuelle du commerce intra-africain (infrastructures inadaptées) et la nécessité de projets comme le \cle{Train Régional Express (TER) Dakar} ou la \cle{ligne Abidjan-Ouaga-Kaya} (réhabilitation récente).
\end{exoresolu}

\begin{attention}[Les 5 erreurs qui coûtent des points au Bac (Histoire - Impérialisme)]
\begin{enumerate}
    \item \textbf{Confondre « Partage » (diplomatie, Berlin 1885) et « Conquête » (militaire, terrain, 1885-1914).} \rightarrow \textbf{Sanction :} Hors sujet partiel. \textbf{Correctif :} Distinguer clairement la phase des \cle{traité de partage} (bureaux européens) et la \cle{pacification} (colonnes militaires, résistances africaines). Citer des dates distinctes.
    \item \textbf{Présenter les Africains comme passifs ou absents.} \rightarrow \textbf{Sanction :} Vision colonialiste, note plafonnée. \textbf{Correctif :} Systématiser : \cle{Résistances} (Samory, Behanzin, Menelik, Lat Dior, Kabarega, Mandume) + \cle{Collaborations} (chefs ralliés, évolués, tirailleurs, commerçants) + \cle{Adaptations} (migrations, islam, christianisme). Utiliser des \textbf{noms propres}.
    \item \textbf{Réduire l'impérialisme à une seule cause (ex: « seulement économique »).} \rightarrow \textbf{Sanction :} Analyse réductrice, manque de nuance. \textbf{Correctif :} Appliquer la grille \textbf{Économique / Politique-Stratégique / Idéologique / Juridique (Berlin)}. Montrer l'\emph{interaction} (ex: Kapitalisme $\to$ Rivalités $\to$ Berlin $\to$ Conquête).
    \item \textbf{Mélanger les modèles d'administration (France vs R.-U.) ou les figer.} \rightarrow \textbf{Sanction :} Erreur factuelle grave. \textbf{Correctif :} \textbf{France} = Dualité \cle{Citoyens (4 Communes) / Sujets (Code indigénat)} + \cle{Association} (pratique). \textbf{R.-U.} = \cle{Indirect Rule} (Native Authorities) + \cle{Direct Rule} (Kenya, Rhodésie, colons). \textbf{Allemagne} = \cle{Direct Rule} militarisé. Nuancer : « La France \emph{théorise} l'assimilation mais \emph{pratique} l'association. »
    \item \textbf{Oublier le bilan 1914 (chiffres clés) et l'héritage.} \rightarrow \textbf{Sanction :} Conclusion absente ou vague. \textbf{Correctif :} Mémoriser le \textbf{chiffre clé : 90\% colonisé / 2 États indépendants (Libéria, Éthiopie)}. Conclure obligatoirement sur : \cle{Frontières artificielles}, \cle{Économie extravertie}, \cle{Nouvelles élites}, \cle{Germes du nationalisme}. Phrase choc : « Berlin 1885 a dessiné la carte de l'Afrique de 1960. »
\end{enumerate}
\end{attention}

\newpage

\coursec{Exercices}

\courssub{Je vérifie mes connaissances}

\begin{exercice}[QCM : L'impérialisme et le partage de l'Afrique]
Parmi les propositions suivantes, une seule est exacte. Recopie la lettre correspondante.
\begin{enumerate}
    \item La conférence de Berlin (1884-1885) a pour principal résultat :
    \begin{itemize}
        \item[A.] L'indépendance immédiate des États africains.
        \item[B.] Le tracé des frontières africaines sans tenir compte des réalités ethniques.
        \item[C.] L'interdiction de la traite négrière en Afrique de l'Est.
        \item[D.] La création de la Société des Nations.
    \end{itemize}
    \item L'impérialisme européen en Afrique à la fin du XIX\textsuperscript{e} siècle s'explique principalement par :
    \begin{itemize}
        \item[A.] Le désir de propager la démocratie.
        \item[B.] La recherche de matières premières, de marchés et de prestige national.
        \item[C.] Une volonté purement humanitaire d'abolir l'esclavage.
        \item[D.] L'absence de rivalités entre puissances européennes.
    \end{itemize}
    \item Le traité Germano-Douala de juillet 1884 place le Cameroun sous :
    \begin{itemize}
        \item[A.] Protectorat britannique.
        \item[B.] Protectorat français.
        \item[C.] Protectorat allemand.
        \item[D.] Administration directe de la Société des Nations.
    \end{itemize}
    \item L'incident de Fachoda (1898) oppose :
    \begin{itemize}
        \item[A.] La France et l'Allemagne au Cameroun.
        \item[B.] La France et la Grande-Bretagne au Soudan.
        \item[C.] La Belgique et le Portugal au Congo.
        \item[D.] L'Italie et l'Éthiopie à Adoua.
    \end{itemize}
\end{enumerate}
\end{exercice}

\begin{exercice}[Vrai ou Faux justifié]
Indique si chaque affirmation est vraie (V) ou fausse (F) et justifie brièvement ta réponse par un fait historique précis.
\begin{enumerate}
    \item L'impérialisme est une politique d'expansion territoriale et d'hégémonie économique, politique et culturelle d'un État sur d'autres territoires.
    \item La résistance de Samory Touré a été brisée rapidement en moins d'un an par les Français.
    \item Le roi Béhanzin du Dahomey a collaboré activement avec les Français pour moderniser son armée.
    \item En 1914, seuls l'Éthiopie et le Libéria échappent au partage colonial de l'Afrique.
\end{enumerate}
\end{exercice}

\begin{exercice}[Définitions et notions clés]
Donne une définition précise et historique (2 à 3 lignes) pour chacun des termes suivants :
\begin{enumerate}
    \item L'impérialisme (contexte fin XIX\textsuperscript{e} siècle).
    \item La Conférence de Berlin (1884-1885).
    \item Le protectorat (distinction avec la colonie).
    \item La « mission civilisatrice » (argumentaire de Jules Ferry).
\end{enumerate}
\end{exercice}

\begin{exercice}[Repères chronologiques et géographiques]
\begin{enumerate}
    \item Replace les événements suivants dans l'ordre chronologique en indiquant la date :
    \begin{itemize}
        \item Conférence de Berlin.
        \item Traité Germano-Douala.
        \item Incident de Fachoda.
        \item Bataille d'Adoua.
        \item Capture de Samory Touré.
    \end{itemize}
    \item Sur une frise chronologique (fournie en annexe ou à tracer), positionne ces cinq événements.
\end{enumerate}
\end{exercice}

\courssub{Je m'entraîne}

\begin{exercice}[Analyse de la conquête française en Afrique de l'Ouest]
À partir de tes connaissances, réponds aux questions suivantes :
\begin{enumerate}
    \item Cite les deux grands axes de pénétration française en Afrique de l'Ouest (du Sénégal vers l'intérieur et de la Côte d'Ivoire vers le nord).
    \item Quel général français est l'artisan principal de la conquête du Soudan français (actuel Mali) face à l'empire de Samory Touré ?
    \item Explique pourquoi la résistance d'Almamy Samory Touré a duré près de 16 ans (1882-1898). Cite au moins trois facteurs (organisation militaire, tactique, diplomatie, armement).
\end{enumerate}
\end{exercice}

\begin{exercice}[La conquête britannique et les autres puissances]
\begin{enumerate}
    \item Présente la stratégie britannique du « Cap au Caire » (Cecil Rhodes) et cite deux obstacles majeurs rencontrés.
    \item Quel roi belge est à l'origine de la création de l'État indépendant du Congo (EIC) ? Quelle en fut la particularité administrative au début ?
    \item Cite deux résistances africaines face à la conquête britannique (une en Afrique australe, une en Afrique de l'Est) et leur issue.
\end{enumerate}
\end{exercice}

\begin{exercice}[Typologie des réactions africaines : Résistances vs Collaborations]
Complète le tableau comparatif suivant :

\begin{center}
\begin{tabularx}{\linewidth}{|l|X|X|}
\hline
\textbf{Critères} & \textbf{Résistance (ex : Samory, Béhanzin, Menelik II)} & \textbf{Collaboration (ex : Rois Douala, Chefs peuls, Lewanika)} \\
\hline
Motivation principale & & \\
\hline
Moyens d'action & & \\
\hline
Rapport aux Européens & & \\
\hline
Issues fréquentes & & \\
\hline
\end{tabularx}
\end{center}
\end{exercice}

\begin{exercice}[Étude de cas : La collaboration des rois Douala (Cameroun, 1884)]
Lis le court extrait suivant (simulé) : « Nous, rois Bell, Akwa, Deido, acceptons la protection de l'Empire allemand pour le commerce et la sécurité, conservant nos droits coutumiers sur la terre et la justice. »
\begin{enumerate}
    \item Identifie le type de traité signé (protectorat, cession, alliance).
    \item Pourquoi parle-t-on de « malentendu » entre les signataires africains et l'administration allemande ?
    \item Cite deux conséquences de cette collaboration pour la société douala entre 1884 et 1914.
\end{enumerate}
\end{exercice}

\courssub{Je me prépare au Bac}

\begin{exercice}[Dissertation type Baccalauréat Technique]
\textbf{Sujet :} « Les réactions africaines face à la conquête européenne (résistances et collaborations) furent-elles vouées à l'échec ? »
\begin{enumerate}
    \item Rédige une introduction complète (accroche, définition des termes, analyse du sujet, problématique, annonce du plan).
    \item Développe le plan en deux parties équilibrées :
    \begin{itemize}
        \item \textbf{Partie I :} L'illusion du succès ou les limites structurelles de l'opposition africaine (supériorité technique, désunion, ruse coloniale).
        \item \textbf{Partie II :} Des formes d'adaptation et de survie qui préparent les nationalismes (échec militaire mais héritage symbolique, collaboration comme gestion du quotidien, naissance d'une élite).
    \end{itemize}
    \item Rédige une conclusion synthétique et une ouverture sur la décolonisation.
\end{enumerate}
\textit{Durée conseillée : 1h30. Évaluation sur 20 points : Introduction (4), Développement (12), Conclusion/Ouverture (4).}
\end{exercice}

\begin{exercice}[Étude critique de documents : Le partage de l'Afrique (1885-1914)]
\textbf{Document 1 :} Extrait du discours de Jules Ferry à la Chambre des députés (28 juillet 1885).
« \textit{Il faut le dire ouvertement : en effet, les races supérieures ont un droit vis-à-vis des races inférieures... Elles ont le devoir de civiliser les races inférieures... Dans la concurrence commerciale qui nous presse... les marchés sont nécessaires. } »

\textbf{Document 2 :} Carte schématique de l'Afrique en 1914 (fournie à l'examen : empires coloniaux français, britannique, allemand, belge, portugais, italien, espagnol ; États indépendants : Éthiopie, Libéria).

\textbf{Questions :}
\begin{enumerate}
    \item \textbf{Présentation :} Nature, auteur, date, contexte historique de chaque document.
    \item \textbf{Analyse du Document 1 :} Identifie et explique les trois arguments principaux de Jules Ferry pour justifier la colonisation (argument racial, argument moral/humanitaire, argument économique).
    \item \textbf{Analyse du Document 2 :} À l'aide de la carte, cite les deux puissances coloniales possédant les plus vastes empires en 1914. Nomme les deux États africains restés indépendants et explique le cas particulier de l'Éthiopie (bataille d'Adoua, 1896).
    \item \textbf{Synthèse / Confrontation :} Montre en quoi la carte (Doc 2) est la réalisation concrète du programme idéologique et économique exposé par Jules Ferry (Doc 1). Conclus sur le sort des sociétés africaines face à ce partage.
\end{enumerate}
\end{exercice}

\begin{exercice}[Exercice composite : Cartographie, Données chiffrées et Cas camerounais]
\textbf{Partie A : Croquis commenté (Afrique de l'Ouest, vers 1900)}
À partir d'un fond de carte muet de l'Afrique de l'Ouest (côtes, fleuves Niger/Sénégal fournis), réalise un croquis légendé et titré montrant :
\begin{itemize}
    \item Les limites de l'AOF (Afrique Occidentale Française) créée en 1895.
    \item Les deux grands axes de conquête française (flèches : Saint-Louis $\rightarrow$ Niger ; Côte d'Ivoire $\rightarrow$ Niger).
    \item L'étendue maximale de l'empire de Samory Touré (vers 1885) et le lieu de sa capture (Guélemou, 1898).
    \item Les capitales des colonies : Saint-Louis (puis Dakar), Bamako, Conakry, Abidjan, Cotonou.
\end{itemize}
\textit{La légende doit être organisée, hiérarchisée (symboles de surface, linéaires, ponctuels) et le croquis soigné.}

\textbf{Partie B : Analyse statistique des empires coloniaux en 1914}
\begin{center}
\begin{tabularx}{\linewidth}{|l|X|X|X|}
\hline
\textbf{Puissance} & \textbf{Superficie (km$^2$)} & \textbf{Population (hab.)} & \textbf{Principales colonies africaines} \\
\hline
France & 10 500 000 & 35 000 000 & AOF, AEF, Madagascar, Maghreb \\
\hline
Grande-Bretagne & 9 000 000 & 45 000 000 & Égypte, Soudan, Afrique de l'Est, Afrique du Sud, Nigéria, Gold Coast \\
\hline
Allemagne & 2 500 000 & 15 000 000 & Cameroun, Togo, Afrique du Sud-Ouest, Afrique de l'Est \\
\hline
Belgique & 2 300 000 & 15 000 000 & Congo belge \\
\hline
Portugal & 2 000 000 & 9 000 000 & Angola, Mozambique, Guinée \\
\hline
Italie & 1 900 000 & 1 500 000 & Libye, Érythrée, Somalie \\
\hline
\end{tabularx}
\end{center}
\begin{enumerate}
    \item Construis un diagramme en barres (sur papier millimétré ou via logiciel) comparant les \textbf{superficies} des empires coloniaux français et britannique en Afrique.
    \item Calcule la part (en \%) de la superficie totale africaine colonisée (environ 30 000 000 km$^2$) détenue par la France et par la Grande-Bretagne.
    \item Tire deux conclusions argumentées sur la rivalité franco-britannique à partir de ces données (superficie vs population, continuité territoriale).
\end{enumerate}

\textbf{Partie C : Texte argumenté — Le traité Germano-Douala (12 juillet 1884)}
Rédige un texte argumenté d'une quinzaine de lignes (environ 150 mots) expliquant pourquoi le traité Germano-Douala illustre à la fois la \textbf{collaboration} (intérêts communs initiaux, reconnaissance mutuelle) et les \textbf{malentendus} (notion de souveraineté vs protectorat, spoliation foncière ultérieure, travail forcé, résistance de Rudolf Duala Manga Bell en 1914). Utilise les termes : \cle{protectorat}, \cle{souveraineté}, \cle{malentendu}, \cle{spoliation}, \cle{résistance}.
\end{exercice}

\newpage

\coursec{Corrigés des exercices}

\coursec{Corrigés des exercices — L'impérialisme en Afrique (Première Technique)}

\begin{corrige}[Exercice 1 — QCM : L'impérialisme et le partage de l'Afrique]
\begin{enumerate}
    \item \textbf{Réponse B.} La conférence de Berlin (novembre 1884 – février 1885), réunie à l'initiative de Bismarck, fixe les règles du partage de l'Afrique : occupation effective des territoires, notification des prises de possession, liberté de navigation sur le Niger et le Congo. Les frontières sont tracées à la règle sur des cartes, sans tenir compte des ethnies, des royaumes et des réseaux commerciaux africains. A est faux (la conférence organise au contraire la conquête), C est faux (la traite n'est pas l'objet principal), D est faux (la SDN est créée en 1919).
    \item \textbf{Réponse B.} Les causes de l'impérialisme sont économiques (matières premières : caoutchouc, huile de palme, coton ; débouchés pour les produits industriels ; placement des capitaux), politiques (prestige national, rivalités entre puissances) et idéologiques (prétendue « mission civilisatrice »). A et C sont des justifications de façade, D est le contraire de la réalité (les rivalités sont vives : Fachoda).
    \item \textbf{Réponse C.} Le 12 juillet 1884, les rois Douala (Bell, Akwa, Deido) signent avec Gustav Nachtigal un traité qui place le Cameroun sous protectorat allemand : naissance du \emph{Kamerun}. Les Britanniques, contactés auparavant par les Douala, arrivent trop tard (le 19 juillet).
    \item \textbf{Réponse B.} En 1898, la mission Marchand (France), venue de l'ouest, rencontre à Fachoda (Soudan) les troupes de Kitchener remontant le Nil depuis l'Égypte. La crise franco-britannique se règle par la retraite française ; la France renonce au Nil et obtient en échange des compensations en Afrique de l'Ouest (accord de 1899).
\end{enumerate}
\textbf{Points de vigilance :} ne pas confondre conférence de Berlin (1884-1885, règles du partage) et traité Germano-Douala (1884, cas camerounais) ; Fachoda oppose la France à la Grande-Bretagne, pas à l'Allemagne.
\end{corrige}

\begin{corrige}[Exercice 2 — Vrai ou Faux justifié]
\begin{enumerate}
    \item \textbf{Vrai.} L'impérialisme est bien la politique d'expansion d'un État visant la domination territoriale, économique, politique et culturelle d'autres peuples. Exemple : la France étend son empire de l'Afrique de l'Ouest à l'Indochine entre 1880 et 1914, en imposant sa langue, son administration et son économie d'échange.
    \item \textbf{Faux.} La résistance de Samory Touré dure près de seize ans (1882-1898). Il reconstitue son empire à l'est après chaque revers, achète des armes modernes via la Sierra Leone et la Côte d'Ivoire, et n'est capturé qu'en septembre 1898 à Guélemou, grâce à la trahison et à la surprise.
    \item \textbf{Faux.} Béhanzin, roi du Dahomey (actuel Bénin), est au contraire un grand résistant : il affronte les troupes du général Dodds lors des guerres de 1890 et 1892-1894, avant d'être vaincu, déporté en Martinique puis en Algérie. C'est son frère Agoli-Agbo, installé par les Français, qui collabore.
    \item \textbf{Vrai.} En 1914, seuls l'Éthiopie (victorieuse de l'Italie à Adoua en 1896 grâce à Menelik II) et le Libéria (État fondé en 1847 par d'anciens esclaves américains, protégé de fait par les États-Unis) conservent leur indépendance.
\end{enumerate}
\textbf{Points de vigilance :} la justification par un fait précis (date, acteur, lieu) est exigée ; une réponse V/F sans justification ne vaut que la moitié des points.
\end{corrige}

\begin{corrige}[Exercice 3 — Définitions et notions clés]
\begin{enumerate}
    \item \textbf{L'impérialisme :} politique d'expansion menée par les puissances industrielles européennes à la fin du XIX\textsuperscript{e} siècle, fondée sur la conquête de territoires, la domination économique (matières premières, marchés, capitaux) et l'hégémonie politique et culturelle sur des peuples jugés « inférieurs ».
    \item \textbf{La Conférence de Berlin (1884-1885) :} réunion diplomatique des puissances européennes (avec les États-Unis), organisée par Bismarck, qui fixe les règles du partage de l'Afrique : occupation effective, notification entre puissances, liberté de commerce sur le Congo et le Niger. Aucun Africain n'y est représenté.
    \item \textbf{Le protectorat :} forme de domination coloniale où l'État africain conserve théoriquement sa souveraineté interne et ses chefs, tandis que la puissance « protectrice » contrôle les affaires extérieures, la défense et l'économie. À la différence de la colonie de peuplement ou d'exploitation directe, l'administration locale subsiste formellement — mais dans les faits, la souveraineté est vidée de son contenu.
    \item \textbf{La « mission civilisatrice » :} argumentaire idéologique défendu notamment par Jules Ferry (discours du 28 juillet 1885) : les « races supérieures » auraient le « devoir » d'apporter la civilisation, l'instruction et le christianisme aux « races inférieures ». Ce discours moral sert de paravent aux intérêts économiques et stratégiques de la colonisation.
\end{enumerate}
\textbf{Points de vigilance :} une définition historique comporte toujours un contexte daté et un exemple ; éviter les définitions circulaires (« le protectorat est une protection »).
\end{corrige}

\begin{corrige}[Exercice 4 — Repères chronologiques et géographiques]
\begin{enumerate}
    \item Ordre chronologique :
    \begin{enumerate}
        \item \textbf{12 juillet 1884 :} traité Germano-Douala (protectorat allemand sur le Cameroun).
        \item \textbf{Novembre 1884 – février 1885 :} conférence de Berlin (règles du partage de l'Afrique).
        \item \textbf{1\textsuperscript{er} mars 1896 :} bataille d'Adoua (victoire de l'Éthiopie de Menelik II sur l'Italie).
        \item \textbf{Juillet 1898 :} incident de Fachoda (crise franco-britannique au Soudan).
        \item \textbf{29 septembre 1898 :} capture de Samory Touré à Guélemou par les Français.
    \end{enumerate}
    \emph{Remarque :} Fachoda (juillet 1898) précède de quelques semaines la capture de Samory (septembre 1898). Le traité Germano-Douala (juillet 1884) précède de quelques mois l'ouverture de la conférence de Berlin (novembre 1884).
    \item Frise chronologique (à tracer par l'élève) :

\begin{popfigure}
\begin{tikzpicture}[scale=1.0]
% Axe principal : 1884.0 à 1898.8 ~ 14.8 ans. Largeur 14.5cm. 1 an ~ 0.98cm. Origine x=0.5 pour marge.
\def\xorigin{0.5}
\def\scale{0.98}
\draw[->,thick,popdark] (\xorigin,0) -- (15,0) node[right]{\small Années};
% Graduations principales
\foreach \y/\label in {1884/1884, 1885/1885, 1890/1890, 1895/1895, 1898/1898}
  \draw[popdark] ({\xorigin+(\y-1884)*\scale},0.1) -- ({\xorigin+(\y-1884)*\scale},-0.1) node[below]{\small \label};
% Événements
% 1. Traité Germano-Douala : juillet 1884 ~ 1884.5
\draw[->,popblue,thick] ({\xorigin+(1884.5-1884)*\scale},0) -- ({\xorigin+(1884.5-1884)*\scale},1.0) node[above,align=center,font=\footnotesize]{Traité Germano-Douala\\(juillet 1884)};
% 2. Conférence Berlin : nov 1884 - fév 1885 ~ 1884.9
\draw[->,poporange,thick] ({\xorigin+(1884.9-1884)*\scale},0) -- ({\xorigin+(1884.9-1884)*\scale},1.8) node[above,align=center,font=\footnotesize]{Conférence de Berlin\\(1884-1885)};
% 3. Adoua : mars 1896 ~ 1896.25
\draw[->,popgreen,thick] ({\xorigin+(1896.25-1884)*\scale},0) -- ({\xorigin+(1896.25-1884)*\scale},1.0) node[above,align=center,font=\footnotesize]{Bataille d'Adoua\\(1\textsuperscript{er} mars 1896)};
% 4. Fachoda : juillet 1898 ~ 1898.5
\draw[->,poppink,thick] ({\xorigin+(1898.5-1884)*\scale},0) -- ({\xorigin+(1898.5-1884)*\scale},1.8) node[above,align=center,font=\footnotesize]{Incident de Fachoda\\(juillet 1898)};
% 5. Capture Samory : sept 1898 ~ 1898.75
\draw[->,poppurple,thick] ({\xorigin+(1898.75-1884)*\scale},0) -- ({\xorigin+(1898.75-1884)*\scale},1.0) node[above,align=center,font=\footnotesize]{Capture de Samory\\(sept. 1898)};
\end{tikzpicture}
\legende{Frise chronologique respectant l'échelle du temps : cinq événements clés de la conquête de l'Afrique (1884-1898).}
\end{popfigure}
\end{enumerate}
\textbf{Points de vigilance :} respecter l'échelle du temps sur la frise (les événements de 1884-1885 sont très proches, ceux de 1898 aussi) ; ne pas inverser Adoua (1896) et Fachoda (1898).
\end{corrige}

\begin{corrige}[Exercice 5 — Analyse de la conquête française en Afrique de l'Ouest]
\begin{enumerate}
    \item Les deux grands axes de pénétration française sont :
    \begin{itemize}
        \item l'axe \textbf{Saint-Louis du Sénégal $\rightarrow$ vallée du Sénégal $\rightarrow$ haut Niger} (vers Bamako et le Soudan), poussée vers l'est ;
        \item l'axe \textbf{Côte d'Ivoire / Golfe de Guinée $\rightarrow$ nord} (Grand-Bassam, puis Kong et Ouagadougou), visant à rejoindre le Niger et à couper l'expansion britannique depuis la Gold Coast.
    \end{itemize}
    L'objectif final est la jonction Sénégal-Niger et la constitution d'un ensemble continu de l'Atlantique au Tchad (AOF créée en 1895).
    \item Le principal artisan de la conquête du Soudan français est le colonel (puis général) \textbf{Louis Archinard}, qui dirige les campagnes contre l'empire toucouleur d'Ahmadou et contre Samory Touré à la fin des années 1880 et au début des années 1890 (prise de Ségou en mars 1890). On accepte aussi la mention des colonels Gallieni et Combes pour les campagnes postérieures contre Samory (1891-1898).
    \item Trois facteurs expliquant la longévité de la résistance de Samory Touré (1882-1898) :
    \begin{itemize}
        \item \textbf{Organisation militaire :} une armée permanente et hiérarchisée (la \emph{sofa}, cavalerie et infanterie), encadrée par des cadres fidèles, avec un système de recrutement dans l'empire du Wassoulou.
        \item \textbf{Tactique :} guerre de mouvement et stratégie du « repli-reconstitution » : Samory évite les batailles rangées perdues d'avance, pratique la politique de la terre brûlée et déplace son empire vers l'est (vers le nord de la Côte d'Ivoire actuelle) après chaque revers.
        \item \textbf{Armement et diplomatie :} il achète des fusils modernes (à répétition) via la Sierra Leone britannique et les comptoirs de la côte, fabrique des munitions, et joue des rivalités franco-britanniques ; il négocie aussi des trêves (traités de 1886, 1887, 1889) pour gagner du temps.
    \end{itemize}
    Il n'est capturé que par surprise, à Guélemou, le 29 septembre 1898.
\end{enumerate}
\textbf{Points de vigilance :} ne pas écrire que Samory « n'a jamais négocié » : il alterne guerre et diplomatie ; situer correctement le Wassoulou entre le haut Niger et la Côte d'Ivoire.
\end{corrige}

\begin{corrige}[Exercice 6 — La conquête britannique et les autres puissances]
\begin{enumerate}
    \item La stratégie du « Cap au Caire », théorisée par Cecil Rhodes, vise à relier par une bande continue de territoires britanniques (et un chemin de fer) Le Cap (Afrique du Sud) au Caire (Égypte), du sud au nord du continent. Deux obstacles majeurs :
    \begin{itemize}
        \item la \textbf{France} et son projet transversal ouest-est (Dakar-Djibouti), qui provoque la crise de \textbf{Fachoda} (1898) sur le Nil ;
        \item l'\textbf{Allemagne}, dont l'Afrique de l'Est allemande (Tanganyika) barre le passage au centre du continent, et les \textbf{Boers} en Afrique du Sud (guerre des Boers, 1899-1902).
    \end{itemize}
    \item Le roi \textbf{Léopold II de Belgique} crée l'État indépendant du Congo (EIC) en 1885, reconnu à Berlin. Particularité : c'est un domaine \textbf{personnel du roi}, et non une colonie de l'État belge ; Léopold II l'exploite (caoutchouc, ivoire) par des concessions privées, avec des exactions massives (travail forcé, mutilations). Ce n'est qu'en 1908, après le scandale international, que la Belgique annexe le Congo, qui devient le Congo belge.
    \item Deux résistances face aux Britanniques :
    \begin{itemize}
        \item \textbf{Afrique australe :} la résistance des \textbf{Zoulous} du roi Cetshwayo (victoire zouloue à Isandhlwana en 1879, puis défaite à Ulundi la même année ; le royaume est annexé).
        \item \textbf{Afrique de l'Est / Afrique de l'Ouest :} la résistance des \textbf{Nandi} du Kenya (Koitalel arap Samoei, 1895-1905, écrasée après l'assassinat du chef) ou celle des \textbf{Ashanti} de Prempeh I (Gold Coast, guerre de 1900, exil du roi).
    \end{itemize}
\end{enumerate}
\textbf{Points de vigilance :} bien distinguer l'EIC (propriété personnelle de Léopold II, 1885-1908) du Congo belge (après 1908) ; la révolte Maji-Maji (1905-1907) vise l'Allemagne au Tanganyika, pas la Grande-Bretagne — ne pas la citer comme résistance anti-britannique.
\end{corrige}

\begin{corrige}[Exercice 7 — Typologie des réactions africaines : Résistances vs Collaborations]
\begin{center}
\begin{tabularx}{\linewidth}{|l|X|X|}
\hline
\rowcolor{popblueL}
\textbf{Critères} & \textbf{Résistance (Samory, Béhanzin, Menelik II)} & \textbf{Collaboration (Rois Douala, chefs peuls, Lewanika)} \\
\hline
Motivation principale & Défendre la souveraineté, le territoire, la religion et les structures politiques face à l'agression étrangère. & Préserver le pouvoir du chef, obtenir des avantages commerciaux et militaires, éliminer des rivaux locaux grâce à l'allié européen. \\
\hline
Moyens d'action & Guerre (armées organisées, guérilla), mobilisation religieuse (mahdisme), diplomatie d'achat d'armes, terre brûlée. & Signature de traités, mise à disposition de guides, de porteurs et de soldats auxiliaires, fourniture de renseignements, adoption partielle des normes européennes. \\
\hline
Rapport aux Européens & Hostilité frontale ou alternance guerre/négociation ; refus de la domination. & Alliance recherchée, souvent vécue comme un partenariat entre égaux (malentendu sur le sens des traités). \\
\hline
Issues fréquentes & Défaite militaire (supériorité technique européenne), exil ou mort du chef (Samory déporté au Gabon, Béhanzin en Algérie) ; rare victoire (Adoua, 1896). & Maintien temporaire du chef comme relais de l'administration, puis marginalisation, spoliation foncière ou destitution (exil et mort de Rudolf Duala Manga Bell, pendu en 1914). \\
\hline
\end{tabularx}
\end{center}
\textbf{Points de vigilance :} la collaboration n'est pas une « trahison » à juger moralement : c'est un choix stratégique dans un rapport de force défavorable ; la frontière entre les deux attitudes est parfois floue (Samory négocie aussi ; les Douala finissent par résister en 1914).
\end{corrige}

\begin{corrige}[Exercice 8 — Étude de cas : La collaboration des rois Douala (Cameroun, 1884)]
\begin{enumerate}
    \item Il s'agit d'un \textbf{traité de protectorat} : les rois Bell, Akwa et Deido acceptent la « protection » de l'Empire allemand (traité Germano-Douala du 12 juillet 1884, signé avec Gustav Nachtigal). Ce n'est ni une cession totale du territoire, ni une simple alliance militaire : le texte mentionne la conservation des droits coutumiers sur la terre et la justice.
    \item On parle de « malentendu » parce que les deux parties n'entendent pas le traité de la même manière :
    \begin{itemize}
        \item pour les \textbf{rois Douala}, il s'agit d'un contrat commercial et diplomatique entre souverains : l'Allemand protège le commerce contre les concurrents, et les Douala conservent leur souveraineté interne (terre, justice, coutumes) ;
        \item pour l'\textbf{administration allemande}, le protectorat signifie la souveraineté effective du Reich : l'Allemagne s'arroge le droit d'administrer, de prélever l'impôt, d'exproprier et d'étendre son autorité à l'intérieur du pays.
    \end{itemize}
    Le mot « protection » recouvre donc deux conceptions incompatibles de la souveraineté.
    \item Deux conséquences pour la société douala (1884-1914) :
    \begin{itemize}
        \item \textbf{Spoliation foncière et dépossession :} expropriations massives (notamment le projet d'expulsion des Douala de la Janga-Plateau à Douala en 1910-1914), perte du monopole du commerce de l'arrière-pays au profit des comptoirs allemands.
        \item \textbf{Marginalisation politique et bascule vers la résistance :} les chefs perdent leurs prérogatives judiciaires et fiscales ; le roi Rudolf Duala Manga Bell proteste, envoie une délégation à Berlin, est accusé de haute trahison et pendu le 8 août 1914 avec son secrétaire Ngoso Din.
    \end{itemize}
\end{enumerate}
\textbf{Points de vigilance :} le traité est signé avec l'Allemagne, pas avec la France ou la Grande-Bretagne ; dater précisément (12 juillet 1884) ; montrer que la collaboration initiale se transforme en résistance (cas Manga Bell), ce qui illustre la fluidité des réactions africaines.
\end{corrige}

\begin{corrige}[Exercice 9 — Dissertation type Baccalauréat Technique]
\textbf{Barème indicatif (sur 20) :} introduction (4 pts), développement structuré et documenté (12 pts, soit 6 par partie), conclusion et ouverture (4 pts).

\textbf{1. Introduction (corrigé rédigé) :}

\emph{Accroche :} En 1914, à l'issue de trois décennies de conquêtes, l'Afrique est presque entièrement partagée entre les puissances européennes ; seuls l'Éthiopie et le Libéria échappent à la domination coloniale.
\emph{Définition des termes :} les « réactions africaines » désignent l'ensemble des attitudes des sociétés et des chefs africains face à la pénétration européenne (1880-1914) : résistances armées ou diplomatiques, et collaborations sous forme de traités et d'alliances.
\emph{Analyse du sujet :} le sujet invite à interroger la notion d'« échec » : échec militaire et politique immédiat, ou réussite différée à long terme ?
\emph{Problématique :} les réactions africaines, militairement dépassées, ont-elles été stériles, ou ont-elles préparé, par l'héritage symbolique et l'apprentissage du jeu colonial, les luttes pour l'indépendance ?
\emph{Annonce du plan :} si les rapports de force condamnaient l'opposition frontale (I), les réactions africaines n'en constituent pas moins des formes d'adaptation et de survie qui nourriront les nationalismes (II).

\textbf{2. Développement (éléments attendus) :}

\emph{Partie I — L'illusion du succès : les limites structurelles de l'opposition africaine.}
\begin{itemize}
    \item Supériorité technique et logistique européenne : fusils à répétition, mitrailleuses Maxim, artillerie, télégraphe, vapeurs fluviaux, quinine contre le paludisme ; écrasement des armées africaines en bataille rangée (Dahomey, 1892-1894 ; Omdurman, 1898).
    \item Désunion africaine : rivalités entre États et ethnies exploitées par les conquérants (alliances des Français avec les ennemis de Samory ; guerres fratricides au Dahomey entre Béhanzin et Agoli-Agbo).
    \item Ruse et droit coloniaux : traités ambigus (malentendu du traité Germano-Douala, 1884), stratégie du « diviser pour régner », trahisons (capture de Samory par surprise à Guélemou, 1898).
    \item Bilan : presque toutes les résistances armées sont vaincues avant 1914 ; les collaborateurs eux-mêmes sont dépossédés (Manga Bell pendu en 1914). Sous l'angle immédiat, l'échec semble total.
\end{itemize}

\emph{Partie II — Des formes d'adaptation et de survie qui préparent les nationalismes.}
\begin{itemize}
    \item L'échec militaire n'efface pas l'héritage symbolique : Samory, Béhanzin, Menelik II deviennent des figures fondatrices des nationalismes et des panafricanismes du XX\textsuperscript{e} siècle ; la victoire d'Adoua (1896) prouve qu'une défaite européenne était possible.
    \item La collaboration comme gestion du quotidien : les chefs collaborateurs (Lewanika du Barotseland, certains chefs peuls) préservent partiellement leurs populations, leurs structures et leurs terres ; la collaboration est une stratégie de survie, non une adhésion.
    \item Naissance d'une élite : écoles de mission, fonctionnaires et commerçants africains (évolués, élites de Douala, de Dakar, de Lagos) s'approprient les armes culturelles et juridiques de l'Occident ; les protestations de Manga Bell (délégation à Berlin, recours au droit) annoncent les revendications politiques.
    \item Les résistances se prolongent sous d'autres formes : révoltes fiscales, mouvements religieux, grèves, puis partis politiques après 1945.
\end{itemize}

\textbf{3. Conclusion (corrigé rédigé) :}

Militairement, les réactions africaines furent bien vouées à l'échec : le fossé technologique, la désunion et la duplicité des traités rendaient la victoire presque impossible, Adoua exceptée. Mais politiquement et symboliquement, résistances et collaborations ont permis aux sociétés africaines de traverser la colonisation sans disparaître : elles ont légué des héros, des mémoires et une élite capable de contester l'ordre colonial avec ses propres armes. \emph{Ouverture :} c'est cette double filiation — mémoire des résistants et apprentissage des institutions — qui nourrira les nationalismes et aboutira aux indépendances des années 1960.

\textbf{Points de vigilance :} problématique explicite obligatoire ; chaque partie doit s'appuyer sur au moins trois exemples datés et localisés ; éviter le récit descriptif (il faut argumenter sur la notion d'« échec ») ; l'ouverture doit rester dans le prolongement du sujet (décolonisation).
\end{corrige}

\begin{corrige}[Exercice 10 — Étude critique de documents : Le partage de l'Afrique (1885-1914)]
\textbf{Barème indicatif (sur 20) :} présentation des documents (4), analyse du Doc 1 (6), analyse du Doc 2 (5), synthèse (5).

\textbf{1. Présentation des documents.}
\begin{itemize}
    \item \textbf{Document 1 :} texte politique, extrait du discours prononcé par Jules Ferry, président du Conseil français, devant la Chambre des députés le 28 juillet 1885, pour défendre la politique coloniale de la France (contexte : conquête du Tonkin, débuts du partage de l'Afrique après la conférence de Berlin).
    \item \textbf{Document 2 :} carte historique schématique représentant le découpage colonial de l'Afrique en 1914, à la veille de la Première Guerre mondiale (contexte : achèvement du partage, crise de Fachoda résolue, empires constitués).
\end{itemize}

\textbf{2. Analyse du Document 1 : les trois arguments de Jules Ferry.}
\begin{itemize}
    \item \textbf{Argument racial (pseudo-scientifique) :} « les races supérieures ont un droit vis-à-vis des races inférieures » : Ferry hiérarchise les peuples et prétend fonder la domination sur une supériorité naturelle — théorie raciste sans fondement, héritée du darwinisme social.
    \item \textbf{Argument moral et humanitaire :} « elles ont le devoir de civiliser les races inférieures » : c'est la « mission civilisatrice » (instruction, christianisme, lutte contre l'esclavage), qui transforme la conquête en obligation morale.
    \item \textbf{Argument économique :} « les marchés sont nécessaires » : la France industrielle a besoin de débouchés pour ses produits, de matières premières et de champs d'investissement dans la concurrence internationale.
\end{itemize}
Ferry articule donc une justification idéologique (races, civilisation) à une motivation matérielle (économie), la première servant de paravent à la seconde.

\textbf{3. Analyse du Document 2.}
\begin{itemize}
    \item Les deux puissances possédant les plus vastes empires en 1914 sont la \textbf{France} (environ 10,5 millions de km$^2$ : AOF, AEF, Madagascar, Maghreb) et la \textbf{Grande-Bretagne} (environ 9 millions de km$^2$ : Égypte, Soudan, Nigéria, Afrique de l'Est et du Sud).
    \item Les deux États restés indépendants sont l'\textbf{Éthiopie} et le \textbf{Libéria}.
    \item Cas de l'Éthiopie : l'empereur \textbf{Menelik II}, ayant modernisé son armée (achat de fusils européens) et unifié le pays, écrase l'armée italienne à la \textbf{bataille d'Adoua le 1\textsuperscript{er} mars 1896} ; l'Italie reconnaît l'indépendance éthiopienne par le traité d'Addis-Abeba (octobre 1896). C'est la seule victoire décisive et durable d'un État africain sur une puissance coloniale au XIX\textsuperscript{e} siècle.
\end{itemize}

\textbf{4. Synthèse / Confrontation.}
La carte de 1914 est la traduction territoriale du programme de 1885 : les « marchés nécessaires » de Ferry sont devenus des empires continus (AOF, AEF, bande britannique Le Cap-Le Caire presque réalisée) ; la « mission civilisatrice » a fourni le vocabulaire moral d'une occupation effective décidée à Berlin. La confrontation des deux documents révèle le décalage entre le discours (droit, devoir, civilisation) et la réalité (conquête armée, frontières arbitraires, exploitation). Pour les sociétés africaines, le partage signifie la perte de souveraineté, le bouleversement des frontières ethniques, le travail forcé et la spoliation — mais aussi, en réaction, la naissance des résistances puis des nationalismes qui déféront ces empires au XX\textsuperscript{e} siècle.

\textbf{Points de vigilance :} toujours présenter un document avant de l'analyser (nature, auteur, date, contexte) ; citer le texte pour appuyer l'analyse ; ne pas écrire que Ferry « a raison » : l'étude critique suppose de dénoncer le caractère pseudo-scientifique de l'argument racial.
\end{corrige}

\begin{corrige}[Exercice 11 — Exercice composite : Cartographie, données chiffrées et cas camerounais]
\textbf{Partie A — Croquis commenté (éléments de correction).}
Le croquis doit comporter un titre (« La conquête française et l'empire de Samory en Afrique de l'Ouest, vers 1900 »), une flèche du nord et une légende hiérarchisée :
\begin{itemize}
    \item \textbf{Symboles de surface :} aplat coloré pour l'AOF (créée en 1895, regroupant Sénégal, Soudan français, Guinée, Côte d'Ivoire, Dahomey, puis Mauritanie, Niger, Haute-Volta) ; aplat hachuré d'une autre couleur pour l'empire de Samory à son extension maximale (vers 1885, du haut Niger au nord de la Côte d'Ivoire).
    \item \textbf{Symboles linéaires :} deux flèches épaisses de conquête : Saint-Louis $\rightarrow$ vallée du Sénégal $\rightarrow$ Bamako (Niger) ; côte ivoirienne $\rightarrow$ Kong $\rightarrow$ Ouagadougou (vers le Niger).
    \item \textbf{Symboles ponctuels :} points nommés pour Saint-Louis/Dakar, Bamako, Conakry, Abidjan (Grand-Bassam accepté), Cotonou ; étoile ou croix pour Guélemou (lieu de la capture de Samory, 1898, au sud-est de l'empire).
\end{itemize}
\textbf{Points de vigilance :} la légende doit être organisée (surfaces, lignes, points) ; les flèches doivent partir de la côte vers l'intérieur ; l'empire de Samory ne doit pas être confondu avec l'AOF (superposition partielle à gérer par des hachures).

\textbf{Partie B — Analyse statistique.}
\begin{enumerate}
    \item Diagramme en barres attendu : deux barres verticales, axe des ordonnées gradué en millions de km$^2$ (de 0 à 12), barre « France » à 10,5 et barre « Grande-Bretagne » à 9,0. (Sur papier millimétré : échelle conseillée 1 cm pour 1 million de km$^2$.)
    \item Calculs des parts (superficie totale colonisée $\approx$ \num{30000000} km$^2$) :
    \begin{align*}
    \text{Part de la France} &= \frac{\num{10500000}}{\num{30000000}} \times 100 = 35\,\% \\
    \text{Part de la Grande-Bretagne} &= \frac{\num{9000000}}{\num{30000000}} \times 100 = 30\,\%
    \end{align*}
    À elles deux, la France et la Grande-Bretagne contrôlent donc $35 + 30 = 65\,\%$ de l'Afrique colonisée.
    \item Deux conclusions argumentées :
    \begin{itemize}
        \item \textbf{Superficie contre population :} la France possède le plus vaste empire (35\,\%) mais moins peuplé (35 millions d'habitants) que l'empire britannique (45 millions pour 30\,\%) : la France domine de grands espaces sahéliens et sahariens peu denses, tandis que la Grande-Bretagne contrôle des territoires plus riches et plus peuplés (Égypte, Nigéria) — rivalité entre quantité d'espace et qualité économique.
        \item \textbf{Continuité territoriale :} les deux empires forment des blocs continus (France : bande ouest-est Dakar-Djibouti presque réalisée ; Grande-Bretagne : axe nord-sud Le Cap-Le Caire), ce qui explique leur affrontement là où les deux axes se croisent : le Nil, d'où la crise de Fachoda (1898), réglée sans guerre par le partage des zones d'influence (1899, puis Entente cordiale de 1904).
    \end{itemize}
\end{enumerate}
\textbf{Points de vigilance :} vérifier deux fois les divisions (10,5/30 = 0,35 ; 9/30 = 0,30) ; le diagramme doit porter un titre, des axes légendés avec unités et une source ; ne pas additionner superficie et population.

\textbf{Partie C — Texte argumenté (corrigé rédigé, environ 150 mots).}

Le traité Germano-Douala du 12 juillet 1884 illustre parfaitement l'ambiguïté des réactions africaines à la conquête. Il naît d'une \textbf{collaboration} réelle : les rois Bell, Akwa et Deido cherchent un allié européen pour sécuriser leur commerce et contrer leurs rivaux, tandis que l'Allemagne obtient une base sur le golfe de Guinée ; le traité reconnaît formellement les droits coutumiers des signataires. Mais il repose sur un \cle{malentendu} fondamental : pour les Douala, le \cle{protectorat} est un contrat entre souverains qui préserve leur \cle{souveraineté} interne ; pour Berlin, il fonde la souveraineté du Reich sur le \emph{Kamerun}. Dès lors, l'administration allemande impose l'impôt, le travail forcé et procède à la \cle{spoliation} foncière (projet d'expulsion des Douala en 1910). La collaboration se retourne alors en \cle{résistance} : Rudolf Duala Manga Bell conteste juridiquement ces abus, délègue à Berlin, et est pendu le 8 août 1914. Le traité montre ainsi comment une alliance inégale, fondée sur deux lectures opposées du même texte, débouche sur la dépossession puis la révolte.

\textbf{Points de vigilance :} employer obligatoirement les cinq termes imposés (\cle{protectorat}, \cle{souveraineté}, \cle{malentendu}, \cle{spoliation}, \cle{résistance}) ; respecter la longueur demandée ($\pm$ 20\,\%) ; articuler le texte (connecteurs : « mais », « dès lors », « ainsi ») et conclure.
\end{corrige}

\newpage

\coursec{Fiche bilan}

\courssub{Carte mentale de synthèse}

\begin{popfigure}
\begin{center}\tcbox[colback=poporange,colframe=poporange,coltext=white,arc=9pt,boxsep=4pt,left=6pt,right=6pt]{\bfseries\small L'impérialisme en Afrique (1870--1914)}\end{center}
\vspace{-2pt}
\begin{tcbraster}[raster columns=3,raster equal height=rows,raster column skip=6pt,raster row skip=6pt,enhanced,blankest]
\begin{tcolorbox}[enhanced,colback=white,colframe=poporange,boxrule=0.9pt,arc=3pt,title={Contexte et causes},fonttitle=\bfseries\scriptsize,coltitle=white,colbacktitle=poporange,left=4pt,right=4pt,top=2pt,bottom=2pt,boxsep=1pt,toptitle=1.5pt,bottomtitle=1.5pt,halign=left]
\begin{itemize}[nosep,leftmargin=*,itemsep=1pt,font=\scriptsize]
\item Crise économique 1873
\item Darwinisme social \& racisme scientifique
\item Recherche matières premières \& débouchés
\item Nationalismes et prestige
\end{itemize}
\end{tcolorbox}
\begin{tcolorbox}[enhanced,colback=white,colframe=popgreen,boxrule=0.9pt,arc=3pt,title={Conférence de Berlin (1884--85)},fonttitle=\bfseries\scriptsize,coltitle=white,colbacktitle=popgreen,left=4pt,right=4pt,top=2pt,bottom=2pt,boxsep=1pt,toptitle=1.5pt,bottomtitle=1.5pt,halign=left]
\begin{itemize}[nosep,leftmargin=*,itemsep=1pt,font=\scriptsize]
\item Partage pacifique entre Européens
\item Libre navigation Congo / Niger
\item Effectivité de l'occupation
\item Traité de protection
\end{itemize}
\end{tcolorbox}
\begin{tcolorbox}[enhanced,colback=white,colframe=poppurple,boxrule=0.9pt,arc=3pt,title={Conquêtes françaises},fonttitle=\bfseries\scriptsize,coltitle=white,colbacktitle=poppurple,left=4pt,right=4pt,top=2pt,bottom=2pt,boxsep=1pt,toptitle=1.5pt,bottomtitle=1.5pt,halign=left]
\begin{itemize}[nosep,leftmargin=*,itemsep=1pt,font=\scriptsize]
\item AOF (1895) Gouv. général Dakar
\item AEF (1910) Gouv. général Brazza
\item Mission civilisatrice Assimilation / Association
\item Administration directe (chefs canton)
\end{itemize}
\end{tcolorbox}
\begin{tcolorbox}[enhanced,colback=white,colframe=poppink,boxrule=0.9pt,arc=3pt,title={Conquêtes anglaises \& autres},fonttitle=\bfseries\scriptsize,coltitle=white,colbacktitle=poppink,left=4pt,right=4pt,top=2pt,bottom=2pt,boxsep=1pt,toptitle=1.5pt,bottomtitle=1.5pt,halign=left]
\begin{itemize}[nosep,leftmargin=*,itemsep=1pt,font=\scriptsize]
\item Indirect Rule (Lugard) Nigeria, Afrique Est
\item Afrique du Sud Cecil Rhodes, Boers
\item Allemagne Kamerun, Togo, etc.
\item Belgique (Congo), Portugal, Italie, Espagne
\end{itemize}
\end{tcolorbox}
\begin{tcolorbox}[enhanced,colback=white,colframe=popgold,boxrule=0.9pt,arc=3pt,title={Résistances africaines},fonttitle=\bfseries\scriptsize,coltitle=white,colbacktitle=popgold,left=4pt,right=4pt,top=2pt,bottom=2pt,boxsep=1pt,toptitle=1.5pt,bottomtitle=1.5pt,halign=left]
\begin{itemize}[nosep,leftmargin=*,itemsep=1pt,font=\scriptsize]
\item Samory Touré Mandinka (1881--1898)
\item Béhanzin Dahomey (1892--1894)
\item Ahmadou Tall Toucouleur (1890--1893)
\item Menelik II Éthiopie (Adoua 1896)
\item Maji Maji Tanganyika (1905--1907)
\item Lat Dior Cayor (1860--1886)
\end{itemize}
\end{tcolorbox}
\begin{tcolorbox}[enhanced,colback=white,colframe=popblue!70!black,boxrule=0.9pt,arc=3pt,title={Collaborations africaines},fonttitle=\bfseries\scriptsize,coltitle=white,colbacktitle=popblue!70!black,left=4pt,right=4pt,top=2pt,bottom=2pt,boxsep=1pt,toptitle=1.5pt,bottomtitle=1.5pt,halign=left]
\begin{itemize}[nosep,leftmargin=*,itemsep=1pt,font=\scriptsize]
\item Chefs traditionnels pouvoir / protection
\item Évolués / Lettrés emplois, statut social
\item Marchands / Courtiers intérêts commerciaux
\item Rivalités locales instrumentalisées
\end{itemize}
\end{tcolorbox}
\end{tcbraster}

\legende{Carte mentale récapitulative : causes, partage (Berlin), conquêtes (France, UK, autres), réactions africaines (résistance vs collaboration).}
\end{popfigure}

\courssub{Bilan essentiel}

\begin{aretenir}[Bilan essentiel -- Leçon n°2 : L'impérialisme en Afrique]

\textbf{Définitions clés}
\begin{itemize}
  \item \cle{Impérialisme} : politique d'expansion d'un État visant à dominer d'autres territoires (économique, politique, culturel). Phase 1870--1914 = \cle{néo-impérialisme} (capitalisme monopoliste, analyses de Hobson et Lénine).
  \item \cle{Conférence de Berlin} (nov. 1884 -- fév. 1885) : acte de naissance du partage de l'Afrique. Principes : libre navigation Niger/Congo, effectivité de l'occupation (occupation effective), notification des acquisitions aux autres puissances.
  \item \cle{Administration directe} (France) : pouvoir centralisé (Gouverneur général \rightarrow Commandants de cercle \rightarrow Chefs de canton auxiliaires nommés/révoqués), assimilation culturelle (écoles laïques, code civil), Code de l'indigénat (sanctions administratives), travail forcé (prestations, corvées).
  \item \cle{Administration indirecte} / \cle{Indirect Rule} (UK, Lugard) : pouvoir délégué aux chefs traditionnels maintenus (\emph{Native Authorities}), respect des coutumes (si non « répugnantes »), justice coutumière, moindre coût administratif, écoles souvent missionnaires.
  \item \cle{Résistance} : refus armé, diplomatique ou culturel de la domination coloniale. \cle{Collaboration} : alliance tactique ou idéologique avec le colonisateur pour préserver intérêts, autorité ou survivre.
\end{itemize}

\textbf{Repères chronologiques (à situer sur frise)}
\begin{center}
\begin{tabularx}{\linewidth}{|c|X|}\hline
\textbf{Date} & \textbf{Événement majeur} \\\hline
1870--1880 & Déclenchement de la « course à l'Afrique » (explorateurs : Stanley, Brazza, Livingstone) \\\hline
1884--1885 & Conférence de Berlin (Otto von Bismarck) : règles du partage \\\hline
1890--1898 & Grande résistance de Samory Touré (empire mandinka, capture 1898) \\\hline
1892--1894 & Guerre du Dahomey (Béhanzin vs général Dodds) \\\hline
1895 & Création de l'AOF (Afrique Occidentale Française), capitale Dakar \\\hline
1896 & Victoire éthiopienne d'Adoua (Menelik II bat l'Italie) \\\hline
1898 & Crise de Fachoda (conflit franco-britannique, retrait français) \\\hline
1902 & Fin guerre des Boers (UK annexe Transvaal, État d'Orange) \\\hline
1905--1907 & Révolte Maji Maji (Afrique orientale allemande, répression massive) \\\hline
1910 & Création de l'AEF (Afrique Équatoriale Française), capitale Brazzaville \\\hline
1914 & Quasi-totalité de l'Afrique colonisée (Éthiopie, Libéria indépendants) \\\hline
\end{tabularx}
\end{center}

\textbf{Méthode express : Commenter un document d'histoire (texte / carte / image / caricature)}
\begin{enumerate}
  \item \textbf{Identifier} : nature (texte officiel, discours, carte, photo, caricature), auteur, date, source, contexte historique précis (quoi ? qui ? quand ? où ? pourquoi ?).
  \item \textbf{Analyser} : idées forces / éléments visuels (légende, couleurs, symboles, mise en page), vocabulaire clé, arguments développés, thèse défendue.
  \item \textbf{Confronter} : croiser avec les connaissances du cours (confirmation, nuance, contradiction), situer le point de vue de l'auteur (acteur, témoin, propagande).
  \item \textbf{Conclure} : portée du document (témoignage, acte administratif, propagande, justification), limites (subjectivité, parti pris), apport à la problématique étudiée.
\end{enumerate}

\textbf{Tableau comparatif : Modèles coloniaux français vs britannique}
\begin{center}
\begin{tabularx}{\linewidth}{|l|X|X|}\hline
\rowcolor{popblueL}
\textbf{Critère} & \textbf{France (Directe / Assimilation $\rightarrow$ Association)} & \textbf{Royaume-Uni (Indirecte / Indirect Rule)} \\\hline
Organisation territoriale & Fédérations centralisées : AOF (1895), AEF (1910), Gouverneur général & Colonies distinctes (Crown Colonies, Protectorats), Gouverneur, \emph{Indirect Rule} \\\hline
Chefferie locale & Chefs de canton auxiliaires d'administration, nommés/révoqués par l'administration & Chefs traditionnels maintenus (\emph{Native Authorities}), rôle judiciaire/fiscal délégué \\\hline
Droit \& Justice & Code de l'indigénat (sanctions adm.), justice française, droit coutumier toléré & Coutumes locales appliquées (si pas « répugnantes »), cours indigènes (\emph{Native Courts}) \\\hline
Éducation & École laïque publique, français langue unique, formation d'élites « évoluées » & Écoles souvent missionnaires, langues vernaculaires + anglais, élites locales \\\hline
Économie & Mise en valeur par l'État, concessions, travail forcé (prestations, corvées, indigénat) & Compagnies chartées (début), puis culture de rente (coton, cacao, mines), travail migrant \\\hline
Idéologie officielle & \emph{Mission civilisatrice}, universalisme républicain, assimilation (théorique) & \emph{Dual Mandate} (Lugard) : développement du territoire + intérêts métropole, « association » \\\hline
\end{tabularx}
\end{center}

\end{aretenir}

\courssub{Checklist avant le Bac}

\begin{aretenir}[Checklist avant le Bac -- Leçon n°2 : L'impérialisme en Afrique]
\begin{itemize}
  \item $\square$ Je sais définir l'impérialisme et expliquer les causes (économiques : crise 1873, matières premières, capitaux ; politiques : nationalismes, stratégie ; idéologiques : darwinisme social, mission civilisatrice ; techniques : médecine, armement, transport) de la poussée des années 1870--1914.
  \item $\square$ Je connais les décisions majeures de la Conférence de Berlin (1884--1885) et leur portée : libre-échange bassin Congo/Niger, effectivité de l'occupation, abolition traite des noirs (affichée), notification des possessions.
  \item $\square$ Je sais décrire les conquêtes françaises : acteurs clés (Faidherbe \emph{(précurseur Sénégal)}, Gallieni, Brazza, Marchand, Gouraud), étapes (Sénégal, Soudan français, Dahomey, Congo, Tchad, Maroc 1912), création AOF (1895) et AEF (1910).
  \item $\square$ Je sais décrire les conquêtes britanniques : Afrique du Sud (Cecil Rhodes, guerre des Boers 1899--1902), Afrique de l'Est (Ouganda, Kenya), Afrique de l'Ouest (Nigeria \emph{Lugard}, Gold Coast), Égypte/Soudan (Kitchener), principe \emph{Indirect Rule} (Lugard).
  \item $\square$ Je cite au moins 3 autres puissances coloniales et leurs territoires : \textbf{Allemagne} (Kamerun, Togo, Deutsch-Ostafrika, Südwestafrika) ; \textbf{Belgique} (État indépendant du Congo \emph{Léopold II}, puis Congo belge 1908) ; \textbf{Portugal} (Angola, Mozambique, Guinée) ; \textbf{Italie} (Libye 1912, Érythrée, Somalie) ; \textbf{Espagne} (Rio de Oro, Guinée espagnole, Maroc nord).
  \item $\square$ Je distingue clairement \textbf{administration directe} (France : centralisation, chefs auxiliaires, code indigénat, assimilation) et \textbf{administration indirecte} (UK : décentralisation, chefs traditionnels \emph{Native Authorities}, coutumes, dual mandate).
  \item $\square$ Je connais les grandes figures de la \textbf{résistance africaine} : \textbf{Samory Touré} (Mandinka, 1881--1898, guerre de mouvement, empire déplacé) ; \textbf{Béhanzin} (Dahomey, 1892--1894, amazones, résistance farouche) ; \textbf{Ahmadou Tall} (Toucouleur, 1890--1893) ; \textbf{Menelik II} (Éthiopie, victoire d'Adoua 1896, seule victoire africaine durable) ; \textbf{Maji Maji} (Tanganyika, 1905--1907, résistance de masse spirituelle) ; \textbf{Lat Dior} (Cayor, Sénégal, 1886) ; \textbf{Rabah} (Tchad/Bornou, 1900).
  \item $\square$ J'explique les formes et motivations de la \textbf{collaboration} : chefs traditionnels (maintien pouvoir, protection contre rivaux) ; lettrés/évolués (accès emplois administration, statut social, écoles) ; marchands/courtiers (monopoles commerciaux, crédit) ; rivalités interethniques ou dynastiques instrumentalisées par le colonisateur.
  \item $\square$ Je sais situer sur une carte muette : les empires coloniaux en 1914 (France, UK, Allemagne, Belgique, Portugal, Italie, Espagne), les deux États indépendants (Éthiopie, Libéria), les zones de crises majeures (Fachoda, Maroc, Afrique du Sud, Congo).
  \item $\square$ Je suis capable de rédiger une introduction (accroche, définition du sujet, problématique, annonce du plan), un développement structuré (2 ou 3 parties équilibrées, arguments + exemples précis + dates), et une conclusion (bilan synthétique, ouverture) sur ce thème.
  \item $\square$ Je sais analyser un document (texte de loi, discours, carte de partage, caricature, photo, statistique) en appliquant la méthode en 4 étapes (Identifier, Analyser, Confronter, Conclure).
\end{itemize}
\end{aretenir}

\findecours

\end{document}
